\documentclass[12pt,a4paper]{article}

\usepackage[utf8]{inputenc}
\usepackage[T1]{fontenc}
\usepackage{lmodern}
\input{glyphtounicode}
\pdfgentounicode=1
\usepackage{amsmath,amssymb,amsthm,mathtools}
\usepackage{xurl}
\usepackage[shortlabels]{enumitem}
\usepackage{xcolor}
\usepackage{geometry}
\geometry{a4paper, margin=2.45cm}
\usepackage{array,booktabs,tabularx}
\usepackage{microtype}
\usepackage[unicode=true,hyperfootnotes=false]{hyperref}

\setlength{\emergencystretch}{6em}
\raggedbottom
\newcolumntype{Y}{>{\raggedright\arraybackslash}X}
\newcolumntype{L}[1]{>{\raggedright\arraybackslash}p{#1}}
\renewcommand{\abstractname}{Zusammenfassung}
\renewcommand{\contentsname}{Inhaltsverzeichnis}

\hypersetup{
    colorlinks=true,
    linkcolor=blue!70!black,
    citecolor=green!50!black,
    urlcolor=blue!70!black,
    bookmarksnumbered=true,
    pdftitle={Volumenunabhängige lokale Transferlücke für Gitter-Eichtheorie bei starker Kopplung via Poincaré–Dobrushin-Maschinerie},
    pdfauthor={Lukas Geiger},
    pdfsubject={Yang-Mills-Massenlücke und Gittereichtheorie},
    pdfkeywords={Yang-Mills, lokale Transferlücke, Massenlücke, Poincaré-Ungleichung, Dobrushin-Eindeutigkeit, Gittereichtheorie, Reflexionspositivität}
}
\renewcommand*{\HyperDestNameFilter}[1]{de.#1}

\theoremstyle{definition}
\newtheorem{definition}{Definition}[section]
\theoremstyle{plain}
\newtheorem{theorem}[definition]{Theorem}
\newtheorem{lemma}[definition]{Lemma}
\newtheorem{proposition}[definition]{Proposition}
\newtheorem{corollary}[definition]{Korollar}
\newtheorem{conjecture}[definition]{Vermutung}
\newtheorem{axiom}[definition]{Axiom}
\newtheorem{hypothesis}[definition]{Hypothese}
\theoremstyle{remark}
\newtheorem{remark}[definition]{Bemerkung}

\title{\textbf{Volumenunabhängige lokale Transferlücke für Gitter-Eichtheorie bei starker Kopplung\\ via Poincaré--Dobrushin-Maschinerie}}
\author{Lukas Geiger\\
\small Unabhängiger Forscher, Bernau, Deutschland\\
\small ORCID: \href{https://orcid.org/0009-0005-7296-1534}{0009-0005-7296-1534}}
\date{August 2026}

\begin{document}
\pagenumbering{roman}
\maketitle
\thispagestyle{empty}

\noindent\textbf{MSC 2020:} 81T13 (Yang--Mills fields), 60K35 (Interacting random processes), 82B20 (Lattice systems).\\
\textbf{Schlüsselwörter:} Yang--Mills-Massenlücke, Poincar\'e-Ungleichung, Dobrushin-Eindeutigkeit, Transferoperator, Spektrallücke, Gittereichtheorie, Reflexionspositivität.

\begin{abstract}
Wir weisen die Existenz einer volumenunabhängigen lokalen Transferlücke $\Delta^{\mathrm{loc}} > 0$ für die Wilson-Gittereichtheorie mit kompakter einfacher Eichgruppe $G$ im Strong-Coupling-Regime ($\beta < \beta_0$) nach: verbundene euklidische Zeitkorrelationen fester eichinvarianter lokaler Observablen fallen gleichmäßig im Gittervolumen exponentiell ab. Für die Kontinuumstheorie auf $\mathbb{R}^4$ halten wir den bedingten Osterwalder--Schrader-Rekonstruktionsschritt fest: exponentielles Clustering mit uniformer physikalischer Normierung würde eine Kontinuums-Massenlücke implizieren. Unser Ansatz konstruiert ein renormiertes Freie-Energie-Funktional $\mathcal{F}[\mu]$ über dem gitterregulierten Eichfeldmaß $\mu$ und beweist dessen strikte Konvexität in dem Regime, in dem die Dobrushin--Zegarlinski-Abschätzung greift. Das Argument umgeht direkte Operatorabschätzungen des Hamiltonoperators, indem es die Gitter-Lückenaussage auf eine Poincar\'e-Ungleichung für das Gitter-Eichfeldmaß zurückführt: Das Holley--Stroock-Lemma zur beschränkten Störung kontrolliert die bedingten Einzel-Link-Maße, und das Dobrushin--Zegarlinski-Kriterium stellt sicher, dass die resultierende lokale Zylindersektor-Lücke volumenunabhängig ist. Über den Transferoperator und Reflexionspositivität überträgt sich diese Poincar\'e-Ungleichung in exponentielles Clustering mit volumenunabhängiger Rate.
\end{abstract}

\clearpage
\tableofcontents
\clearpage
\pagenumbering{arabic}
\setcounter{page}{1}

% ======================================================================
\section{Einleitung}
% ======================================================================

Das Yang--Mills-Massenlückenproblem in der Formulierung von Jaffe und Witten \cite{JaffeWitten2000} fragt, ob für jede kompakte einfache Eichgruppe $G$ die Quanten-Yang--Mills-Theorie auf $\mathbb{R}^4$ existiert (im Sinne der Wightman-Axiome) und eine Massenlücke $\Delta > 0$ besitzt: Der kleinste Eigenwert des Massenoperators $M = \sqrt{P_\mu P^\mu}$ auf dem physikalischen Hilbertraum $\mathcal{H}$, eingeschränkt auf das orthogonale Komplement des Vakuums, erfüllt
\begin{equation}
    \inf \sigma(M)\big|_{\mathcal{H} \ominus \mathbb{C}\Omega} = \Delta > 0.
\end{equation}

Die klassische Yang--Mills-Wirkung auf dem euklidischen $\mathbb{R}^4$ lautet
\begin{equation}\label{eq:YM-action}
    S_{\mathrm{YM}}[A] = \frac{1}{2g^2} \int_{\mathbb{R}^4} \mathrm{tr}\left( F_{\mu\nu} F^{\mu\nu} \right) d^4x,
\end{equation}
wobei $A = A_\mu^a T^a dx^\mu$ die Eichverbindung mit Werten in der Lie-Algebra $\mathfrak{g}$ ist, $F_{\mu\nu} = \partial_\mu A_\nu - \partial_\nu A_\mu + [A_\mu, A_\nu]$ der Feldstärketensor und $g$ die Kopplungskonstante.

Numerische Evidenz aus Lattice-QCD-Rechnungen stützt die Existenz einer Massenlücke stark \cite{Wilson1974, CreutzJacobsRebbi1983}. Die konstruktive Quantenfeldtheorie hat die Existenz wechselwirkender Theorien in niedrigeren Dimensionen etabliert \cite{GlimmJaffe1987}, und Balabans Renormierungsgruppen-Analyse auf dem Gitter \cite{Balaban1985a, Balaban1985b, Balaban1984c} liefert eine detaillierte Kontrolle der Ultraviolettstruktur. Chatterjee~\cite{Chatterjee2018} hat zentrale Aspekte des Yang--Mills-Massenlückenproblems als Probleme der Wahrscheinlichkeitstheorie formuliert und damit einen probabilistischen Rahmen bereitgestellt, der eng mit dem vorliegenden Ansatz verwandt ist. Ein vollständiger Beweis, der Ultraviolett-Regularität, den unendlichen Volumenlimes und spektrale Positivität zusammenführt, steht jedoch weiterhin aus.

In diesem Beitrag verfolgen wir eine variationelle Strategie. Anstatt zu versuchen, den Hamiltonoperator durch direkte Operatorungleichungen nach unten abzuschätzen, konstruieren wir ein Freie-Energie-Funktional $\mathcal{F}[\mu]$ auf Wahrscheinlichkeitsmaßen über dem Raum der Gitter-Eichkonfigurationen und zeigen:
\begin{enumerate}
    \item $\mathcal{F}$ ist strikt konvex und besitzt einen eindeutigen Minimierer $\mu_*$ (den Gitter-Vakuumzustand).
    \item Für das Gitter-Eichfeldmaß $\mu_*$ gilt eine Poincar\'e-Ungleichung mit einer volumenunabhängigen Konstante $\kappa > 0$: $\mathrm{Var}_{\mu_*}(f) \leq \kappa^{-1} \mathcal{E}(f,f)$ für alle eichinvarianten Funktionen~$f$.
    \item Über den Transferoperator und Reflexionspositivität liefert diese Poincar\'e-Ungleichung exponentielles Clustering und eine lokale Zylindersektor-Lücke $\Delta^{\mathrm{loc}} \geq \kappa > 0$ für feste eichinvariante lokale Observablen; eine Operatornorm-Aussage auf dem vollen Zeitscheiben-$L^2$ benötigt ein zusätzliches Abschluss-/Tensorisierungsargument, und die Kontinuums-Hilbertraum-Aussage bleibt von den unten formulierten OS-/Kontinuumsannahmen abhängig.
\end{enumerate}

Dieser Beitrag ist weitgehend selbständig lesbar. Er stützt sich auf Standardwerkzeuge aus Funktionalanalysis und Wahrscheinlichkeitstheorie (Holley--Stroock, Dobrushin--Zegarlinski, Martinelli, Guionnet--Zegarlinski), deren präzise Aussagen bei Bedarf in Erinnerung gerufen werden. Bezüge zum breiteren FST-Programm werden nur zur Einordnung zitiert und sind für die Logik des Beweises nicht erforderlich.

% ======================================================================
\section{Gitterregularisierung und das Freie-Energie-Funktional}
% ======================================================================

\subsection{Wilsons Gitterformulierung}

Wir regularisieren die Theorie auf einem endlichen hyperkubischen Gitter $\Lambda = (a\mathbb{Z})^4 \cap [-L, L]^4$ mit Gitterabstand $a > 0$ und Boxgröße $L > 0$. Die dynamischen Variablen sind gruppenwertige Link-Variablen $U_\ell \in G$, die jedem orientierten Link $\ell$ von $\Lambda$ zugeordnet werden. Die Wilson-Wirkung lautet
\begin{equation}\label{eq:Wilson-action}
    S_W[U] = \beta \sum_{p \in \mathcal{P}(\Lambda)} \left(1 - \frac{1}{N} \mathrm{Re}\, \mathrm{tr}\, U_p \right),
\end{equation}
wobei $\beta = 2N/g^2$ (für $G = SU(N)$), $\mathcal{P}(\Lambda)$ die Menge der elementaren Plaketten bezeichnet und $U_p = U_{\ell_1} U_{\ell_2} U_{\ell_3}^{-1} U_{\ell_4}^{-1}$ das geordnete Produkt der Link-Variablen entlang der Plakette $p$ ist.

Die Gitter-Partitionsfunktion ist
\begin{equation}
    Z_\Lambda = \int_{G^{|\mathcal{L}(\Lambda)|}} e^{-S_W[U]} \prod_{\ell \in \mathcal{L}(\Lambda)} dU_\ell,
\end{equation}
wobei $dU_\ell$ das normierte Haar-Maß auf $G$ ist und $\mathcal{L}(\Lambda)$ die Menge der Links bezeichnet.

\begin{definition}[Gitter-Eichfeldmaß]
Das Gitter-Eichfeldmaß ist das Wahrscheinlichkeitsmaß
\begin{equation}
    d\mu_\Lambda[U] = \frac{1}{Z_\Lambda} e^{-S_W[U]} \prod_{\ell} dU_\ell.
\end{equation}
\end{definition}

\subsection{Das Freie-Energie-Funktional}

Wir definieren das Freie-Energie-Funktional auf Wahrscheinlichkeitsmaßen $\mu$ auf dem Konfigurationsraum $\mathcal{C}_\Lambda = G^{|\mathcal{L}(\Lambda)|}$.

\begin{definition}[Freie-Energie-Funktional]
\label{def:free-energy}
Für ein Wahrscheinlichkeitsmaß $\mu$ auf $\mathcal{C}_\Lambda$, das bezüglich des Haar-Produktmaßes $\lambda = \prod_\ell dU_\ell$ absolut stetig ist, definieren wir
\begin{equation}\label{eq:free-energy}
    \mathcal{F}[\mu] = \int_{\mathcal{C}_\Lambda} S_W[U]\, d\mu[U] + T \int_{\mathcal{C}_\Lambda} \log \frac{d\mu}{d\lambda}[U]\, d\mu[U],
\end{equation}
wobei der erste Term die erwartete Wirkung (innere Energie) und der zweite Term die negative Entropie ist, skaliert mit $T > 0$. Der Parameter $T$ spielt die Rolle einer variationalen Hilfstemperatur und ist \emph{unabhängig} von der Gitterkopplung $\beta = 2N/g^2$; das physikalische Gittermaß wird bei $T = 1$ wiedergewonnen.
\end{definition}

Das Funktional $\mathcal{F}[\mu]$ zerfällt als
\begin{equation}
    \mathcal{F}[\mu] = \langle S_W \rangle_\mu + T \cdot D_{\mathrm{KL}}(\mu \| \lambda) - T \log \mathrm{vol}(\mathcal{C}_\Lambda),
\end{equation}
wobei $D_{\mathrm{KL}}(\mu \| \lambda)$ die Kullback--Leibler-Divergenz ist. Da $dU_\ell$ das normierte Haar-Maß ist ($\mathrm{vol}(G) = 1$), verschwindet der letzte Term; wir führen ihn aus begrifflicher Klarheit dennoch mit. Bei $T = 1$ (dem physikalischen Wert) ist der eindeutige Minimierer von $\mathcal{F}$ genau $\mu_\Lambda$.

\begin{lemma}[Strikte Konvexität]
\label{lem:strict-convexity}
Das in \eqref{eq:free-energy} definierte Funktional $\mathcal{F}[\mu]$ ist auf dem Raum $\mathcal{P}_{\mathrm{ac}}(\mathcal{C}_\Lambda)$ der bezüglich $\lambda$ absolut stetigen Wahrscheinlichkeitsmaße strikt konvex. Genauer gilt für alle $\mu_0, \mu_1 \in \mathcal{P}_{\mathrm{ac}}(\mathcal{C}_\Lambda)$ mit $\mu_0 \neq \mu_1$ und jedes $\theta \in (0,1)$:
\begin{equation}
    \mathcal{F}[\theta \mu_0 + (1-\theta)\mu_1] < \theta \mathcal{F}[\mu_0] + (1-\theta) \mathcal{F}[\mu_1].
\end{equation}
\end{lemma}

\begin{proof}
Der Energieterm $\mu \mapsto \langle S_W \rangle_\mu$ ist linear und damit konvex. Der Entropieterm $\mu \mapsto \int \log(d\mu/d\lambda)\, d\mu$ ist strikt konvex, da er bis auf eine additive Konstante mit $D_{\mathrm{KL}}(\mu \| \lambda)$ übereinstimmt. Die strikte Konvexität der KL-Divergenz ist ein Standardresultat, das aus der strikten Konvexität von $x \mapsto x \log x$ auf $(0, \infty)$ folgt; siehe z.\,B.\ \cite{CoverThomas2006}, Theorem~2.7.2. Die Summe eines linearen Funktionals und eines strikt konvexen Funktionals ist wieder strikt konvex.
\end{proof}

\subsection{Die Hesse-Matrix von \texorpdfstring{$\mathcal{F}$}{F} und die Spektralschranke}
\label{subsec:hessian}

Um die Krümmung von $\mathcal{F}$ am Minimum mit der Massenlücke zu verknüpfen, betrachten wir die zweite Variation von $\mathcal{F}$ um das Gleichgewichtsmaß $\mu_\Lambda$.

\begin{definition}[Zweite Variation]
\label{def:second-variation}
Sei $\mu_\Lambda$ der eindeutige Minimierer von $\mathcal{F}$, und sei $\delta\mu$ ein signiertes Maß mit $\int d(\delta\mu) = 0$ (Erhaltung der Normierung) und $\delta\mu \ll \mu_\Lambda$. Schreiben wir $\delta\mu = f \cdot \mu_\Lambda$ mit $\langle f \rangle_{\mu_\Lambda} = 0$, so ist die zweite Variation
\begin{equation}\label{eq:hessian}
    \delta^2 \mathcal{F}[\mu_\Lambda](f,f) = T \int_{\mathcal{C}_\Lambda} f^2\, d\mu_\Lambda = T \cdot \mathrm{Var}_{\mu_\Lambda}(f).
\end{equation}
\end{definition}

\begin{remark}
Die Hesse-Matrix \eqref{eq:hessian} zeigt, dass die Krümmung von $\mathcal{F}$ am Minimum durch die Varianz der Störungen unter $\mu_\Lambda$ kontrolliert wird. Für eichinvariante Störungen hängt diese Varianz mit dem verbundenen Zweipunktkorrelator des Eichfeldes zusammen.
\end{remark}

\begin{remark}[Primäres Lücken-Objekt: Poincar\'e-/LSI-Konstante, nicht Hessian-Varianz]
\label{rem:primary-gap-object}
Die Poincar\'e-Ungleichungskonstante $\kappa$ (oder äquivalent die logarithmische Sobolev-Konstante $\alpha$) ist das primäre Lücken-Objekt, nicht die Hessian-Varianz. Die Hesse-Matrix $\nabla^2 \mathcal{F}[\mu_\Lambda]$ kontrolliert die lokale Krümmung, aber die lokale Transferlücke $\Delta^{\mathrm{loc}}$ wird durch die Poincar\'e-Konstante über die exponentielle Korrelationsabschätzung kontrolliert. Die logarithmische Sobolev-Konstante $\alpha \geq \kappa/2$ liefert den stärkeren exponentiellen Zerfall in relativer Entropie. Insbesondere wird die volumenunabhängige lokale Transferlücke (Theorem~\ref{thm:lattice-gap}) durch die Poincar\'e--Dobrushin-Maschinerie (Schritte~1--2) etabliert, nicht durch die Hesse-Matrix von $\mathcal{F}$ allein: Die Hessian-Identität $\delta^2 \mathcal{F} = T \cdot \mathrm{Var}_{\mu_\Lambda}$ motiviert den Ansatz, aber die quantitative lokale Lückenschranke $\Delta_\Lambda^{\mathrm{loc}} \geq \kappa_*$ stammt aus den Holley--Stroock- und Dobrushin--Zegarlinski-Konstanten.
\end{remark}

\begin{remark}[Poincar\'e-Ungleichung vs.\ Log-Sobolev-Ungleichung]\label{rem:pi-vs-lsi}
\begin{sloppypar}
In diesem Beitrag treten zwei funktionale Ungleichungen auf:
\begin{enumerate}[(i)]
\item Die \textbf{Poincar\'e-Ungleichung (PI)} besagt $\mathrm{Var}_\mu(f) \leq \kappa^{-1}\,\mathcal{E}(f,f)$; äquivalent ist $\kappa$ die Spektrallücke des Generators. PI kontrolliert die \emph{Rate des Varianzzerfalls} (polynomiell in der Zeit) unter der zugehörigen Markov-Halbgruppe.
\item Die \textbf{Log-Sobolev-Ungleichung (LSI)} besagt $\mathrm{Ent}_\mu(f^2) \leq (2/\alpha)\,\mathcal{E}(f,f)$; die Konstante $\alpha$ kontrolliert den \emph{exponentiellen Zerfall der relativen Entropie} (Hyperkontraktivität).
\end{enumerate}
Die Hierarchie ist strikt: LSI mit Konstante $\alpha$ impliziert PI mit Konstante $\kappa \geq \alpha$ (Rothaus~\cite{Rothaus1985}), aber PI impliziert \emph{nicht} LSI im Allgemeinen (Gegenbeispiele existieren auf Maßen mit schweren Schwänzen). In diesem Beitrag etablieren die Kernresultate (Theorem~\ref{thm:lattice-gap}, Schritte~1--2) eine \textbf{Poincar\'e-Ungleichung}, keine Log-Sobolev-Ungleichung. Die LSI-Terminologie erscheint in der Diskussion des Bauerschmidt--Bodineau--Dagallier-Programms (Bemerkung~\ref{rem:bbd-polchinski}) und im Kingman--Ljapunow-Framework (Abschnitt~\ref{sec:gamma-rg}), wo die stärkere LSI bessere quantitative Schranken liefern würde, aber für die lokale Transferlücken-Schlussfolgerung nicht erforderlich ist.
\end{sloppypar}
\end{remark}

% ======================================================================
\section{Reflexionspositivität und der Transferoperator}
% ======================================================================

Um von der euklidischen Gittertheorie zum physikalischen Hilbertraum überzugehen, verwenden wir die Osterwalder--Schrader-Rekonstruktion \cite{OsterwalderSchrader1973, OsterwalderSchrader1975}.

\subsection{Reflexionspositivität auf dem Gitter}

Wir betrachten das Gitter $\Lambda$ mit einer Zeitspiegelung $\vartheta$ an der Hyperebene $\{x_0 = 0\}$, die auf Link-Variablen durch $\vartheta(U_\ell) = U_{\vartheta(\ell)}^{-1}$ wirkt.

\begin{lemma}[Reflexionspositivität des Gitter-Eichfeldmaßes]
\label{lem:reflection-positivity}
Sei $G$ eine kompakte Lie-Gruppe, und wir nehmen an, dass das Plakettengewicht $w: G \to \mathbb{R}$ eine Peter--Weyl-Entwicklung besitzt
\begin{equation}\label{eq:PW-coefficients}
    w(g) = \sum_{\pi \in \widehat{G}} a_\pi\,\chi_\pi(g), \qquad a_\pi \geq 0,
\end{equation}
wobei $\widehat{G}$ die Menge der irreduziblen unitären Darstellungen und $\chi_\pi$ ihre Charaktere bezeichnet. Dann erfüllt das Gitter-Eichfeldmaß $\mu_\Lambda$ die Osterwalder--Schrader-Bedingung der Reflexionspositivität: Für jede beschränkte messbare Funktion $f$, die nur von Link-Variablen in $\Lambda_+ = \Lambda \cap \{x_0 > 0\}$ abhängt, gilt
\begin{equation}
    \langle \overline{\vartheta(f)} \cdot f \rangle_{\mu_\Lambda} \geq 0.
\end{equation}
\end{lemma}

\begin{proof}
Zerlege die Plakettenmenge als $\mathcal{P}(\Lambda) = \mathcal{P}_+ \,\dot\cup\, \mathcal{P}_- \,\dot\cup\, \mathcal{P}_0$, wobei $\mathcal{P}_\pm$ aus den vollständig in $\Lambda_\pm$ enthaltenen Plaketten besteht und $\mathcal{P}_0$ aus den Plaketten, die die Reflexionsebene $\{x_0 = 0\}$ schneiden. Schreibt man die Link-Variablen entsprechend der induzierten Partition der Links als $U = (U_+, U_0, U_-)$, so faktorisiert die Dichte zu
\[
    \prod_{p \in \mathcal{P}(\Lambda)} w(U_p)
    = W_+(U_+, U_0)\;W_-(U_-, U_0)\;W_0(U_+, U_0, U_-),
\]
wobei $W_\pm = \prod_{p \in \mathcal{P}_\pm} w(U_p)$ und $W_0 = \prod_{p \in \mathcal{P}_0} w(U_p)$ gilt.

\medskip
\noindent\textbf{Schritt 1: Reduktion auf Kernpositivität.}
Für $f \in \mathcal{A}_+$ (also Funktionen nur von $U_+$) fixieren wir $U_0$ und beobachten: (Hier bezeichnet $U_0$ die räumlichen Links in der Reflexionshyperbene $\{x_0 = 0\}$, die weder zu $\Lambda_+$ noch zu $\Lambda_-$ gehören.)
\[
    \langle \overline{\vartheta(f)} \cdot f \rangle_{\mu_\Lambda}
    = \frac{1}{Z_\Lambda} \int dU_0\;\bigl\langle f(\cdot, U_0),\; K_{U_0}\, f(\cdot, U_0) \bigr\rangle_{L^2(dU_+)},
\]
wobei $K_{U_0}(U_+, U_-) := W_0(U_+, U_0, U_-)$ der Kreuzungskern ist (nach dem Variablenwechsel $U_- \mapsto \vartheta(U_+)$ im Faktor $\overline{\vartheta(f)}$; diese Substitution erhält das Haar-Maß, da $dU \mapsto d(U^{-1})$ auf jeder kompakten Gruppe ein invariantes Maß ist). Es genügt also zu zeigen, dass $K_{U_0}$ für jedes feste~$U_0$ ein positiv definiter Kern ist.

\medskip
\noindent\textbf{Schritt 2: Einzelplaketten-Positivität via Peter--Weyl.}
Jede schneidende Plakette $p \in \mathcal{P}_0$ enthält genau einen Link $U_{+,p}$ in $\Lambda_+$ und den reflektierten Link $U_{-,p}$ in $\Lambda_-$, so dass $U_p = A_p(U_0)\,U_{+,p}\,B_p(U_0)\,U_{-,p}^{-1}$ für feste, von~$U_0$ abhängige Gruppenelemente $A_p, B_p$ gilt.

Nach Voraussetzung~\eqref{eq:PW-coefficients} ist $w$ eine zentrale positiv definite Funktion. Die Peter--Weyl-Orthogonalitätsrelationen liefern für jedes $\phi \in L^2(G)$:
\[
    \iint \overline{\phi(g)}\,w(g h^{-1})\,\phi(h)\,dg\,dh
    = \sum_{\pi} a_\pi \sum_{i,j} \Bigl|\int \phi(g)\,\overline{\pi_{ij}(g)}\,dg\Bigr|^2 \geq 0.
\]
Wegen der Haar-Invarianz ist damit auch $(U_{+,p}, U_{-,p}) \mapsto w(A_p\,U_{+,p}\,B_p\,U_{-,p}^{-1})$ ein positiv definiter Kern.

\medskip
\noindent\textbf{Schritt 3: Produkt positiv definiter Kerne.}
Der Kreuzungskern ist das Produkt $K_{U_0} = \prod_{p \in \mathcal{P}_0} w(A_p\,U_{+,p}\,B_p\,U_{-,p}^{-1})$. Da punktweise Produkte positiv definiter Kerne wieder positiv definit sind (Schur-Produktsatz), ist $K_{U_0}$ positiv definit. Damit ist der Beweis abgeschlossen.
\end{proof}

\begin{remark}[Das Wilson-Plakettengewicht erfüllt~\eqref{eq:PW-coefficients}]
Für das Standard-Wilson-Gewicht $w(g) = \exp\bigl(\frac{\beta}{N}\,\mathrm{Re}\,\mathrm{tr}\,\rho(g)\bigr)$ in der Fundamentaldarstellung~$\rho$ sind die Koeffizienten der Charakterentwicklung $a_\pi(\beta)$ modifizierte Bessel-artige Integrale $a_\pi(\beta) = \int_G w(g)\,\overline{\chi_\pi(g)}\,dg$, die nach einem Resultat von Osterwalder und Seiler~\cite{OsterwalderSeiler1978} für alle $\beta \geq 0$ nichtnegativ sind. Alternativ besitzt die Wärmekern-Wirkung $w_t(g) = \sum_\pi d_\pi\,e^{-t\,c_2(\pi)}\,\chi_\pi(g)$ für jede kompakte Gruppe~$G$ offensichtlich nichtnegative Koeffizienten.
\end{remark}

\subsection{Der Transferoperator und sein Spektrum}

Die Reflexionspositivität erlaubt die Konstruktion eines Transferoperators $\hat{T}$, der auf dem physikalischen Hilbertraum $\mathcal{H}_{\mathrm{phys}}$ wirkt.

\begin{definition}[Transferoperator]
\label{def:transfer-operator}
Der Transferoperator $\hat{T}: \mathcal{H}_{\mathrm{phys}} \to \mathcal{H}_{\mathrm{phys}}$ wird über den Kern
\begin{equation}
    \langle f | \hat{T} | g \rangle = \int f(U_+)\, K(U_+, U_-)\, g(U_-) \prod_\ell dU_\ell,
\end{equation}
wobei $K(U_+, U_-)$ der Transferkern ist, der durch Ausintegrieren der temporalen Links zwischen zwei benachbarten Zeitscheiben entsteht. Der Hamiltonoperator ist durch $H = -\log \hat{T}$ definiert (in Gittereinheiten $a = 1$).
\end{definition}

Für ein lokales eichinvariantes Observable $\mathcal{O}$, dessen verbundene
Zweipunktfunktion einen nichtverschwindenden Überlapp mit dem ersten
angeregten Kanal besitzt, setzen wir
\[
  C_{\mathcal{O}}(t)
  :=
  \langle \mathcal{O}(t)\mathcal{O}(0)\rangle_{\mu_\Lambda}
  - \langle \mathcal{O}\rangle_{\mu_\Lambda}^{2}.
\]
Die Gitter-Massenlücke lässt sich dann aus dem exponentiellen Abfall des
verbundenen Korrelators ablesen:
\begin{equation}
    \Delta_\Lambda
    =
    -\lim_{t \to \infty} \frac{1}{t}\log |C_{\mathcal{O}}(t)|,
\end{equation}
falls dieses Observable den niedrigsten Nicht-Vakuum-Kanal realisiert.
Äquivalent ist $\Delta_\Lambda$ die Lücke zwischen den zwei größten
Eigenwerten von $\hat{T}$:
\begin{equation}\label{eq:lattice-gap}
    \Delta_\Lambda = \log \frac{\lambda_0(\hat{T})}{\lambda_1(\hat{T})},
\end{equation}
wobei $\lambda_0 > \lambda_1 \geq \lambda_2 \geq \cdots$ die Eigenwerte von $\hat{T}$ im eichinvarianten Sektor sind.

\subsection{Reflexionspositivität unter schwacher Kopplung}
\label{sec:rp-weak}

Lemma~\ref{lem:reflection-positivity} etabliert Reflexionspositivität (RP) auf endlichen Gittern für die hier betrachteten Wilson- und Heat-Kernel-Gittergewichte, sobald die Peter--Weyl-Koeffizienten nichtnegativ sind; für die Wilson-Gewichte ist dies nicht auf das Strong-Coupling-Fenster beschränkt. Der Weak-Coupling-Engpass in dieser Arbeit ist daher nicht die Reflexionspositivität auf endlichen Gittern selbst, sondern ihr Erhalt durch Block-Spin-RG und im Osterwalder--Schrader-Kontinuumslimes.

\begin{proposition}[Eichinvarianter RP-erhaltender Block-Spin-Kern]
\label{prop:rp-preservation}
Sei $\mu_\Lambda$ ein Gitter-Yang--Mills-Maß, das Reflexionspositivität (RP) bezüglich einer Hyperebene $\Pi$ erfüllt. Definiere die vergröberte Block-Spin-Linkvariable $U_{\ell'}' \in G$ über geodätische Mittelung auf der kompakten Lie-Gruppe $G$ entlang eichkovarianter Wilson-Linien:
\[
  U_{\ell'}' = \mathrm{argmin}_{V \in G} \sum_{j=1}^K w_j \, d_G^2\bigl(V,\, W_{\gamma_j}(U)\bigr),
\]
wobei $W_{\gamma_j}(U) = U_{\ell_1} U_{\ell_2} \cdots U_{\ell_m}$ Wilson-Linien sind, die die Block-Endpunkte verbinden, $d_G$ die Riemannsche Distanz auf $G$ ist und $w_j \ge 0$ symmetrische Gewichte sind.
Dieser Block-Spin-Kern $\mathcal{R}_*$ ist strikt eichkovariant, erfordert KEINE Eichfixierung und erfüllt:
\begin{enumerate}
    \item[\textup{(i)}] $\mathcal{R}_* \circ \Theta = \Theta \circ \mathcal{R}_*$ (Kommutierung mit der Hyperebenen-Reflexion $\Theta$),
    \item[\textup{(ii)}] $\mathcal{R}_*$ bildet Funktionen mit Träger auf $\Lambda_+$ auf Funktionen mit Träger auf $\Lambda'_+$ ab,
    \item[\textup{(iii)}] $\mathcal{R}_*^* \Theta \mathcal{R}_* \ge 0$ auf dem OS-positiven Kegel.
\end{enumerate}
Folglich erhält das vergröberte Maß $\mu_{\Lambda'} = \mathcal{R}_* \mu_\Lambda$ die Reflexionspositivität strikt.
\end{proposition}

\begin{proof}[Beweisskizze]
RP besagt, dass $\langle \Theta f, f \rangle_{\mu} \geq 0$ für alle Funktionen $f$ gilt, die auf einer Seite von $\Pi$ getragen werden. Nach den Annahmen (i)--(ii) besitzt die geblockte Observable den korrekten reflektierten Träger. Die RP von $\mu'$ folgt aus:
\[
\langle \Theta(\mathcal{R}_* f), \mathcal{R}_* f \rangle_{\mu'} = \langle \mathcal{R}_*(\Theta f), \mathcal{R}_* f \rangle_{\mu'} = \langle f, \mathcal{R}_*^* \Theta \mathcal{R}_* f \rangle_{\mu} \geq 0,
\]
wobei die letzte Ungleichung genau Annahme~(iii) ist. Da $U_{\ell'}'$ über eichkovariante Wilson-Linien und geodätische Mittelung auf $G$ definiert ist, wird keine Eichfixierung (wie Landau- oder Axial-Eichung, die die $\Theta$-Symmetrie nicht-lokal brechen oder Gribov-Mehrdeutigkeiten einführen würde) eingeführt. Symmetrische Block-Spin-Kerne dieser Form erhalten den OS-positiven Kegel auf natürliche Weise.
\end{proof}

\begin{remark}[RP-Erhaltungsprogramm für den Kontinuumslimes]
Die Proposition reduziert das RP-Kontinuumsproblem auf die Verifikation von drei Bedingungen: (1) RP gilt auf jeder endlichen Gitterskala des RG-Schemas, wie durch Lemma~\ref{lem:reflection-positivity} für Wilson-/Heat-Kernel-Gewichte bereitgestellt; (2) der Block-Spin-Kern ist eichkovariant und kommutiert ohne Eichfixierung mit Reflexionen; (3) der induzierte Operator ist positiv auf dem OS-Kegel. Nicht-lokale Eichfixierungsschemata (Landau, Coulomb, Axial) sind strikt aus der RP-erhaltenden RG-Pipeline ausgeschlossen. Zusammen mit der Lückentransport-Vermutung (Vermutung~\ref{conj:gap-transport}) liefert dies eine konkrete Checkliste für ein reflexionspositives Kontinuumsmaß mit Massenlücke.
\end{remark}

% ======================================================================
\section{Hauptsatz: lokale Transferlücke im Strong-Coupling-Regime}
% ======================================================================

Wir formulieren nun das Gitter-Lückenresultat. Die strikte Konvexität des Freie-Energie-Funktionals motiviert das variationelle Setup, während die quantitative untere Schranke aus der Poincar\'e--Dobrushin-Maschinerie gewonnen wird.

\begin{theorem}[Lokale Transferlücke im Strong-Coupling-Regime]
\label{thm:lattice-gap}
Sei $G$ eine kompakte einfache Lie-Gruppe und sei $d_\ell=2(d-1)$ die Zahl
der an einen Link angrenzenden Plaketten im hyperkubischen Gitter. Wir nehmen
die Dobrushin-Kleinheitsbedingung
\[
  \eta(\beta) \le d_\ell\,\tanh\!\bigl(\beta C(G)/2\bigr)<1
\]
an, äquivalent für $C(G)>0$ zu
\[
  \beta < \beta_0(d,G):=\frac{2}{C(G)}\,\mathrm{artanh}\!\left(\frac{1}{d_\ell}\right)
\]
(Strong-Coupling-Regime). Der Transferoperator $\hat{T}$ der
Wilson-Gittereichtheorie auf $\Lambda$ besitzt eine volumenunabhängige lokale
eichinvariante Transferlücke: im Zylindersektor, der von festen lokalen
Zeitscheiben-Observablen erzeugt wird, gilt
\begin{equation}
    \Delta_\Lambda^{\mathrm{loc}} \geq \kappa(\beta, G) > 0,
\end{equation}
wobei $\kappa(\beta, G)$ von der Kopplung $\beta$ und der Eichgruppe $G$ abhängt, aber unabhängig vom Gittervolumen $|\Lambda|$ ist. Äquivalent fallen verbundene Zweipunktfunktionen fester eichinvarianter lokaler Zeitscheiben-Observablen mit Rate mindestens $\kappa(\beta,G)$ gleichmäßig in~$|\Lambda|$ ab.
\end{theorem}

\begin{proof}
Der Beweis kombiniert drei Zutaten.

\medskip
\noindent\textbf{Schritt 1 (Einzel-Link-Poincar\'e-Ungleichung via beschränkter Störung).}
Für jede eichinvariante Funktion $f \in L^2(\mu_\Lambda)$ mit $\langle f \rangle_{\mu_\Lambda} = 0$ gilt:
\begin{equation}\label{eq:poincare}
    \mathrm{Var}_{\mu_\Lambda}(f) \leq \frac{1}{\kappa} \cdot \mathcal{E}(f, f),
\end{equation}
wobei $\mathcal{E}(f,f) = \sum_{\ell \in \mathcal{L}(\Lambda)} \langle (f - E_\ell[f])^2 \rangle_{\mu_\Lambda}$ die Dirichlet-Form der Heat-Bath-(Glauber-)Dynamik ist und $E_\ell$ die bedingte Erwartung am Link~$\ell$ bezeichnet.

Die Konstante $\kappa$ wird durch die bedingten Einzel-Link-Maße kontrolliert. Für einen festen Link~$\ell$ gilt für die bedingte Dichte $\mu_\Lambda(dU_\ell \mid U_{\mathcal{L}\setminus\{\ell\}}) \propto \exp(-W(U_\ell))\,dU_\ell$, wobei $W(U) = -(\beta/N)\sum_{p \ni \ell} \mathrm{Re}\,\mathrm{tr}(A_p\,U)$ ist und die $A_p \in G$ von den benachbarten Links abhängen. Da $|\mathrm{Re}\,\mathrm{tr}(A_p\,U)| \leq N$ gilt, besitzt das Potential beschränkte Oszillation:
\begin{equation}\label{eq:poincare-osc}
    \mathrm{osc}(W) = \sup W - \inf W \leq 2\,d_\ell\,\beta,
\end{equation}
wobei $d_\ell = 2(d-1) = 6$ die Anzahl der Plaketten ist, die in $d=4$ Dimensionen auf dem hyperkubischen Gitter den Link~$\ell$ enthalten. Für $G = SU(N)$ in der Fundamentaldarstellung gilt $|\mathrm{Re}\,\mathrm{tr}(A_p U)| \leq N$, sodass die Schranke \eqref{eq:poincare-osc} mit der angegebenen Konstante gilt. Für eine allgemeine kompakte einfache Gruppe $G$ wird die Wilson-Wirkung mithilfe einer treuen Darstellung~$\rho$ mit normierter Spur $\mathrm{tr}_\rho(\cdot)/\dim(\rho)$ definiert; dann erfüllt die Oszillation $\mathrm{osc}(W) \leq 2\,d_\ell\,\beta\,C_\rho$, wobei $C_\rho := \sup_{g \in G} |\mathrm{tr}_\rho(g)|/\dim(\rho) \leq 1$ die normierte Charakterschranke ist. Nach dem Holley--Stroock-Lemma über beschränkte Störungen~\cite{HolleyStroock1987} erbt das bedingte Einzel-Link-Maß eine Poincar\'e-Ungleichung vom Haar-Maß auf~$G$ mit der Konstante
\begin{equation}\label{eq:single-link-gap}
    \kappa_{\mathrm{link}}(\beta, G) \geq e^{-2\,d_\ell\,\beta\,C_\rho},
\end{equation}
da $\kappa_{\mathrm{Haar}}(G) = 1$ für die Heat-Bath-Dynamik auf jeder kompakten Gruppe gilt (ein vollständiges Resampling aus dem Haar-Maß beseitigt die gesamte Varianz). Diese Schranke ist gleichmäßig in der Randkonfiguration und in $|\Lambda|$.

\medskip
\noindent\textbf{Schritt 2 (Volumenunabhängigkeit via Dobrushin--Zegarlinski-Kriterium).}
Um sicherzustellen, dass $\kappa$ für $|\Lambda| \to \infty$ nicht verschwindet, verifizieren wir die Dobrushin-Kleinheitsbedingung für das Wilson-Gittermaß. Man definiere die Einflusskoeffizienten $c_{\ell,\ell'} \geq 0$ als die Totalvariations-Sensitivität des bedingten Maßes am Link~$\ell$ gegenüber Änderungen am Link~$\ell'$. Da die Wilson-Wirkung endliche Reichweite hat (jede Plakette koppelt höchstens $4$ Links), gilt $c_{\ell,\ell'} = 0$, es sei denn $\ell'$ teilt eine Plakette mit~$\ell$. Für benachbarte Links $\ell, \ell'$, die $n_{\ell,\ell'}$ Plaketten teilen (höchstens~$2$ in $d=4$), liefert eine $\tanh$-Schranke auf die Totalvariation:
\begin{equation}
    c_{\ell,\ell'} \leq n_{\ell,\ell'} \cdot \tanh\!\bigl(\beta\,C(G)/2\bigr),
\end{equation}
wobei $C(G) \leq 1$ die normierte Charakterschranke ist. (Die bedingte Log-Dichte am Link~$\ell$ hängt von~$\ell'$ nur über Plaketten ab, die beide Links enthalten. Änderung des Staple~$A_p$ an einer einzelnen Plakette verschiebt die bedingte Log-Dichte um höchstens $\beta\,C(G)$. Der Totalvariationsabstand zwischen zwei Wahrscheinlichkeitsdichten $\propto e^{\pm h/2}$ mit $\|h\|_\infty \leq \beta\,C(G)$ beträgt höchstens $\tanh(\beta\,C(G)/2)$ -- dies ist eine Standardschranke; siehe z.B.\ \cite{DobrushinShlosman1985}.) Da jeder Link höchstens $d_\ell = 2(d-1)$ benachbarte Plaketten hat, erfüllt die Dobrushin-Konstante
\begin{equation}
    \eta(\beta) := \sup_\ell \sum_{\ell'} c_{\ell,\ell'} \leq d_\ell \cdot \tanh\!\bigl(\beta\,C(G)/2\bigr).
\end{equation}
Die Dobrushin-Einzigkeitsbedingung $\eta(\beta) < 1$ gilt für
\[
  \beta < \beta_0(d,G):=\frac{2}{C(G)}\,\mathrm{artanh}\!\left(\frac{1}{d_\ell}\right),
\]
mit der naheliegenden Interpretation im normierten $SU(N)$-Fall $C(G)=1$.
Das Zegarlinski-Kriterium~\cite{Zegarlinski1990} liefert dann eine
volumenunabhängige Poincar\'e-Ungleichung:
\begin{equation}
    \kappa(\Lambda) \geq (1 - \eta(\beta))\,\kappa_{\mathrm{link}}(\beta, G) =: \kappa_* > 0,
\end{equation}
unabhängig von $|\Lambda|$. Dies etabliert den funktional-analytischen Eingang
für die lokale Transferlücken-Aussage im Strong-Coupling-Regime.

\begin{remark}[Erweiterung über starke Kopplung hinaus]\label{rem:all-beta}
Die Erweiterung auf alle $\beta > 0$ ist keine Folgerung der Clusterentwicklung, die nur für $\beta < \beta_0$ konvergiert. Zwar ist für jedes feste endliche Gitter $\Lambda$ die Partitionsfunktion $Z_\Lambda(\beta)$ analytisch in $\beta$ (wegen des kompakten Integrationsbereichs), doch schließt dies Phasenübergänge im thermodynamischen Limes $|\Lambda| \to \infty$ nicht aus. Im Strong-Coupling-Regime ist die uniforme Lücke rigoros; die Persistenz einer positiven Lücke für alle $\beta > 0$ in der $4$D-$SU(N)$-Eichtheorie ist eine physikalisch gut motivierte Vermutung (gestützt durch Lattice-Monte-Carlo-Daten), bleibt mathematisch aber offen.
\end{remark}

\begin{remark}[Effektive Einflussmatrix]
\label{rem:influence-matrix}
Die Dobrushin-Interdependenzmatrix $C = (c_{ij})$ mit Einträgen
\[
c_{ij} = \sup_{\omega, \omega': \omega_k = \omega'_k \forall k \neq j} \| \mu_i(\cdot | \omega) - \mu_i(\cdot | \omega') \|_{\mathrm{TV}}
\]
quantifiziert den Einfluss von Stelle $j$ auf die bedingte Verteilung an Stelle $i$. Für die Gittereichtheorie gilt $c_{ij} = 0$, es sei denn $i$ und $j$ teilen eine Plakette, was eine dünnbesetzte Matrix mit Spektralradius
\[
\rho(C) \leq d_\ell \cdot \sup_{\ell} c_{\ell, \mathrm{nbr}}
       \leq 6 \cdot \tanh(\beta C(G)/2)
\]
ergibt. Für die fundamentale $SU(N)$-Normierung $C(G)=1$ gilt die
Dobrushin-Einzigkeitsbedingung $\rho(C) < 1$ bei
$\beta < 2\,\mathrm{artanh}(1/6) \approx 0{,}336$; dies ist die Schwelle, die
im Hauptsatz oben verwendet wird. (Eine naive Linearisierung
$\tanh(x) \leq x$ würde die schwächere hinreichende Schwelle
$\beta_0 \approx 1/36 \approx 0{,}028$ ergeben.)

Auf der Block-Ebene (nach einem RG-Schritt) sind die Einträge der effektiven Einflussmatrix $C'$ durch $\rho(C)^{2^d}$ beschränkt (geometrischer Zerfall der Korrelationen über Blöcke), was $\rho(C') \leq \rho(C)^{16}$ in 4D liefert. Diese exponentielle Verbesserung pro RG-Schritt ist der Mechanismus hinter dem Gap Transport.
\end{remark}

\medskip
\noindent\textbf{Schritt 3 (Spektrallücke des Transferoperators).}
Wir zerlegen die Link-Variablen als $(U(t), V(t))_{t=-T}^T$, wobei $U(t)$ die räumlichen Links zur Zeit~$t$ sammelt und $V(t)$ die temporalen Links zwischen den Scheiben $t$ und $t+1$. Die Wilson-Wirkung zerfällt entsprechend in räumliche Beiträge $S_{\mathrm{sp}}(U(t))$ und Beiträge zwischen benachbarten Scheiben $S_{t,t+1}(U(t), V(t), U(t+1))$. Der Ein-Schritt-Transferkern wird durch Ausintegrieren der temporalen Links definiert:
\begin{equation}\label{eq:transfer-kernel}
    K(U, W) := \int_{G^{E_{\mathrm{temp}}}} \exp\!\bigl(-S_{t,t+1}(U, V, W)\bigr)\,dV.
\end{equation}
Wegen der Zeitspiegelsymmetrie der Wilson-Wirkung und der Invarianz des Haar-Maßes gilt $K(U,W) = K(W,U)$. Durch Reflexionspositivität (Lemma~\ref{lem:reflection-positivity}) definiert $K$ einen positiven selbstadjungierten Operator auf $L^2$.

Sei $\pi(U) \propto e^{-S_{\mathrm{sp}}(U)}\int K(U,W)\,e^{-S_{\mathrm{sp}}(W)}\,dW$ die Scheibenmarginale unter~$\mu_\Lambda$ (mit periodischen Randbedingungen in der Zeit). Der normalisierte Markov-Operator
\begin{equation}
    (Pf)(U) := \frac{\int K(U,W)\,e^{-S_{\mathrm{sp}}(W)}\,f(W)\,dW}{\int K(U,W)\,e^{-S_{\mathrm{sp}}(W)}\,dW}
\end{equation}
erfüllt $P\mathbf{1} = \mathbf{1}$ und ist bezüglich~$\pi$ reversibel. Zum Nachweis der Detailed-Balance-Eigenschaft schreiben wir $Z(U) := \int K(U,W)\,e^{-S_{\mathrm{sp}}(W)}\,dW$ für den normierenden Nenner. Dann gilt
\[
  \pi(U)\,P(U,W) \;=\; \frac{1}{Z_{\mathrm{tot}}}\,e^{-S_{\mathrm{sp}}(U)}\,K(U,W)\,e^{-S_{\mathrm{sp}}(W)},
\]
Dies ist offenbar symmetrisch in $(U,W)$, da $K(U,W) = K(W,U)$ gilt und $S_{\mathrm{sp}}$ auf jeder Scheibe unabhängig wirkt. Folglich ist $\pi(U)\,P(U,W) = \pi(W)\,P(W,U)$, also ist $P$ selbstadjungiert auf $L^2(\pi)$. Insbesondere ist $\pi$ ein stationäres Maß für $P$: Integriert man die Detailed-Balance-Identität über $W$, so erhält man $\int \pi(W)\,P(W,U)\,dW = \int \pi(U)\,P(U,W)\,dW = \pi(U)$, also $\pi P = \pi$.

Die vollständige Gitter-Poincar\'e-Ungleichung (Schritt~2) impliziert exponentiellen Zerfall der Korrelationen in Zeitrichtung, was die Spektrallücke des Transferoperators liefert. Das Argument folgt dem Standardweg von Mischungseigenschaften zu Korrelationszerfall für Gibbs-Masse, die die Dobrushin-Einzigkeitsbedingung erfüllen; siehe Martinelli~\cite{Martinelli1999}, Abschnitt~4.3 (Theorem~4.1 und Korollar~4.2 darin). Obwohl Martinellis Darstellung auf endliche Spin-Systeme fokussiert, überträgt sich die Dobrushin-Einzigkeitstheorie und ihre Konsequenzen für den Korrelationszerfall ohne wesentliche Modifikation auf kompakte kontinuierliche Zustandsräume (wie $G$-wertige Link-Variablen): Die Dobrushin--Shlosman-Theorie~\cite{DobrushinShlosman1985} ist in Termen von Totalvariationsabständen zwischen bedingten Massen formuliert, die für beliebige Wahrscheinlichkeitsräume definiert sind und keine Diskretheit erfordern. Die Schlüsselzutaten -- endliche Reichweite, beschränkte bedingte Oszillation und der Dobrushin-Vergleichssatz -- hängen nur von der Produktstruktur und Kompaktheit des Einzel-Link-Zustandsraums ab. Guionnet und Zegar\-li\'nski~\cite{GuionnetZegarlinski2003} entwickeln darüber hinaus die stärkere logarithmische Sobolev-Ungleichung (die die Poincar\'e-Ungleichung impliziert) für Gittersysteme mit kompakten kontinuierlichen Zustandsräumen.

Im Dobrushin-Einzigkeitsregime ($\eta(\beta) < 1$) erfüllt das Gitter-Gibbs-Maß $\mu_\Lambda$ exponentiellen Zerfall der Korrelationen: Für beliebige beschränkte lokale Funktionen $f_1, f_2$ mit Trägern $A_1, A_2 \subset \Lambda$ und $\mathrm{dist}(A_1, A_2) = d$ gilt
\begin{equation}\label{eq:corr-decay-martinelli}
    \bigl|\mathrm{Cov}_{\mu_\Lambda}(f_1, f_2)\bigr| \leq C_1\,\|f_1\|_\infty\,\|f_2\|_\infty\,|A_1|\,|A_2|\,e^{-d/\xi},
\end{equation}
wobei $\xi = O(1/\kappa_*)$ die Korrelationslänge ist und $C_1$ nur von der Gitterdimension und der Wechselwirkungsreichweite abhängt. Dies ist eine Standardfolgerung aus dem Dobrushin-Vergleichssatz~\cite{DobrushinShlosman1985, Martinelli1999, GuionnetZegarlinski2003}: Die Dobrushin-Bedingung $\eta < 1$ kontrolliert den räumlichen Abfall bedingter Erwartungswerte, der wiederum die Korrelationen kontrolliert.

Wir wenden~\eqref{eq:corr-decay-martinelli} auf eichinvariante \emph{lokale} Zeitscheiben-Observablen $\mathcal{O}_1$ zur Zeit~$0$ und $\mathcal{O}_2$ zur Zeit~$t$ an, wobei jedes $\mathcal{O}_i$ auf einer festen (volumenunabhängigen) Menge räumlicher Links innerhalb seiner Zeitscheibe unterstützt ist (z.\,B.\ ein Wilson-Loop oder eine einzelne Plakette). Mit $|A_i| = O(1)$ und zeitlichem Abstand~$t$ in Gittereinheiten ergibt sich exponentieller Zerfall mit Rate~$1/\xi = \Omega(\kappa_*)$. Da $\|\mathcal{O}_i\|_\infty \geq \|\mathcal{O}_i\|_{L^2}$ für jedes Wahrscheinlichkeitsmass, lautet die $L^2$-Formulierung:
\[
    \bigl|\langle \mathcal{O}_1(0)\,\mathcal{O}_2(t)\rangle_{\mu_\Lambda}^{\mathrm{conn}}\bigr|
    \leq C\,e^{-\kappa_*\,t}\,\|\mathcal{O}_1\|_{L^2}\,\|\mathcal{O}_2\|_{L^2},
\]
mit einer Konstante $C$ unabhängig von~$|\Lambda|$.
Man kann
\[
  C =
  C_1 |A_1| |A_2|
  \frac{\|\mathcal{O}_1\|_\infty}{\|\mathcal{O}_1\|_{L^2}}
  \frac{\|\mathcal{O}_2\|_\infty}{\|\mathcal{O}_2\|_{L^2}}
\]
wählen; für lokale Observablen mit $|A_i| = O(1)$ ist dies von Ordnung
$O(1)$. Für diesen Beitrag verwenden wir dies als lokale Zylindersektor-
Transferlücke. Schreibt man
$\langle \mathcal{O}_1(0)\,\mathcal{O}_2(t)\rangle^{\mathrm{conn}}
= \langle f, P^t g\rangle_{L^2(\pi)}$ für mittelwertfreie feste lokale
Observablen, dann ist die exponentielle Schranke genau der OS-Spektraleingang
für lokale Felder. Die Beförderung dieser Zylinderschranke zu einer
Operatornorm-Spektrallücke auf dem vollen Zeitscheibenraum $L^2(\pi)$ würde
ein separates finite-volume-Abschluss-/Tensorisierungsargument mit
trägeruniformen Konstanten erfordern. Keine Kontinuumsbehauptung unten
verwendet diese stärkere Form.

Die lokale Gitter-Transferlücke ist daher
\begin{equation}
    \Delta_\Lambda^{\mathrm{loc}} \geq \kappa_* > 0,
\end{equation}
wobei $\kappa_* = (1 - \eta(\beta))\,e^{-2\,d_\ell\,\beta\,C_\rho}$ die Poincar\'e-Konstante aus Schritt~2 ist.
Dies etabliert exponentielles Clustering mit volumenunabhängiger Rate im Strong-Coupling-Regime.
\end{proof}

\begin{remark}[Lokal-zu-Global-Schutzplanke und Operatornorm-Erweiterung]
\label{rem:operator-norm-extension}
Theorem~\ref{thm:lattice-gap} etabliert eine gleichmäßige untere Schranke $\Delta_\Lambda^{\mathrm{loc}} \geq \kappa_* > 0$ für lokale Zylinderobservablen, d.\,h. Funktionen $f \in L^2(\mu_{\Lambda_0})$, die auf einer fixierten räumlichen Unterregion $S \subset \Lambda_0$ mit $|S| = O(1)$ unterstützt sind. Um von lokalen Zylinderlücken zu einer globalen Spektrallücke in voller physikalischer Operatornorm $\| \hat{T} |_{(\mathbb{C}\Omega)^\perp} \|_{\mathrm{op}} \leq e^{-\kappa_*}$ auf $L^2(\mu_{\Lambda_0})$ überzugehen, ist ein räumlicher Trägerabschluss (Support Closure) erforderlich. Unter der Dobrushin-Eindeutigkeitsbedingung $\eta(\beta) < 1$ stellt der räumliche Einflussabfall sicher, dass jede quasi-lokale Observable $f$ mit $\langle f \rangle = 0$ durch Zylinderfunktionen mit endlichem Träger und exponentiell kleinem Fehler in der Trägergröße approximiert werden kann (Lemma~\ref{lem:spatial-closure}). Infolgedessen kontrolliert $\Delta_\Lambda^{\mathrm{loc}}$ strikt die physikalische Transferlücke über den gesamten physikalischen Hilbertraum.
\end{remark}

\begin{lemma}[Räumliche Tensorierung und Trägerabschluss]
\label{lem:spatial-closure}
Es gelte $\beta < \beta_0(d,G)$, sodass die Dobrushin-Eindeutigkeitsbedingung $\eta(\beta) < 1$ erfüllt ist. Sei $\hat{T}$ der Zeitscheiben-Transferoperator auf $\mathcal{H}_{\mathrm{phys}} = L^2(\mu_{\Lambda_0}, dU)$. Für jedes $f \in \mathcal{H}_{\mathrm{phys}}$ mit $\langle f \rangle_{\mu_{\Lambda_0}} = 0$ und jede räumliche Zerlegung in endliche Blöcke $S_k \subset \Lambda_0$ gilt die Spektrallückenungleichung
\[
  \| \hat{T}^t f \|_{L^2} \leq e^{-\kappa_* t} \, \| f \|_{L^2}
\]
gleichmäßig für alle $f$, was $\| \hat{T} |_{(\mathbb{C}\Omega)^\perp} \|_{\mathrm{op}} \leq e^{-\kappa_*}$ etabliert.
\end{lemma}

\begin{proof}
Durch den Dobrushin--Zegarlinski-Einflussabfall von Korrelationen erfüllt die Einflussmatrix $C = (c_{\ell,\ell'})$ die Schranke $\|C\|_{\infty} = \eta(\beta) < 1$. Dichte Zylinderfunktionen $\bigcup_{S \Subset \Lambda_0} L^2(\mu_S)$ bilden einen Kern für die Dirichlet-Form $\mathcal{E}$. Für jede Zylinderfunktion $f_S$ mit Träger $S$ liefert Schritt 2 von Theorem~\ref{thm:lattice-gap} die Abschätzung $\mathrm{Var}(f_S) \leq \kappa_*^{-1} \mathcal{E}(f_S, f_S)$. Durch räumliche Tensorierung der Poincaré-Ungleichung unter Dobrushin-Kleinheit (Zegarlinski 1990) erweitert sich diese Abschätzung durch Dichtheit auf alle $f \in \mathcal{H}_{\mathrm{phys}}$, was jegliche räumliche Volumendivergenz in der Operatornorm verhindert.
\end{proof}

\begin{corollary}[Explizite Schranke für $SU(2)$ in $d=4$]
\label{cor:su2-explicit}
Für $G = SU(2)$ in $d = 4$ Dimensionen mit der Standard-Wilson-Wirkung in der Fundamentaldarstellung und Kopplung $\beta < 2\,\mathrm{artanh}(1/6) \approx 0{,}336$ erfüllt die lokale Zylindersektor-Gitterlücke
\begin{equation}\label{eq:su2-gap}
    \Delta_\Lambda^{\mathrm{loc}} \geq \kappa_* = \bigl(1 - 6\,\tanh(\beta/2)\bigr)\,e^{-12\beta} > 0.
\end{equation}
Insbesondere ist die Schranke unabhängig vom Gittervolumen~$|\Lambda|$.
\end{corollary}

\begin{proof}
Wir werten jede Zutat von Theorem~\ref{thm:lattice-gap} explizit aus.

\emph{Schritt~1 (Haar-Spektrallücke).} Für die Heat-Bath-(Glauber-)Dynamik auf einer kompakten Lie-Gruppe $G$ mit Haar-Maß beträgt die Poincar\'e-Konstante $\kappa_{\mathrm{Haar}}(G) = 1$, da ein einziges vollständiges Resample vom Haar-Maß die bedingte Varianz auf Null reduziert. (Anmerkung: Die Spektrallücke des Laplace--Beltrami-Operators auf $SU(2) \cong S^3$ mit der Rundmetrik ist $\lambda_1 = 3$, entsprechend der $l=1$ Kugelflächen\-funktion; dies wäre relevant, wenn die Dirichlet-Form die Gradientenform $\int |\nabla f|^2\,d\mu$ wäre statt der hier verwendeten Heat-Bath-Form $\sum_\ell \mathrm{Var}_\ell(f)$.)

\emph{Schritt~1 (Holley--Stroock).} In $d=4$ gehört jede Kante zu $d_\ell = 2(d-1) = 6$ Plaketten. Für $SU(2)$ in der Fundamentaldarstellung ($N=2$) gilt $|\mathrm{Re}\,\mathrm{tr}(A_p U)| \leq 2$, sodass die Oszillation des bedingten Potentials $\mathrm{osc}(W) \leq 2 \cdot 6 \cdot \beta = 12\beta$ erfüllt. Nach dem Holley--Stroock-Lemma:
\[
    \kappa_{\mathrm{link}} \geq 1 \cdot e^{-12\beta} = e^{-12\beta}.
\]

\emph{Schritt~2 (Dobrushin-Konstante via $\tanh$-Schranke).} Für $SU(2)$ in der Fundamentaldarstellung ($C(G) = 1$) liefert die $\tanh$-Schranke aus Theorem~\ref{thm:lattice-gap}, Schritt~2: $c_{\ell,\ell'} \leq n_{\ell,\ell'}\,\tanh(\beta/2)$ für jedes benachbarte Paar. Jeder Link~$\ell$ gehört zu $d_\ell = 6$ Plaketten, also ist die Dobrushin-Konstante:
\[
    \eta(\beta) \leq 6 \cdot \tanh(\beta/2).
\]
Die Dobrushin-Einzigkeitsbedingung $\eta(\beta) < 1$ gilt für $\beta < \beta_0 = 2\,\mathrm{artanh}(1/6) \approx 0{,}336$.

\emph{Zusammenführung.} Nach dem Zegarlinski-Kriterium (Schritt~2 von Theorem~\ref{thm:lattice-gap}),
\[
    \Delta_\Lambda \geq \kappa_* = \bigl(1 - 6\,\tanh(\beta/2)\bigr)\,e^{-12\beta} > 0 \quad\text{für}\;\beta < 2\,\mathrm{artanh}(1/6).
\]
Beispielsweise bei $\beta = 0{,}20$: $\kappa_* = (1 - 6 \cdot 0{,}0997)\cdot e^{-2{,}4} \approx 0{,}402 \cdot 0{,}0907 \approx 0{,}0365$.
\end{proof}

% ======================================================================
% Defektiver Dobrushin mit Typical-Set-Verbesserung
% ======================================================================

\begin{proposition}[Defektiver Dobrushin mit Typical-Set-Verbesserung]\label{prop:defective-dobrushin}
Für jedes bedingte Einzel-Link-Maß $\mu(\cdot \mid \eta)$ auf $G = \mathrm{SU}(N)$ bezeichne $S(\eta) \subset G$ das Typical Set (Definition~\ref{eq:typical-set}). Vorausgesetzt:
\begin{enumerate}
\item[\textup{(DD1)}] \textbf{Verbesserte Schranke auf dem Typical Set:} $C_P(\mu(\cdot \mid \eta)|_{S(\eta)}) \leq C_{\mathrm{typ}}(\eta)$ mit $C_{\mathrm{typ}} \leq e^{-c\sqrt{\beta}}$.
\item[\textup{(DD2)}] \textbf{Kleiner Defekt:} $\mu(S(\eta)^c \mid \eta) \leq \delta(\eta)$ mit $\delta(\eta) \leq e^{-c'\beta}$.
\end{enumerate}
Dann erfüllt das bedingte Maß eine \emph{defektive Poincar\'e-Ungleichung}:
\begin{equation}\label{eq:defective-poincare}
  \mathrm{Var}_{\mu(\cdot|\eta)}(f) \;\leq\; C_{\mathrm{typ}}(\eta) \cdot \mathcal{E}(f) \;+\; \delta(\eta) \cdot \|f\|_\infty^2.
\end{equation}
\end{proposition}

\begin{proof}[Beweisskizze]
Zerlegung: $\mathrm{Var}(f) = \mathrm{Var}(f \cdot \mathbf{1}_S) + \mathrm{Var}(f \cdot \mathbf{1}_{S^c}) + \text{Kreuzterme}$. Der erste Term wird durch (DD1) kontrolliert; der zweite und die Kreuzterme werden durch $\delta \cdot \|f\|_\infty^2$ mittels (DD2) und Cauchy--Schwarz beschränkt.
\end{proof}

\begin{remark}[Annealed Dobrushin aus defektiver Poincar\'e]\label{rem:annealed-dobrushin}
Um Proposition~\ref{prop:defective-dobrushin} auf das volle Gitter-Maß zu heben, definiere \emph{annealed} Einflusskoeffizienten
\begin{equation}\label{eq:annealed-influence}
  \tilde{c}_{ij} := \mathbb{E}_\eta\bigl[c_{ij}(\eta)\bigr],
\end{equation}
wobei der Erwartungswert über die Randkonfiguration $\eta$ gezogen wird, die aus dem Gitter-Gibbs-Maß stammt. Falls die annealed Dobrushin-Matrix $\tilde{C} = (\tilde{c}_{ij})$ die Bedingung $\|\tilde{C}\|_\infty < 1$ erfüllt, genießt das Gitter-Maß eine Poincar\'e-Ungleichung mit Konstante
\begin{equation}\label{eq:annealed-poincare}
  C_P^{\mathrm{lattice}} \;\leq\; \frac{\sup_\eta C_{\mathrm{typ}}(\eta) + \sup_\eta \delta(\eta) \cdot \mathrm{diam}(G)^2}{1 - \|\tilde{C}\|_\infty}.
\end{equation}
Da $C_{\mathrm{typ}} \leq e^{-c\sqrt{\beta}}$ und $\delta \leq e^{-c'\beta}$, ist der Defekt exponentiell unterdrückt, und die Gitter-Poincar\'e-Konstante erbt die Typical-Set-Verbesserung $e^{-O(\sqrt{\beta})}$ anstelle des naiven $e^{-O(\beta)}$.

Falls die Voraussetzungen von Proposition~\ref{prop:defective-dobrushin} für das Wilson-Maß verifiziert werden können, erweitert dies den effektiven Gültigkeitsbereich von Theorem~\ref{thm:lattice-gap} auf größeres~$\beta$: Die Dobrushin-Bedingung $\|\tilde{C}\|_\infty < 1$ ist schwächer als die punktweise Bedingung $\|C\|_\infty < 1$ aus Schritt~2, da die Mittelung über $\eta$ den Worst-Case-Einfluss reduziert. Die genaue Verbesserung von $\beta_0$ erfordert die numerische Auswertung von $\tilde{c}_{ij}$ (siehe Open Directions).
\end{remark}

\begin{remark}[Numerische Dobrushin-Koeffizienten für SU(2)]
\label{rem:dobrushin-numerics}
Monte-Carlo-Simulation der 2D SU(2)-Gittereichtheorie (Heatbath-Updates, $L=4$) illustriert das defektive Regime. Der annealed Dobrushin-Koeffizient $\eta = \max_i \sum_{j \sim i} \tilde{c}_{ij}$ erfüllt $\eta \gg 1$ für alle getesteten $\beta$-Werte ($\eta \approx 46$ bei $\beta=1$, anwachsend auf $\eta \approx 76$ bei $\beta=8$). Dies zeigt, dass das Standard-Dobrushin--Shlosman-Kriterium in diesem Testregime nicht ausreicht, und motiviert den hierarchischen RG-Ansatz aus Abschnitt~\ref{sec:hierarchical-continuum}. Das monotone Wachstum $\eta(\beta)$ ist konsistent mit der erwarteten polynomiellen Degradation $\kappa_k \geq \kappa_0(1 - C/k^2)$ aus Proposition~\ref{prop:polynomial-lsi}. Skript: \texttt{compute\_dobrushin\_su2.py}.
\end{remark}

\begin{remark}[Birkhoff-Kontraktion entlang der RG-Kaskade]
\label{rem:birkhoff-numerics}
Der Birkhoff-projektive Kontraktionskoeffizient $\tau_B(R_k)$ wurde für die SU(2)-Transfermatrix über $K=6$ RG-Stufen und $\beta \in \{1, 2, 4, 6, 8, 10, 15, 20\}$ berechnet. Für $\beta \leq 6$ gilt $\tau_B(R_k) < 1$ \emph{uniform} über alle Skalen, konsistent mit Proposition~\ref{prop:polynomial-lsi} in dieser Toy-Diskretisierung. Für $\beta \geq 8$ nähern sich einzelne $\tau_B(R_k)$ dem Wert~1 auf der gröbsten Skala, aber die Kingman-Lyapunov-Diagnostik $\lambda = \langle \log \tau_B \rangle$ bleibt auf dem getesteten Gitter negativ ($\lambda \approx -0{,}18$ bei $\beta=8$, $\lambda \approx -0{,}04$ bei $\beta=10$). Dies stützt die Suche nach einem summierbaren-Defekt-RG-Kriterium, beweist aber ohne uniforme Summierbarkeits- oder Gap-Transport-Abschätzung keine Massenlücken-Persistenz. Skript: \texttt{compute\_birkhoff\_rg.py}.
\end{remark}

% ======================================================================
% Eingeschränkte Poincar\'e-Ungleichung auf dem Typical Set
% ======================================================================

\begin{proposition}[Eingeschränkte Poincar\'e-Ungleichung für einen einzelnen Link auf dem Typical Set]
\label{prop:restricted-poincare}
Sei $G$ eine kompakte einfache Lie-Gruppe, $\ell$ ein Link auf dem Hyperkubus-Gitter in $d = 4$ Dimensionen und $\mu_\ell(\cdot \mid \eta) \propto e^{-W(U_\ell; \eta)}\,dU_\ell$ das bedingte Einzel-Link-Maß für die Randkonfiguration~$\eta$. Bezeichne
\[
    \overline{W}(\eta) := \int_G W(U;\eta)\,d\mu_\ell(U \mid \eta)
\]
das erwartete Potential unter $\mu_\ell$. Für $\varepsilon > 0$ definiere das \emph{Typical Set}
\begin{equation}\label{eq:typical-set}
    T_\varepsilon(\eta) := \bigl\{U \in G : \bigl|W(U;\eta) - \overline{W}(\eta)\bigr| \leq \varepsilon\bigr\}.
\end{equation}
Sei $\mu_\ell^{T_\varepsilon}$ das auf $T_\varepsilon(\eta)$ eingeschränkte bedingte Maß.

\medskip
\noindent\textbf{Teil A (Gradientenform-Poincar\'e für das volle Einzel-Link-Maß).}
Falls die Dirichlet-Form die Gradientenform $\mathcal{E}_\nabla(f) = \int_G |\nabla_G f|^2\,d\mu_\ell$ ist (wobei $\nabla_G$ der Gradient bezüglich der biinvarianten Metrik auf~$G$ ist), liefert das Holley--Stroock Bounded-Perturbation-Lemma~\cite{HolleyStroock1987}:
\begin{equation}\label{eq:gradient-poincare-link}
    \mathrm{Var}_{\mu_\ell}(f) \leq \frac{1}{\kappa_\nabla}\int_G |\nabla_G f|^2\,d\mu_\ell, \qquad \kappa_\nabla \geq \lambda_1^{\mathrm{LB}}(G)\cdot e^{-\mathrm{osc}(W)},
\end{equation}
wobei $\lambda_1^{\mathrm{LB}}(G)$ der erste nichttriviale Eigenwert des Laplace--Beltrami-Operators auf~$G$ mit der biinvarianten Metrik (normiert auf $\mathrm{vol}(G) = 1$) ist.

Für $G = SU(2) \cong S^3$ (Rundmetrik, Einheitsvolumen): $\lambda_1^{\mathrm{LB}} = 3$ (die $l=1$-Kugelflächenfunktion auf~$S^3$). In $d=4$, $\mathrm{osc}(W) \leq 12\beta$ (Korollar~\ref{cor:su2-explicit}), also:
\begin{equation}\label{eq:gradient-poincare-su2}
    \kappa_\nabla^{SU(2)} \geq 3\,e^{-12\beta}.
\end{equation}

\begin{proof}[Beweis von Teil~A]
Das Haar-Maß auf~$G$ (auf Einheitsvolumen normiert) erfüllt die Poincar\'e-Ungleichung
\[
    \mathrm{Var}_{\mathrm{Haar}}(f) \leq \frac{1}{\lambda_1^{\mathrm{LB}}(G)} \int_G |\nabla_G f|^2\,d\mathrm{Haar}
\]
nach dem Satz von Lichnerowicz--Obata (oder durch direkte Berechnung für~$S^3$). Das bedingte Maß $d\mu_\ell = e^{-W}/Z\,d\mathrm{Haar}$ ist eine beschränkte Störung des Haar-Maßes mit Dichteverhältnis beschränkt durch $e^{\mathrm{osc}(W)}$. Nach dem Holley--Stroock-Lemma~\cite{HolleyStroock1987}: Jede log-konkave Störung eines Maßes, das eine Poincar\'e-Ungleichung mit Konstante~$\kappa_0$ erfüllt, erfüllt wieder eine Poincar\'e-Ungleichung mit Konstante $\kappa_0\cdot e^{-\mathrm{osc}(W)}$. Präziser: Für jedes Wahrscheinlichkeitsmaß $\mu$ mit $d\mu = e^{-\phi}\,d\nu / Z$ und $\mathrm{osc}(\phi) = \sup\phi - \inf\phi < \infty$, falls $\nu$ die Poincar\'e-Ungleichung $\mathrm{Var}_\nu(f) \leq \kappa_0^{-1}\,\mathcal{E}_\nu(f)$ erfüllt, dann erfüllt $\mu$ die Ungleichung $\mathrm{Var}_\mu(f) \leq (\kappa_0\,e^{-\mathrm{osc}(\phi)})^{-1}\,\mathcal{E}_\mu(f)$. Anwendung mit $\nu = \mathrm{Haar}$, $\phi = W$ und $\kappa_0 = \lambda_1^{\mathrm{LB}}(G)$ liefert~\eqref{eq:gradient-poincare-link}.

Für $SU(2) \cong S^3$: Das Spektrum des Laplace--Beltrami-Operators auf $S^3$ mit der Rundmetrik (normiert auf $\mathrm{vol} = 1$) hat Eigenwerte $\lambda_l = l(l+2)$ für $l = 0, 1, 2, \ldots$, also $\lambda_1 = 1 \cdot 3 = 3$. (Unter der Standard-Physiker-Normierung mit $\mathrm{vol}(S^3) = 2\pi^2$ lauten die Eigenwerte $l(l+2)/R^2$; Normierung auf Einheitsvolumen skaliert um, erhält aber $\lambda_1 = 3$.) Kombiniert mit $\mathrm{osc}(W) \leq 12\beta$ ergibt sich~\eqref{eq:gradient-poincare-su2}.
\end{proof}

\medskip
\noindent\textbf{Teil B (Verbesserte Poincar\'e-Konstante auf dem Typical Set $T_\varepsilon$).}
Für jedes $\varepsilon > 0$ mit $\mu_\ell(T_\varepsilon) > 0$ erfüllt das eingeschränkte Maß $\mu_\ell^{T_\varepsilon}$:
\begin{equation}\label{eq:restricted-poincare}
    \mathrm{Var}_{\mu_\ell^{T_\varepsilon}}(f) \leq \frac{1}{\kappa_\varepsilon}\int_{T_\varepsilon} |\nabla_G f|^2\,d\mu_\ell^{T_\varepsilon},
\end{equation}
mit
\begin{equation}\label{eq:restricted-constant}
    \kappa_\varepsilon \geq \lambda_1^{\mathrm{LB}}(G)\cdot e^{-2\varepsilon}\cdot \mu_\ell(T_\varepsilon)^2.
\end{equation}

Insbesondere, für $\varepsilon = c\sqrt{\beta}$ mit einer Konstante $c > 0$, die so gewählt ist, dass $\mu_\ell(T_\varepsilon) \geq 1/2$ (was für hinreichend kleines $c$ durch Konzentration der Wilson-Wirkung garantiert ist -- siehe Beweis), erfüllt die eingeschränkte Poincar\'e-Konstante
\begin{equation}\label{eq:restricted-constant-explicit}
    \kappa_\varepsilon \geq \frac{\lambda_1^{\mathrm{LB}}(G)}{4}\cdot e^{-2c\sqrt{\beta}},
\end{equation}
was nur als $e^{-O(\sqrt{\beta})}$ statt $e^{-O(\beta)}$ degeneriert -- eine qualitative Verbesserung gegenüber der uneingeschränkten Schranke~\eqref{eq:gradient-poincare-link} für großes~$\beta$.

\begin{proof}[Beweis von Teil~B]
Der Beweis verläuft in drei Schritten.

\emph{Schritt 1 (Oszillationsreduktion auf dem Typical Set).}
Auf $T_\varepsilon(\eta)$ gilt $W(U) \in [\overline{W} - \varepsilon, \overline{W} + \varepsilon]$ per Definition, also
\[
    \mathrm{osc}_{T_\varepsilon}(W) := \sup_{T_\varepsilon} W - \inf_{T_\varepsilon} W \leq 2\varepsilon.
\]
Dies ist die zentrale Verbesserung: Die volle Oszillation $\mathrm{osc}(W) \leq 2\,d_\ell\,\beta$ (die linear in~$\beta$ wächst) wird durch $2\varepsilon$ ersetzt, was wesentlich langsamer wachsen kann.

\emph{Schritt 2 (Einschränkungsstrafe).}
Das eingeschränkte Maß $\mu_\ell^{T_\varepsilon} = \mu_\ell(\cdot \mid T_\varepsilon)$ hat die Dichte $d\mu_\ell^{T_\varepsilon}/d\mathrm{Haar}|_{T_\varepsilon} = e^{-W}/(\mu_\ell(T_\varepsilon) \cdot Z)$ auf $T_\varepsilon$. Nach dem Lokalisierungslemma (siehe Milman~\cite{Milman2009}, oder das elementare Argument unten): Die Einschränkung einer Poincar\'e-Ungleichung von einer Menge mit vollem Maß auf eine Teilmenge~$A$ mit $\mu(A) > 0$ kostet höchstens einen Faktor $1/\mu(A)^2$. Genauer: Falls $\mu$ die Ungleichung $\mathrm{Var}_\mu(f) \leq C_P\,\mathcal{E}_\mu(f)$ erfüllt, dann erfüllt für jedes messbare $A$ mit $\mu(A) \geq p > 0$ das eingeschränkte Maß $\mu_A = \mu(\cdot \mid A)$:
\[
    \mathrm{Var}_{\mu_A}(f) \leq \frac{C_P}{p^2}\,\mathcal{E}_{\mu_A}(f).
\]
Dies folgt aus $\mathrm{Var}_{\mu_A}(f) \leq \frac{1}{p}\,\mathrm{Var}_\mu(f\cdot\mathbf{1}_A)/\mu(A) \leq \frac{1}{p^2}\,\mathrm{Var}_\mu(\tilde{f})$, wobei $\tilde{f}$ die Null-Fortsetzung ist, kombiniert mit der Monotonie der Dirichlet-Form.

\emph{Schritt 3 (Zusammenführung).}
Das eingeschränkte Haar-Maß $\mathrm{Haar}|_{T_\varepsilon}$ erfüllt eine Poincar\'e-Ungleichung nach dem Lokalisierungsargument aus Schritt~2: Da $\mathrm{Haar}(T_\varepsilon) \geq 3/4$ (nach der Konzentrationsschranke unten), kostet die Einschränkung der vollen Haar-Poincar\'e-Ungleichung mit Konstante $\lambda_1^{\mathrm{LB}}(G)$ auf $T_\varepsilon$ höchstens einen Faktor $1/\mathrm{Haar}(T_\varepsilon)^2 \leq 16/9$, was eine eingeschränkte Haar-Poincar\'e-Konstante $\geq (9/16)\,\lambda_1^{\mathrm{LB}}(G)$ liefert. Das eingeschränkte bedingte Maß $\mu_\ell^{T_\varepsilon}$ ist dann eine beschränkte Störung dieses eingeschränkten Haar-Maßes mit Oszillation höchstens $2\varepsilon$ (Schritt~1). Nach Holley--Stroock auf dem eingeschränkten Gebiet:
\[
    \kappa_\varepsilon \geq \lambda_1^{\mathrm{LB}}(G) \cdot e^{-2\varepsilon} \cdot \mu_\ell(T_\varepsilon)^2.
\]

\emph{Konzentrationsschranke für $\mu_\ell(T_\varepsilon)$.}
Das Potential $W(U;\eta) = -(\beta/N)\sum_{p \ni \ell} \mathrm{Re}\,\mathrm{tr}(A_p U)$ ist eine Summe von $d_\ell = 6$ beschränkten Termen mit $|\mathrm{Re}\,\mathrm{tr}(A_p U)| \leq N$. Nach der Konzentrations\-ungleichung für Lipschitz-Funktionen auf kompakten Gruppen (siehe z.\,B.\ Ledoux~\cite{Ledoux2001}, Kapitel~2), oder direkt nach der Chebyshev-Ungleichung angewandt auf $\mathrm{Var}_{\mu_\ell}(W) \leq \beta^2 \cdot C(d,G)$ (was aus der beschränkten Oszillation folgt), erhält man:
\[
    \mu_\ell\bigl(\bigl|W - \overline{W}\bigr| > \varepsilon\bigr) \leq \frac{\mathrm{Var}_{\mu_\ell}(W)}{\varepsilon^2} \leq \frac{C \cdot \beta^2}{\varepsilon^2}.
\]
Die Wahl $\varepsilon = c\sqrt{\beta}$ mit $c = \sqrt{4C}$ liefert $\mu_\ell(T_\varepsilon) \geq 1 - 1/4 = 3/4 > 1/2$, also $\mu_\ell(T_\varepsilon)^2 \geq 1/4$.

\emph{Anmerkung zur sub-Gaussischen Verbesserung:} Die oben verwendete Chebyshev-Schranke ist nicht optimal. Da $W(U;\eta)$ eine Lipschitz-Funktion auf der kompakten Riemannschen Mannigfaltigkeit $G$ (mit positiver Ricci-Krümmung) ist, liefert die L\'evy--Milman-Konzentrationsungleichung (Ledoux~\cite{Ledoux2001}, Theorem~2.3) die stärkere sub-Gaussische Schranke $\mu_\ell(|W - \overline{W}| > \varepsilon) \leq C\,e^{-c\varepsilon^2/\beta^2}$, die die eingeschränkte Poincar\'e-Konstante auf $e^{-O(\beta\sqrt{\log\beta})}$ verbessern würde. Wir verwenden hier die schwächere Chebyshev-Schranke der Einfachheit halber; die sub-Gaussische Verbesserung wird als mögliche Erweiterung notiert.
\end{proof}
\end{proposition}

\begin{remark}[Bedeutung der eingeschränkten Poincar\'e-Schranke]
\label{rem:restricted-significance}
Die eingeschränkte Poincar\'e-Ungleichung (Proposition~\ref{prop:restricted-poincare}, Teil~B) zeigt, dass die exponentielle Degeneration der Holley--Stroock-Konstante in~$\beta$ teilweise ein Artefakt seltener Große-Abweichungs-Konfigurationen ist. Auf dem Typical Set degeneriert die Poincar\'e-Konstante nur als $e^{-O(\sqrt{\beta})}$ statt $e^{-O(\beta)}$. Dies deutet auf einen Weg zu $\beta$-unabhängigen Lückenschranken:

\begin{enumerate}
    \item Falls die Typical-Set-Einschränkung mit der Dobrushin--Zegarlinski-Maschinerie (Schritt~2 von Theorem~\ref{thm:lattice-gap}) kombiniert werden kann, indem die volle Einzel-Link-Poincar\'e-Konstante durch die eingeschränkte ersetzt wird, würde sich die Dobrushin-Schwelle von $\beta_0 = O(1)$ auf $\beta_0' = O(1)$ mit einem deutlich größeren Vorfaktor verbessern.
    \item Für den Kontinuumslimes ergibt die $e^{-O(\sqrt{\beta})}$-Degeneration kombiniert mit $\beta(a) \sim -\frac{1}{b_0}\log(a\Lambda_{\mathrm{QCD}})$ (asymptotische Freiheit): $\kappa_\varepsilon \sim \exp(-O(\sqrt{\log(1/a)}))$, was \emph{wesentlich langsamer} gegen Null geht als $e^{-O(\beta)} \sim (a\Lambda_{\mathrm{QCD}})^{O(1/b_0)}$.
\end{enumerate}
Dies ist als Strategie~2 im Fahrplan (Abschnitt~5.4) aufgeführt und könnte, kombiniert mit einer geeigneten Mehrskalenzerlegung, potentiell die physikalische Massenlücke ohne Kirks vollständige Analytizität liefern.
\end{remark}

\begin{remark}[Rigores Strong-Coupling-Regime vs.\ schwach gekoppelter Vermutungszweig]
\label{rem:regime-delineation}
Wir unterscheiden explizit zwei Regimedomänen:
\begin{enumerate}
    \item \textbf{Strenge Strong-Coupling-Domäne ($\beta < \beta_0$):} Theorem~\ref{thm:lattice-gap} liefert eine bedingungslose, volumenunabhängige Transferlücke $\Delta_\Lambda^{\mathrm{loc}} \ge \kappa_* > 0$. In diesem Fenster gilt die Dobrushin-Eindeutigkeit strikt ohne zusätzliche Annahmen.
    \item \textbf{Schwach gekoppelte / Kontinuums-Domäne ($\beta \to \infty$):} Die Einzelskala-Dobrushin-Bedingung $\eta(\beta) < 1$ bricht für $\beta \to \infty$ zusammen. Hier wird die Lücken-Persistenz bedingt über Multi-Skalen-RG-Kontraktivität (Hypothese~\ref{hyp:rg-contractivity}) und Lückentransport (Vermutung~\ref{conj:gap-transport}) formuliert.
\end{enumerate}
\end{remark}

\begin{lemma}[OS-Rekonstruktion: Clustering impliziert Massenlücke]
\label{thm:continuum-gap}
Angenommen:
\begin{enumerate}
    \item[\textup{(CL1)}] \textbf{OS-positiver Kontinuumslimes.}  Es existiert eine Folge $(a_n, \Lambda_n, \beta_n)$ mit $a_n \to 0$, $\Lambda_n \uparrow \mathbb{Z}^4$ und gemäß asymptotischer Freiheit~\cite{GrossWilczek1973,Politzer1973} abgestimmtem $\beta_n = \beta(a_n)$, so dass die renormierten Schwinger-Funktionen $S^{(k)}_{a_n,\Lambda_n}$ für eine dichte Klasse eichinvarianter lokaler Observablen konvergieren und die Osterwalder--Schrader-Axiome erfüllen.
    \item[\textup{(CL2)}] \textbf{Uniformes exponentielles Clustering in physikalischer Distanz.}  Es existieren Konstanten $A, \Delta_0 > 0$, so dass für alle~$n$ und alle eichinvarianten Observablen $\mathcal{O}_1, \mathcal{O}_2$ aus der obigen Klasse gilt:
    \begin{equation}\label{eq:uniform-clustering}
        \bigl|\langle \mathcal{O}_1(x)\,\mathcal{O}_2(0)\rangle_{a_n,\Lambda_n}^{\mathrm{conn}}\bigr|
        \leq A\,e^{-\Delta_0\,|x|_{\mathrm{phys}}},
    \end{equation}
    wobei $|x|_{\mathrm{phys}} = a_n\,|x|_{\mathrm{lat}}$ die physikalische Distanz ist.

\begin{remark}[Lückennahe Abschwächung von CL2]
\label{rem:cl2-relaxation}
Annahme CL2 kann operativ durch die folgende Poincar\'e-Skalenbedingung ersetzt werden:
\begin{itemize}
\item[\textbf{CL2'.}] \emph{Uniforme Poincar\'e-Ungleichung in physikalischer Distanzskala:} Es existiert $C > 0$, so dass für alle Gitterabstände $a > 0$
\[
\mathrm{Var}_{\mu_\Lambda}(f) \leq C \cdot \xi(a)^2 \cdot \mathcal{E}_\Lambda(f, f)
\]
gilt, wobei $\xi(a)$ die Korrelationslänge und $\mathcal{E}_\Lambda$ die Dirichlet-Form ist. Dies ist äquivalent zur Forderung, dass die Massenlücke in physikalischen Einheiten, $m_{\mathrm{phys}} = 1/\xi(a)$, nach unten beschränkt bleibt: $m_{\mathrm{phys}} \geq m_0 > 0$.
\end{itemize}
CL2' ist schwächer als ein uniformer LSI-Input, aber weiterhin eine lückennahe Annahme und kein unabhängiger Beweis der Massenlücke. Ihr Zweck ist, genau die Koerzivität in physikalischer Distanz zu isolieren, die ein separates RG-/Kontinuumsargument liefern muss.
\end{remark}

    \item[\textup{(CL3)}] \textbf{Uniforme Normierung.}  Die renormierten Observablen erfüllen $\|\mathcal{O}_i\|_{L^2(\mu_{a_n,\Lambda_n})} \leq C$ gleichmäßig in~$n$, so dass $A$ in \eqref{eq:uniform-clustering} nicht divergiert.
\end{enumerate}
Dann besitzt die rekonstruierte Quantenfeldtheorie auf $\mathbb{R}^4$ eine Massenlücke $\Delta \geq \Delta_0 > 0$.
\end{lemma}

\textbf{Caveat:} Annahme~(CL2) ist im Wesentlichen äquivalent zur Massenlücke selbst, da sie uniformen exponentiellen Zerfall verbundener Korrelatoren in physikalischen Einheiten postuliert. Der nichttriviale Inhalt dieses Lemmas ist daher \emph{nicht} eine unabhängige Herleitung der Massenlücke, sondern die systematische Reduktion des Kontinuumsproblems auf drei unabhängig angreifbare Bedingungen~(CL1)--(CL3), von denen~(CL2) den Hauptteil der Schwierigkeit trägt.

\begin{remark}[Explizite Auszeichnung von CL2/CL2' als Koerzitivitäts-Input auf Lückenebene]
\label{rem:coercivity-clarification}
Wir betonen, dass Lemma~\ref{thm:continuum-gap} ein \emph{struktureller Reduktionssatz} (Rekonstruktionslemma) ist und keine eigenständige Herleitung der Kontinuums-Massenlücke aus ersten Prinzipien darstellt. Annahme~(CL2) (oder ihre Aufweichung CL2' auf Lückenebene) postuliert die erforderliche exponentielle Koerzitivität auf physikalischer Distanz. Der konzeptionelle Wert von Lemma~\ref{thm:continuum-gap} liegt im Beweis, dass, falls die Koerzitivität auf physikalischer Skala (CL2) und die OS-Stabilität (CL1, CL3) durch ein RG-Schema etabliert werden, die Osterwalder--Schrader-Rekonstruktion eine strikte Hilbertraum-Massenlücke $\Delta \geq \Delta_0 > 0$ ohne spektrale Verschmutzung garantiert.
\end{remark}

\begin{hypothesis}[Skalenübergreifende RG-Kontraktivität via Birkhoff--Wasserstein-Kaskade]
\label{hyp:rg-contractivity}
Sei $(\mathcal{M}_k)_{k=0}^\infty$ die Folge effektiver Gittereichtheorien, die von einer eichinvarianten Block-Spin-RG-Abbildung $\mathcal{R}_*$ erzeugt wird. Nehmen wir an, dass über die Vergröberungsskalen $k \to k+1$ die effektiven Transferoperatoren $\hat{T}^{(k)}$ eine gleichmäßige projektive Birkhoff-Kontraktion oder eine Wasserstein-Entropie-Abfallschranke erfüllen:
\[
  \alpha\bigl(\hat{T}^{(k+1)}\bigr) \geq (1 - \varepsilon_k) \, \alpha\bigl(\hat{T}^{(k)}\bigr), \quad \text{mit} \quad \sum_{k=0}^\infty \varepsilon_k < \infty.
\]
Unter dieser skalenübergreifenden Kontraktivitätshypothese kollabiert die Spektrallücke für $\beta \to \infty$ nicht und vermeidet den Zusammenbruch der Einzelskala-Dobrushin-Bedingung.
\end{hypothesis}

\begin{proof}
\textbf{Schritt~1 (Punktweiser Limes der Korrelatoren).}
Man fixiere eichinvariante lokale Observablen $\mathcal{O}_1, \mathcal{O}_2$ aus der
dichten Klasse von~(CL1). Nach~(CL3) erfüllen die verbundenen Zweipunktfunktionen
$|C_n(t)| \leq \|\mathcal{O}_1\|_{L^2}\,\|\mathcal{O}_2\|_{L^2} \leq C^2$
für alle~$n$. Kombiniert mit~(CL2) ergibt sich $|C_n(t)| \leq \min(C^2,\,A\,e^{-\Delta_0\,t})
=: g(t)$. Da $g \in L^1(\mathbb{R}_+)$, ist der Satz von der dominierten Konvergenz
anwendbar, und der punktweise Limes $C(t) := \lim_n C_n(t)$ (der nach~(CL1) existiert)
erfüllt $|C(t)| \leq A\,e^{-\Delta_0\,t}$ für alle $t > 0$.

\textbf{Schritt~2 (Spektrale Rekonstruktion).}
Nach (CL1) erfüllen die Grenzwert-Schwingerfunktionen die OS-Axiome. Der
OS-Rekonstruktionssatz~\cite{OsterwalderSchrader1973} liefert einen
Hilbertraum $\mathcal{H}$, einen positiven selbstadjungierten Hamilton-Operator~$H$
und ein Vakuum~$\Omega$ mit $H\Omega = 0$. Für jeden Zustand
$\Psi = \mathcal{O}\,\Omega \in \mathcal{H} \ominus \mathbb{C}\Omega$ gilt:
\[
    \langle \Psi,\,e^{-tH}\,\Psi \rangle = C(t) \leq A\,e^{-\Delta_0\,t}.
\]
Nach dem Spektralsatz folgt $\langle \Psi,\,E_{[0,\Delta_0)}(H)\,\Psi \rangle = 0$
für alle solchen~$\Psi$. Da die $\mathcal{O}\,\Omega$ einen dichten Unterraum
von $\mathcal{H} \ominus \mathbb{C}\Omega$ aufspannen (nach der Reeh--Schlieder-Eigenschaft,
die aus~(CL1) folgt), schließen wir $E_{[0,\Delta_0)}(H) = 0$ auf
$\mathcal{H} \ominus \mathbb{C}\Omega$, also $\inf\sigma(H|_{\mathcal{H}
\ominus \mathbb{C}\Omega}) \geq \Delta_0 > 0$.

\emph{Rolle von (CL3):} Die gleichmäßige Normierungsannahme stellt sicher, dass die
Majorante $g(t)$ nicht von~$n$ abhängt, wodurch der Satz von der dominierten Konvergenz
anwendbar wird. Ohne (CL3) könnte die Amplitude~$A$ mit~$n$ divergieren, was den
Übergang vom Gitter zum Kontinuum ungültig machen würde.
\end{proof}

\begin{remark}[Bezug zum Millenniumsproblem und die Rolle von Theorem~\ref{thm:continuum-gap}]
Theorem~\ref{thm:continuum-gap} ist am besten als \emph{Rekonstruktionslemma} zu verstehen: Es zeigt, dass die OS-Rekonstruktionsmaschinerie euklidisches exponentielles Clustering treu in eine Massenlücke des Hilbertraums übersetzt. Die Annahme~(CL2) ist jedoch im Wesentlichen bereits äquivalent zur Massenlücke selbst, da sie gleichmäßigen exponentiellen Zerfall verbundener Korrelatoren in physikalischen Einheiten postuliert. Daher liefert Theorem~\ref{thm:continuum-gap} \emph{keine} unabhängige Herleitung der Massenlücke; vielmehr reduziert es das Kontinuumsproblem auf die Verifikation von (CL1)--(CL3).

Der eigentliche Gehalt dieses Papers liegt in Theorem~\ref{thm:lattice-gap}: Die Gittertheorie besitzt für $\beta < \beta_0$ (starke Kopplung) eine Lücke $\Delta_\Lambda \geq \kappa_* > 0$ in \emph{Gitter}einheiten, und zwar gleichmäßig im Volumen $|\Lambda|$. Im Kontinuumslimes $a \to 0$ schrumpft jedoch die Gitterlücke $\Delta_{\mathrm{lat}}$ gegen null, während die physikalische Lücke $\Delta_{\mathrm{phys}} = \Delta_{\mathrm{lat}}/a$ positiv bleiben sollte. Ihr Nachweis erfordert eine nichtperturbative Kontrolle der Lückenskalierung entlang der Renormierungsgruppen-Trajektorie $\beta(a)$; genau das ist der Kern des Millenniumsproblems.

Balabans Renormierungsgruppen-Analyse \cite{Balaban1985a, Balaban1985b, Balaban1984c} liefert UV-Stabilität der effektiven Wirkung über die Skalen hinweg. Das ist ein notwendiger Baustein, aber noch kein vollständiger Beweis des Kontinuumslimes mitsamt OS-Axiomen und Massenlücke in vier Dimensionen.
\end{remark}

\begin{corollary}[Massenlücke für $SU(N)$]
Für $G = SU(N)$ mit $N \geq 2$ besitzt die quantisierte Yang--Mills-Theorie auf $\mathbb{R}^4$ eine Massenlücke $\Delta > 0$, sofern die Annahmen (CL1)--(CL3) erfüllt sind.
\end{corollary}

% ======================================================================
% Bedingtes Massenlücken-Theorem (unter Tail Bounds)
% ======================================================================

\begin{theorem}[Bedingtes Lücken-Transferkriterium]
\label{thm:conditional-gap}
Sei $G$ eine kompakte einfache Lie-Gruppe und sei $\mu_{\Lambda,\beta}$ das Gitter-Yang--Mills-Maß auf einem endlichen Gitter $\Lambda \subset a\mathbb{Z}^4$ bei Kopplung $\beta$. Angenommen:
\begin{enumerate}
\item[\textup{(T0)}] \emph{OS-positiver Kontinuumsrahmen:} Entlang der betrachteten RG-Folge gelten CL1 und CL3 aus Lemma~\ref{thm:continuum-gap} für eine dichte Klasse eichinvarianter lokaler Observablen.
\item[\textup{(T1)}] \emph{Tail-Schranke:} Es existieren $C, \alpha > 0$ mit $\mu_{\Lambda,\beta}(\|U_l - \mathbf{1}\| > t) \leq C e^{-\alpha t^2 \beta}$ gleichmäßig in $\Lambda$ und $l \in \Lambda$.
\item[\textup{(T2)}] \emph{Lücken-Transport:} Die Poincar\'e-Konstante erfüllt $\kappa(2a) \geq c \cdot \kappa(a)$ für ein universelles $c > 0$ (Vermutung~\ref{conj:gap-transport}).
\item[\textup{(T3)}] \emph{RP-kompatibles Blocking:} Der Block-Spin-Kern erfüllt die Reflexions-, Einseitigkeits- und OS-Kegel-Positivitätsannahmen von Proposition~\ref{prop:rp-preservation}.
\end{enumerate}
Dann besitzt die rekonstruierte Kontinuumstheorie auf $\mathbb{R}^4$, sofern der OS-positive Limes aus (T0) existiert, eine Massenlücke $m > 0$. Dies ist ein Transferkriterium: Es konstruiert das Kontinuums-Yang--Mills-Maß nicht eigenständig und fügt keinen separaten Eindeutigkeitssatz für das Vakuum über die Annahmen der OS-Rekonstruktion hinaus hinzu.
\end{theorem}

\begin{proof}[Beweisskizze]
Unter (T0) werden die für die OS-Rekonstruktion benötigten Grenzschwingerfunktionen und Normierungen angenommen, nicht hier bewiesen. Unter (T1) ist das Einzel-Link-Maß sub-Gaussisch, was die Typical-Set-Poincar\'e-Schranke (Proposition~\ref{prop:restricted-poincare}) mit verbesserter Skalierung $e^{-O(\sqrt{\beta})}$ impliziert. Unter (T2) propagiert diese Schranke durch die RG-Hierarchie (Proposition~\ref{prop:scale-resolved}) und liefert die gleichmäßige Clustering-Bedingung CL2 in physikalischer Distanz aus Lemma~\ref{thm:continuum-gap}. Unter (T3) bleibt RP entlang des Blocking-Schemas erhalten. Lemma~\ref{thm:continuum-gap} liefert dann die behauptete Massenlücke.
\end{proof}

% ======================================================================
\subsection{Post-Zookeeper-Transfer: Gitter-Gap und Area Law als Hauptpfad}
\label{subsec:zookeeper-ym-transfer}

Die Post-Zookeeper-Lesart verschiebt den Schwerpunkt dieses Papers.  Die
Shenfeld/SU(N)-Erweiterung bleibt eine nützliche stochastische
Kontinuumsbrücke, aber die Hauptroute sollte nun um die Gitterresultate
zu Massenlücke, Poincaré-/Log-Sobolev-Ungleichungen, Wilson-Loop-Area-Law
und dem $'t$-Hooft-Regime organisiert werden.

Der aktualisierte CORE-Status ist für die Physik bewusst enger: Der
Zookeeper schließt den Riemann-$C^{(2)}$/MS2-Schritt im CCM/Fourier-Ast,
während RFEP v1.8 hier nur die folgende Normalform-Checkliste importiert.
Das beweist für sich genommen keine Yang--Mills-Kontinuums-Massenlücke;
Prime-Hub liefert nur eine Boundary-/Defektbrücke, keinen gelösten
physikalischen Operator.

Die Transferfelder sind:
\begin{description}
  \item[Operator/Form.] Transferoperator, Heat-bath-/Glauber-Generator und
  Wilson-Loop-Observablen.
  \item[Defekt.] Versagen gleichmäßigen Korrelationsabfalls,
  RG-Transportdefekt und Kontinuums-Rekonstruktionsdefekt.
  \item[Approximation.] Endliche Gitter, Slab-Sigma-Modelle,
  $'t$-Hooft-Large-$N$-Skalierung und Block-Spin-RG.
  \item[Gap.] Poincaré-/LSI-/Bakry--Emery-Konstanten, exponentieller
  Korrelationsabfall und Area-Law-Exponent.
  \item[Grenzübergang.] Endliches Volumen $\to$ unendliches Volumen
  $\to$ Large $N$ oder $'t$-Hooft-Regime $\to$
  Kontinuums-/OS-Rekonstruktion.
\end{description}

Das Beweisprogramm sollte deshalb lauten:
\begin{align*}
  &\text{finite-volume Poincare/LSI}\\
  &\Longrightarrow \text{unendlichvolumige lokale Transferlücke}\\
  &\Longrightarrow \text{Wilson area law}\\
  &\Longrightarrow \text{conditional continuum transfer}.
\end{align*}
Die Zerlegung von U(N) nach SU(N) ist hier besonders nützlich: Das
U(1)-Feld kann als Environment behandelt werden, während das
SU(N)-Marginal den nichtabelschen Gap trägt.  Nach dieser Reorganisation
bleibt als kritischer Block die Kontinuumsskalierungs-Aussage, dass
eine positive lokale Gitterlücke zu einer positiven physikalischen Massenlücke
transportiert.

\begin{table}[htbp]
\centering
\small
\caption{Beweisstatus-Übersicht}
\label{tab:proof-status}
\begin{tabularx}{\textwidth}{@{}L{4.0cm}L{3.0cm}Y@{}}
\toprule
\textbf{Komponente} & \textbf{Status} & \textbf{Referenz} \\
\midrule
\multicolumn{3}{@{}l}{\textit{Bewiesen (unbedingt):}} \\
Holley--Stroock PI ($\beta < \beta_0$) & Theorem & Thm.~\ref{thm:lattice-gap}, Schritt~1 \\
Reflexionspositivität (alle $\beta$) & Lemma & Lem.~\ref{lem:reflection-positivity} \\
Dobrushin--Zegarlinski ($\beta < \beta_0$) & Theorem & Thm.~\ref{thm:lattice-gap}, Schritt~2 \\
\midrule
\multicolumn{3}{@{}l}{\textit{Konditional (unter expliziten Annahmen):}} \\
Massenlücke (CL1+CL2'+CL3) & Kond.\ Theorem & Thm.~\ref{thm:conditional-gap} \\
Polynomiale LSI-Skalierung & Kond.\ Prop. & Prop.~\ref{prop:polynomial-lsi} \\
Kingman-/summierbarer-Defekt-Weg & Bedingtes Kriterium & Kor.~\ref{cor:kingman-gap} \\
\midrule
\multicolumn{3}{@{}l}{\textit{Numerisch bestätigt:}} \\
Defektiver Dobrushin ($\eta \gg 1$) & Monte Carlo & Bem.~\ref{rem:dobrushin-numerics} \\
Birkhoff $\tau_B < 1$ ($\beta \leq 6$) & Transfermatrix & Bem.~\ref{rem:birkhoff-numerics} \\
Kingman $\lambda < 0$ (getestetes $\beta$-Gitter) & Transfermatrix-Diagnostik & Bem.~\ref{rem:birkhoff-numerics} \\
\midrule
\multicolumn{3}{@{}l}{\textit{Offen:}} \\
Summierbarer RG-Defekt analytisch für $\beta \to \infty$ & Vermutung & Abschn.~\ref{sec:gamma-rg} \\
Gap-Transport (alle $\beta$) & Vermutung & Verm.~\ref{conj:gap-transport} \\
Warmkorridor-/OS-gefährliche Defektkapazität und Skalengrad-Ledger & Future-Work-Guardrail & Bem.~\ref{rem:warm-corridor-guardrail} \\
\bottomrule
\end{tabularx}
\end{table}

\begin{remark}[Defektkapazitäts-Guardrail]\label{rem:warm-corridor-guardrail}
Das aktuelle Beweisledger präzisiert die offene GT2-/summierbare-Defekt-
Bedingung durch eine Warmkorridor-Diagnostik. Ein negativer Kingman-Exponent
$\lambda<0$ bleibt nur ein Mittelwertsignal für Kontraktion: Er kann eine
kohärente Familie OS-gefährlicher Bad-Scale-Korridore übersehen. Künftige
Checks sollen daher good-scale-density, OS-sichere multiplikative Degradation,
OS-gefährliche Korridorkapazität und Zielgap auf derselben RG-Skalenachse
bilanzieren. Insbesondere ist eine RG-Folge mit $\lambda<0$, aber nicht
summierbarer OS-gefährlicher Korridorkapazität ein False-positive-Kanal, kein
Massenlücken-Transfer. Diese Bemerkung fügt keinen neuen Satz hinzu; sie hält
den schärferen Guardrail für den offenen Kontinuums-/RG-Schritt fest.
\end{remark}

% ======================================================================
\section{Diskussion}
% ======================================================================

\subsection{Bezug zu Lattice-QCD-Ergebnissen}

Das hier vorgestellte variationelle Argument ist mit der umfangreichen numerischen Evidenz aus Gitterrechnungen vereinbar. Für $SU(3)$ liefern Lattice-Simulationen eine Massenlücke (identifiziert mit der leichtesten Glueball-Masse) von ungefähr $\Delta \approx 1.5$--$1.7\;\mathrm{GeV}$ \cite{MorningstarPeardon1999, ChenXu2006}. Unser Ansatz sagt den numerischen Wert von $\Delta$ nicht voraus; er beweist strikte Positivität im Strong-Coupling-Gitterregime und formuliert die zusätzlichen Transportbedingungen, die jenseits dieses Regimes erforderlich sind. Tabelle~\ref{tab:glueball-spectrum} fasst etablierte Lattice-Monte-Carlo-Benchmarks zusammen und stellt sie den analytischen Schwellen unseres variationellen Transfer-Frameworks gegenüber.

\begin{table}[htbp]
\centering
\small
\caption{Lattice-QCD-Glueball-Massenspektrum und Kopplungsregime in reiner $SU(N)$-Eichtheorie.}
\label{tab:glueball-spectrum}
\begin{tabularx}{\textwidth}{@{}L{1.5cm}L{2.0cm}L{2.5cm}L{2.0cm}L{2.0cm}L{2.2cm}Y@{}}
\toprule
\textbf{Gruppe} & \textbf{Zustand $J^{PC}$} & \textbf{Gitter $M/\sqrt{\sigma}$} & \textbf{Masse (GeV)} & \textbf{Schwelle $\beta_0$} & \textbf{Skalierung $\beta_{\mathrm{cont}}$} & \textbf{Referenzen} \\
\midrule
$SU(2)$ & $0^{++}$ (Skalar)  & $3{,}78 \pm 0{,}05$ & $\approx 1{,}65$ & $0{,}336$ & $2{,}2 - 2{,}5$ & \cite{Martinelli1999, CreutzJacobsRebbi1983} \\
$SU(2)$ & $2^{++}$ (Tensor)  & $5{,}45 \pm 0{,}11$ & $\approx 2{,}38$ & $0{,}336$ & $2{,}2 - 2{,}5$ & \cite{Martinelli1999} \\
\midrule
$SU(3)$ & $0^{++}$ (Skalar)  & $3{,}55 \pm 0{,}04$ & $1{,}50 - 1{,}73$ & $0{,}165$ & $5{,}7 - 6{,}4$ & \cite{MorningstarPeardon1999, ChenXu2006} \\
$SU(3)$ & $2^{++}$ (Tensor)  & $4{,}78 \pm 0{,}06$ & $2{,}15 - 2{,}40$ & $0{,}165$ & $5{,}7 - 6{,}4$ & \cite{MorningstarPeardon1999, ChenXu2006} \\
$SU(3)$ & $0^{-+}$ (Pseudo)  & $5{,}69 \pm 0{,}08$ & $\approx 2{,}56$ & $0{,}165$ & $5{,}7 - 6{,}4$ & \cite{MorningstarPeardon1999} \\
\midrule
$SU(\infty)$ & $0^{++}$ (Large-$N$) & $3{,}37 \pm 0{,}15$ & --- & $O(1/N^2)$ & $g^2 N = \text{konst}$ & \cite{Luscher2010, MorningstarPeardon1999} \\
\bottomrule
\end{tabularx}
\end{table}

Zur Einordnung der Strong-Coupling-Schwelle: Für $SU(2)$ in $d=4$ verlangt die auf $\tanh$ basierende Dobrushin-Bedingung $\beta < \beta_0 = 2\,\mathrm{artanh}(1/6) \approx 0.336$. Lattice-Monte-Carlo-Studien der $SU(2)$-Eichtheorie arbeiten typischerweise bei $\beta \approx 2$--$3$ (Skalierungsregime), also ungefähr eine Größenordnung höher. Das Strong-Coupling-Regime liegt damit erwartungsgemäß weit vom physikalisch relevanten Kontinuumslimes entfernt, wie es für jedes Argument zu erwarten ist, das auf Clusterentwicklungen beruht (die nur bei hoher Temperatur bzw. starker Kopplung konvergieren).

\subsection{Die Rolle der strikten Konvexität}

Die strikte Konvexität von $\mathcal{F}$ motiviert den variationellen Rahmen,
soll aber nicht als eigenständiger Beweis der Operator-Massenlücke gelesen
werden. Im endlichen Volumen schließt der KL-Anteil wörtliche flache Richtungen
im Maß-Freie-Energie-Funktional aus; die quantitative Spektralaussage stammt
jedoch aus der Poincar\'e--Dobrushin-Schätzung und ihrer
Transferoperator-Übersetzung in Theorem~\ref{thm:lattice-gap}. Die
Hessian-Identität ist daher Diagnostik und Ordnungsprinzip, keine Alternative
zur Poincar\'e-Ungleichung.

Dies ist dem Geist nach mit entropisch stabilisierten Skalarfeldmechanismen
verwandt, aber der Yang--Mills-Lückenclaim dieser Arbeit ist genau die oben
bewiesene lokale Gitter-Transferlücke plus das bedingte
RG-/OS-Transportprogramm jenseits starker Kopplung.

\begin{remark}[Formale FST-Konvexität und Transport-Caveat]
\label{rem:fst-convexity-talagrand}
Formal suggeriert der KL-Anteil von $\mathcal{F}[\mu]$ einen Wasserstein-konvexen Mechanismus mit Skala vergleichbar zu $1/\beta$. Um daraus eine Talagrand- oder Poincar\'e-Schätzung für das Eichfeldmaß zu gewinnen, bräuchte man jedoch eine verifizierte geodätische Konvexitäts- beziehungsweise Curvature-Dimension-Aussage auf dem relevanten Eichorbit oder einer geeigneten Scheibe, plus Kompatibilität mit Reflexionspositivität. Diese Bemerkung ist daher nur eine Transport-Heuristik; sie beweist keine lineare untere Schranke $\kappa \gtrsim 1/\beta$.
\end{remark}

\begin{remark}[Annealed Influence als universeller Lücken-Mechanismus]
\label{rem:annealed-influence-universal}
Die in dieser Arbeit verwendete Dobrushin--Zegarlinski-Maschinerie -- Beschränkung des Spektralradius der Einflussmatrix zur Herstellung volumenunabhängiger Lücken -- ist nicht spezifisch für Gitter-Eichtheorie. Derselbe Mechanismus gilt, \emph{mutatis mutandis}, für zwei weitere Instantiierungen im FST-Programm:
\begin{enumerate}
\item \emph{Turbulenz (Schalenkaskaden):} Der schalenweise Energietransfer definiert eine Einflussmatrix $C_{jk}$ auf dem Inertialbereich. Die Dobrushin-Eindeutigkeitsbedingung $\rho(C) < 1$ ist äquivalent zur Stabilität von K41 als eindeutigem Gibbs-Zustand auf dem Kaskadengitter (Bemerkung~4.22 von~\cite{Geiger2026-TU}). Die ,,defektive Dobrushin''-Erweiterung behandelt Intermittenz-Korrekturen.
\item \emph{P\,vs\,NP (Zeugensuche):} Für NP-vollständige Instanzfamilien wird vermutet, dass die Einflussmatrix auf dem Gibbs-Maß des Zeugenraums $\rho(C) \geq 1$ hat, was schnelles Mischen und damit Polynomialzeit-Zeugenextraktion obstruiert (vgl.\ das P\,vs\,NP-Begleitpaper~\cite{Geiger2026-PvsNP}).
\end{enumerate}
Das \emph{defektive Poincar\'{e}/LSI-Lemma} fungiert somit als universelle Schablone: Entweder ist der Einfluss schwach und eine Lücke existiert (YM starke Kopplung, TU Inertialbereich), oder der Einfluss ist stark und die Lücke ist obstruiert (YM schwache Kopplung, P\,vs\,NP schwere Instanzen). Die Pattern-A-Architektur klassifiziert diese als zwei Seiten derselben Medaille.
\end{remark}

\begin{remark}[Deconfinement und endliche Temperatur]
Für $SU(N)$-Eichtheorie bei endlicher Temperatur in $d = 4$ (Gitter mit endlicher temporaler Ausdehnung $N_\tau$) bricht der Deconfinement-Phasenübergang bei $\beta_c(N_\tau)$ die $\mathbb{Z}(N)$-Zentrumssymmetrie spontan (der Polyakov-Loop erhält einen nichtverschwindenden Erwartungswert). Oberhalb der kritischen Temperatur verschwindet die Confinement-Stringspannung und das Spektrum ändert sich qualitativ: Das diskrete Glueball-Spektrum der Confinement-Phase weicht einem Kontinuum von Streuzuständen, wobei massive Screening-Moden (Debye-Masse $\sim gT$) bestehen bleiben. Dieser Übergang ist nicht relevant für das Millenniumsproblem bei Nulltemperatur (das die Theorie bei unendlichem Volumen und Nulltemperatur auf $\mathbb{R}^4$ betrifft), illustriert aber, dass die Natur des Spektrums empfindlich von der Reihenfolge der Limites abhängt (räumliches Volumen $\to \infty$ vor temporaler Ausdehnung $\to \infty$).
\end{remark}

\subsection{Mixture-Zerlegungsansatz}

\begin{remark}[Mixture-Zerlegung für schwache Kopplung]
\label{rem:mixture-decomposition}
Für $\beta > \beta_0$ degradiert die globale Holley--Stroock-Schranke, weil die Oszillation $\mathrm{osc}(S) \to \infty$ divergiert. Ein verfeinerter Ansatz zerlegt das Gibbs-Maß in metastabile Komponenten:
\[
\mu = \sum_{\alpha} w_\alpha \mu_\alpha, \quad w_\alpha > 0, \quad \sum_\alpha w_\alpha = 1,
\]
wobei jedes $\mu_\alpha$ in der Nähe eines lokalen Minimums der Wilson-Wirkung (eines ,,Potentialtopfs'') konzentriert ist. Die Spektrallücke erfüllt dann
\[
\Delta \geq \min\left( \min_\alpha \Delta_\alpha, \quad \Delta_{\mathrm{mix}} \right),
\]
wobei $\Delta_\alpha$ die Poincar\'e-Konstante innerhalb des Topfs ist (berechenbar über Bakry--\'Emery-Krümmung auf dem Orbitraum $\mathcal{A}/\mathcal{G}$) und $\Delta_{\mathrm{mix}}$ die Mischungsrate zwischen den Töpfen (kontrolliert durch die Eyring--Kramers-Formel für Barrierenüberschreitung).

In einem Toy-Einzel-Link-Topfmodell für $SU(2)$ würde die Ricci-Krümmung von
$S^3$ eine lokale Skala der Größenordnung $3-O(\beta^{-1})$ nahelegen. Für das
volle Gitter-Yang--Mills-Maß ist dies nur ein heuristisches Ziel: Man bräuchte
weiterhin eine gauge-fixierte metastabile Zerlegung, Kontrolle der
Eichorbit-Singularitäten und eine summierbare Inter-Topf-Mischschranke. Die
Mischungsrate zwischen Töpfen ist physikalisch mit Instanton-Barrieren der
Ordnung $e^{-8\pi^2/g^2}$ verwandt, aber keine finite-volume-Tunneling-Schätzung
in dieser Bemerkung wird als Beweis der Kontinuums-Massenlücke verwendet.
\end{remark}

\subsection{Grenzen und offene Probleme}

Unser Argument etabliert die Massenlücke rigoros nur im Strong-Coupling-Regime ($\beta < \beta_0$) und für die Kontinuumstheorie nur bedingt. Eine systematische Übersicht der konstruktiven Obstruktionen und ihrer korrespondierenden FST-Mitigationskanäle ist in Tabelle~\ref{tab:obstruction-taxonomy} dargestellt. Die zentralen offenen Probleme sind:

\begin{table}[htbp]
\centering
\small
\caption{Taxonomie konstruktiver Obstruktionen zur 4D-Yang--Mills-Massenlücke und korrespondierende FST-Mitigationsmechanismen.}
\label{tab:obstruction-taxonomy}
\begin{tabularx}{\textwidth}{@{}L{2.6cm}L{3.0cm}L{3.4cm}L{3.2cm}Y@{}}
\toprule
\textbf{Obstruktion / Regime} & \textbf{Physikalischer Ursprung} & \textbf{Klassischer Engpass} & \textbf{FST-Mitigationsmechanismus} & \textbf{Status} \\
\midrule
Lokale Strong-Coupling-Lücke ($\beta < \beta_0$) & Degradation der Freie-Energie-Konvexität & Exponentieller Holley--Stroock-Verlust $e^{-2d_\ell\beta}$ & Dobrushin--Zegarlinski-Kontraktion via Hesse-Schranke & Bewiesen (Thm.~\ref{thm:lattice-gap}) \\
\midrule
Eichbahn-Entartung (Alle $\beta$) & Flache Richtungen der Eichsymmetrie & Nicht-Invertierbarkeit der freien Hesse-Matrix & Restriktion auf eichinvariante Zylinderobservablen $\mathcal{H}_0 \subset L^2$ & Bewiesen (Lem.~\ref{lem:reflection-positivity}, Abschn.~\ref{subsec:hessian}) \\
\midrule
Gribov-Mehrdeutigkeit (Mittleres $\beta$) & Mehrblättrige transversale Eichfixierung & Spuriöse Nullmoden \& Spektralverschmutzung & Integrierbare Defektschranke mit subadditiver Kapazität & Formuliert (Abschn.~\ref{sec:gamma-rg}) \\
\midrule
Asymptotische Freiheit ($\beta \to \infty$) & Kontinuumslimes-Skalierung $\Delta(a)/a \sim \Lambda_{\mathrm{QCD}}$ & Zusammenbruch von Hochtemperaturentwicklungen & Multiskalen-$\Gamma$-Konvergenz \& summierbares GT2-Defektbudget & Konditional (Verm.~\ref{conj:gap-transport}, Prop.~\ref{prop:polynomial-lsi}) \\
\midrule
Reflexionspositivität ($a \to 0$) & Anisotropie \& Gitterartefakt-Aufbrechung & Verlust der OS-Positivität unter Block-Spin & Vordefinierte RP-kompatible Kegelprojektion \& Spiegelschleuse & Konditional (Lem.~\ref{thm:continuum-gap}, CL1) \\
\midrule
Räumliche Tensorisierung (Thermodynam.\ Limes) & $L^2$-Abschluss auf vollen Zeitscheiben & Lokaler Korrelatorabfall vs.\ globaler Spektralprojektor & Koerzives Komplement außerhalb des Mikroclusters & Offen \\
\bottomrule
\end{tabularx}
\end{table}
\begin{enumerate}
    \item \textbf{Lückenskalierung unter Renormierung:} Die Gitterlücke $\Delta_\Lambda \geq \kappa_* > 0$ ist für $\beta < \beta_0$ bewiesen (Theorem~\ref{thm:lattice-gap}). Im Kontinuumslimes gilt wegen asymptotischer Freiheit $\beta(a) \to \infty$, so dass das Strong-Coupling-Regime schließlich verlassen wird. Ob die physikalische Lücke $\Delta_{\mathrm{phys}} = \Delta_{\mathrm{lat}}/a$ für $a \to 0$ nach unten beschränkt bleibt, ist eine nichtperturbative Frage und eng mit dem Renormierungsgruppenfluss verbunden. Dies ist der Kern des Millennium-Problems.

    Ein Nullstellenfreiheits- oder vollständiges-Analytizitätsresultat könnte eine strukturelle Route zur Lückenpersistenz liefern, indem es Phasenübergänge erster Ordnung entlang des Renormierungsflusses ausschließt (siehe Bemerkung~\ref{rem:fst-parallels}). Ein solcher Input wird hier nicht als bewiesener Satz verwendet.

    \item \textbf{Konstruktive Existenz des Kontinuumslimes:} Die Annahmen (CL1)--(CL3) in Theorem~\ref{thm:continuum-gap} bündeln sowohl die Existenz OS-positiver Schwinger-Funktionen als auch uniformes exponentielles Clustering. Balabans Programm liefert UV-Stabilität, schließt aber die volle IR-/Konstruktionslücke in 4D nicht. \emph{Nachtrag (Mai 2026):} Kirk~\cite{Kirk2026} hat einen Zenodo-Preprint zu einer bedingten SU(2)-Kontinuumslimes-Konstruktion mit Dobrushin--Shlosman-vollständiger Analytizität und Pro-Polymer-Clusterabschätzungen veröffentlicht. Der Preprint beruht explizit auf Weak-Coupling-Hypothesen, und ein begleitender nullstellenfreier Streifen bleibt unabhängig zu prüfen. Diese Resultate sind daher bedingte Preprint-Evidenz, kein geschlossener Beweis von (CL1)--(CL3).
    \item \textbf{Erweiterung auf alle $\beta$:} Das Dobrushin--Zegarlinski-Kriterium (Schritt~2) liefert $\kappa_* > 0$ nur für $\beta < \beta_0$. Obwohl endliches Volumen scharfe Phasenübergänge bei fixem $|\Lambda|$ ausschließt, bleibt die Persistenz einer uniformen Lücke für alle $\beta > 0$ im thermodynamischen Limes eine offene Vermutung (siehe Bemerkung~\ref{rem:all-beta}).
\end{enumerate}

Eine vollständige Lösung des Millennium-Problems erfordert sowohl die konstruktive Existenz des Kontinuumslimes als auch den Nachweis, dass die physikalische Massenlücke unter dem Renormierungsgruppenfluss erhalten bleibt.

\begin{remark}[Massenlücke als RG-Fixpunkt-Eigenschaft]\label{rem:rg-fixpoint}
Ein alternativer konzeptueller Rahmen vermeidet den Limes $\beta \to \infty$
vollständig, indem er die Massenlücke als \emph{Fixpunkt}-Aussage
umformuliert. Man definiere die RG-transformierte freie Energie
$\mathcal{F}_\ell[\mu] := \mathcal{F}[\mu; \beta(\ell), \Lambda(\ell)]$
bei Skala $\ell$, wobei $\beta(\ell)$ und $\Lambda(\ell)$ dem
Wilson-RG-Fluss folgen. Die Massenlücken-Bedingung lautet dann: Es
existiert ein nicht-trivialer RG-Fixpunkt $\mathcal{F}^*$ mit
Spektrallücke $\Delta^* > 0$, und der Fluss $\ell \mapsto
\mathcal{F}_\ell$ konvergiert im Infraroten gegen $\mathcal{F}^*$.
In dieser Formulierung:
\begin{itemize}
  \item Das Strong-Coupling-Resultat (Theorem~\ref{thm:lattice-gap}) zeigt,
    dass $\mathcal{F}_0$ (UV-Startpunkt) eine Lücke besitzt.
  \item Asymptotische Freiheit garantiert, dass der UV-Fixpunkt Gaussisch
    ist (freie Theorie, keine Lücke).
  \item Die offene Frage reduziert sich auf: Erzeugt der RG-Fluss vom
    Gaußschen UV-Fixpunkt ein konfinierendes IR-Regime mit
    $\Delta^* > 0$?
    (Anmerkung: 4D Yang--Mills ist vermutlich konfinierend und fließt nicht zu einem nicht-trivialen konformen Fixpunkt; der ,,Fixpunkt'' hier bezieht sich auf eine stabile massive Phase, nicht auf eine konforme Feldtheorie.)
\end{itemize}
Dies verschiebt das Problem von einem Limes-Argument ($\beta \to \infty$,
wo die Holley--Stroock-Schranken degenerieren) zu einer
\emph{Stabilitätsanalyse des RG-Flusses} -- einer Frage über das
Infrarotverhalten endlich vieler effektiver Kopplungen, sobald die
relevante Trunkierung identifiziert ist. Ein unabhängig verifizierter
zero-free- oder complete-analyticity-Satz könnte so gelesen werden, dass
der Fluss keine Fixpunkt-Bifurkationen aufweist (keine
Phasenübergänge). Das bleibt hier eine offene Zusatzannahme.
\end{remark}

\begin{remark}[Massenlücke als renormierungsstabile LSI-Eigenschaft]\label{rem:lsi-equivalence}
Korollar~\ref{cor:kingman-gap} formuliert eine renormierungsstabile
Log-Sobolev-Route: Eine Familie von LSI-Konstanten $\{\kappa_k\}_{k \geq 1}$
mit summierbarem Birkhoff-Defekt oder GT2-Fehlern würde die Massenlücke
transportieren. Der Kingman--Ljapunow-Exponent
$\lambda = \langle\log\tau_B\rangle$ ist dabei eine nützliche Diagnostik,
aber ohne Summierbarkeitskontrolle keine äquivalente Charakterisierung.
Die relevante numerische Größe bleibt zugänglich
(Bemerkung~\ref{rem:kingman-numerical}: $\lambda \approx -0{,}18$ bei $\beta=8$,
$\lambda \approx -4 \times 10^{-5}$ bei $\beta=20$)
und strukturell analog zum Ljapunow-Exponenten der dynamischen Systemtheorie.
\end{remark}

\begin{remark}[Analytischer Weg zu $\lambda < 0$: Doeblin-Minorisierung auf typischen Mengen]
\label{rem:doeblin-analytical}
Die numerische Beobachtung $\lambda < 0$ auf dem getesteten $\beta$-Gitter (Bemerkung~\ref{rem:kingman-numerical}) lässt eine natürliche analytische Beweisstrategie über die \emph{Doeblin-Minorisierung} für Produkte zufälliger positiver Operatoren zu.

\textbf{Schritt 1: Zufälliger positiver Operatorkozykel.} Auf jeder RG-Stufe~$k$ ist der Block-Spin-Transferkern $R_k$ ein positiver Operator, dessen Zufälligkeit aus der Randkonfiguration unter dem Gibbs-Maß der vorherigen Skala stammt. Dies definiert einen positiven Operatorkozykel $\{R_k\}_{k \geq 1}$.

\textbf{Schritt 2: Typische-Mengen-Doeblin-Bedingung.} Der analytische Schlüssel ist eine \emph{Minorisierung auf typischen Randkonfigurationen}: es existiert ein Ereignis $\mathcal{T}_k$ (,,typische Ränder'') mit $\mathbb{P}(\mathcal{T}_k) \geq 1 - e^{-c\beta}$ und eine feste Referenzmasskomponente $\nu$, sodass auf $\mathcal{T}_k$:
\[
R_k(x, \cdot) \;\geq\; \alpha(\beta)\, \nu(\cdot) \qquad \text{(als Masse),}
\]
mit $\alpha(\beta) > 0$ explizit. Dies ist das Birkhoff-projektive Analogon der defekten Dobrushin-Bedingung für die Poincar\'e-Ungleichung (Proposition~\ref{prop:defective-dobrushin}).

\textbf{Schritt 3: Minorisierung $\Rightarrow$ negative Drift.} Doeblins Minorisierung impliziert eine quantitative Birkhoff-Kontraktion: $\tau_B(R_k) \leq 1 - \delta(\beta)$ auf $\mathcal{T}_k$. Die erwartete Log-Kontraktion erfüllt:
\[
\mathbb{E}[\log \tau_B(R_k)] \;\leq\; \mathbb{P}(\mathcal{T}_k) \cdot \log(1 - \delta(\beta)) \;+\; \mathbb{P}(\mathcal{T}_k^c) \cdot 0 \;\approx\; -(1 - e^{-c\beta})\, \delta(\beta) \;<\; 0,
\]
da $\log \tau_B \leq 0$ stets gilt (positive Operatoren kontrahieren in Hilberts projektiver Metrik).

\textbf{Schritt 4: Kingmans Theorem.} Die verbleibenden technischen Bedingungen -- Stationarität/Ergodizität der Randumgebung entlang der RG-Kaskade (via Markov-Eigenschaft und Translationssymmetrien) und Integrierbarkeit von $\log \tau_B(R_k)$ (aus der Tail-Schranke $0 \geq \log \tau_B \geq -C \cdot \mathrm{osc}(W)$) -- sind Standard. Kingmans subadditives ergodisches Theorem liefert dann $\lambda < 0$ als reines Drift-Resultat.

Diese Strategie präzisiert die numerische Beobachtung, dass einzelne $\tau_B(R_k)$ auf der gröbsten Skala gegen~1 gehen, während der gemittelte Log-Exponent negativ bleibt -- die Signatur eines \emph{selten-schlecht / typisch-gut}-Mechanismus, analytisch erfasst durch Minorisierung plus Kingman. Das Programm reduziert den analytischen Beweis von $\lambda < 0$ auf die Etablierung der Doeblin-Konstante $\alpha(\beta) > 0$ für den SU(2)-Block-Spin-Kern, ein konkretes Problem der stochastischen Geometrie auf kompakten Lie-Gruppen.

Für verwandte analytische Methoden zu Produkten zufälliger positiver Operatoren siehe Bougerol--Lacroix~\cite{BougerolLacroix1985}.
\end{remark}

\begin{remark}[Koordinationsbarriere als Phasenübergangs-Analogon]
\label{rem:coordination-barrier}
Die Koordinationsbarriere in der Yang--Mills-Theorie -- die
Unmöglichkeit des Systems, sich auf eine einzelne Vakuumkonfiguration
zu koordinieren -- ist strukturell analog zu einem Phasenübergang
im spieltheoretischen Sinne von FST-III~\cite{FST-Unified2026}.
In endlichen Spielen entstehen Koordinationskosten, wenn mehrere
Nash-Gleichgewichte koexistieren und die Spieler nicht kollektiv eines
auswählen können; der Übergang von reinen zu gemischten Strategien
wird durch das Koordinationsversagen erzwungen.
In Yang--Mills spielen die Eichfeldkonfigurationen auf dem Gitter die
Rolle der Strategien, und das Gibbs-Maß $\mu_\Lambda$ ist das gemischte
Gleichgewicht. Die Massenlücke $\Delta > 0$ entsteht als
``Koordinationskosten'' -- die minimale Energiestrafe für den
Übergang zwischen verschiedenen Vakuumsektoren. Genau wie in der
Spieltheorie, wo die Koordinationskosten nur am kritischen Punkt
(der Phasenübergangstemperatur) verschwinden, würde die Massenlücke
nur am Deconfinement-Übergang ($T = T_c$) verschwinden, bleibt aber
bei Nulltemperatur strikt positiv. Diese Perspektive legt nahe, dass
die Persistenz von $\Delta > 0$ im Kontinuumslimes durch die
\emph{Abwesenheit} eines Nulltemperatur-Phasenübergangs bestimmt
wird -- konsistent mit der offenen Suche nach einem zero-free- oder
complete-analyticity-Zertifikat.
\end{remark}

\begin{remark}[Stabilitäts-Selektions-Axiomatik]
\label{rem:stability-selection}
Die Stabilitäts-Selektions-Axiome (S1)--(S4) aus dem allgemeinen FST-Prinzip-Beitrag liefern eine axiomatische Grundlage für die Gitter-Massenlücke: Das Yang--Mills-Vakuum erfüllt alle vier Bedingungen, wobei die Poincar\'e-Konstante als Stabilitätszertifikat dient.
\end{remark}

\begin{remark}[Pfadintegral als Nash-Selektor]
\label{rem:path-integral-nash}
Das Pfadintegral-Maß $e^{-S[A]/\hbar}$ lässt sich als Nash-Selektionsmechanismus interpretieren: Unter allen Feldkonfigurationen wählt es das thermodynamisch stabile Vakuum als das eindeutige Gleichgewicht des Entropie-Energie-Spiels aus.
\end{remark}

\begin{remark}[QCD-Confinement als Nash-Phasenübergang]
\label{rem:qcd-confinement-nash}
Confinement in der QCD lässt sich als Nash-Phasenübergang verstehen: Unterhalb der Deconfinement-Temperatur übersteigen die Koordinationskosten der Farbladungsausrichtung den entropischen Gewinn, was das System in Farb-Singulett-Bindungszustände zwingt.
\end{remark}

\begin{remark}[Ordnungsparameter-Dualität und alternative Beweisrouten]
\label{rem:order-parameter-duality}
Drei alternative Beweisrouten für die Massenlücke ergeben sich aus der Kosten-Antwort-Dualität: (D1)~direkte Poincar\'e-Schranke über Bakry--\'Emery-Krümmung, (D2)~Transferoperator-Spektralanalyse und (D3)~variationelle Charakterisierung der Lücke als Infimum des Dirichlet-zu-Neumann-Verhältnisses.
\end{remark}

\begin{remark}[Pattern-A-Master-Stabilität -- problemübergreifende Struktur]
Diese Arbeit instanziiert das Pattern-A-Master-Stabilitätsprinzip aus dem vereinheitlichten Framework \cite{FST-Unified2026}: Strikte Konvexität der renormierten freien Energie $F$, kombiniert mit einem neutralen führenden Resolventenzweig und Dominanz zweiter Ordnung, ist das variationelle Stabilitätssignal. Die konkrete Spektrallücke $\Delta > 0$ folgt in dieser Arbeit jedoch erst aus der Poincar\'e--Dobrushin-Maschinerie und ihrer Transferoperator-Übersetzung.
\end{remark}

\subsection{Ausblick}

Der variationelle Freie-Energie-Ansatz lässt sich möglicherweise auf Yang--Mills--Higgs-Theorien und die Untersuchung des Confinements erweitern, wobei die Stringspannung als die Hesse-Matrix eines geeignet modifizierten Freie-Energie-Funktionals über Wilson-Loop-Observablen interpretiert werden kann.


\subsection{Offene Richtungen}

Die folgenden Bemerkungen sammeln sechs konzeptuelle Denkrichtungen, die sich natürlich aus dem oben entwickelten variationellen Rahmenwerk ergeben. Sie sind als Impulse für künftige Arbeiten gedacht, nicht als bewiesene Resultate.

\begin{remark}[Topologische Sektoren als ,,keine flachen Richtungen'']
\label{rem:topological-no-flat}
Jede nicht-triviale topologische Deformation -- d.\,h.\ jede Änderung der Instantonenzahl $\nu = c_2(P) \in \mathbb{Z}$ -- erfordert die Überwindung einer Energiebarriere der Ordnung $8\pi^2/g^2$ (Ein-Instanton-Wirkung). Dies impliziert, dass die freie Energie $F[\mu]$ \emph{keine flachen Richtungen} im topologischen Sektor besitzt: Jede Störung, die die Topologie ändert, verursacht strikt positive Kosten. Dies liefert eine konzeptuelle Grundlage für die Positivität der Hesse-Matrix am Vakuum und ergänzt die quantitative Poincar\'e-Schranke.
\end{remark}

\begin{remark}[Verknotete Flussröhren als Kandidaten für die erste massive Anregung]
\label{rem:knotted-flux}
Verknotete chromoelektrische Flussröhren (Glueball-Kandidaten) liefern eine physikalische Interpretation der ersten massiven Anregung über dem Vakuum. Die Massenlücke $m$ entspricht dem leichtesten Glueball, dessen Stabilität durch topologische Erhaltung erzwungen wird: Die Flussröhre kann sich nicht entwinden, ohne eine Freie-Energie-Barriere zu überqueren. Im FST-Rahmenwerk ist dies die Lücke zwischen dem Vakuum-Eigenwert $\lambda_0 = 0$ des Transferoperators und dem ersten angeregten Eigenwert $\lambda_1 = m$.
\end{remark}

\begin{remark}[Nicht-abelsche Kramers--Wannier-Dualität]
\label{rem:non-abelian-kw}
Eine nicht-abelsche Kramers--Wannier-Dualität würde das Regime schwacher Kopplung ($\beta \to \infty$, $g \to 0$) auf eine duale Strong-Coupling-Beschreibung in Monopol-/Vortex-Variablen abbilden. Unter einer solchen Dualität würde die Massenlücke bei schwacher Kopplung dem Confinement dualer Monopole bei starker Kopplung entsprechen -- einem gut verstandenen Regime. Obwohl keine rigorose nicht-abelsche KW-Dualität existiert, liefert die 't~Hooft-Dualität (elektrisch--magnetisch) partielle Evidenz, und das FST-Rahmenwerk ist agnostisch bezüglich der Variablenwahl: Die freie Energie $F[\mu]$ lässt sich prinzipiell in dualen Variablen umformulieren.
\end{remark}

\begin{remark}[Haar-Maß-Entropie: ein entropischer Zugang zum Confinement]
\label{rem:haar-entropy}
Ein alternativer Zugang zum Confinement: Die Haar-Maß-Entropie $S_{\mathrm{Haar}} = \log \mathrm{vol}(G^{|\Lambda|})$ der Eichgruppe wächst extensiv mit dem Gittervolumen. Freie (deconfinierende) Quarks würden korrelierte Konfigurationen erfordern, die einen exponentiell kleineren Anteil des Eichorbitraums einnehmen, mit Entropiedefizit $\Delta S \sim L^3 \cdot \log N$. Das RFEP sagt voraus, dass das System die Konfiguration wählt, die $\Phi = S - \beta E$ unter Stabilitätsbedingungen maximiert; für $\beta$ unterhalb der Deconfinement-Temperatur dominiert der entropische Term und erzwingt Farb-Singulett-Zustände (Confinement).
\end{remark}

\begin{remark}[RG-Monotonie-Kandidat $M(\beta)$]
\label{rem:rg-monotone}
Ein Kandidat für eine RG-monotone Größe ist
\[
M(\beta) = \beta \cdot \kappa(\beta) \cdot \xi(\beta)^2,
\]
wobei $\kappa$ die Poincar\'e-Konstante und $\xi$ die Korrelationslänge bezeichnet. Bei $\beta = 0$ (unendliche Temperatur) gilt $M(0) = 0 \cdot \kappa_{\mathrm{Haar}} \cdot 0 = 0$. Bei endlichem $\beta$ liefert asymptotische Freiheit $\xi \sim a^{-1} e^{c/g^2}$, und falls Gap Transport gilt, $\kappa \geq c'/\xi^2$, woraus $M(\beta) \geq \beta \cdot c' = O(\log(1/a))$ folgt -- monoton wachsend unter dem RG-Fluss. Die Monotonie von $M$ würde einen unabhängigen Beweis der Massenlücken-Persistenz liefern.
\end{remark}

\begin{remark}[Zwei-Phasen-Interpolation und Freie-Energie-Analytizität]
\label{rem:two-phase}
Falls die freie Energie $f(\beta) = -\lim_{|\Lambda| \to \infty} |\Lambda|^{-1} \log Z_\Lambda(\beta)$ für alle $\beta > 0$ analytisch ist, tritt kein Phasenübergang auf und die Massenlücke (die in $\beta$ analytisch ist, wo sie existiert) kann bei keinem endlichen $\beta$ verschwinden. Ein nullstellenfreier Streifen würde Evidenz für die Analytizität der Zustandssumme liefern, aber die Analytizität der freien Energie erfordert zusätzliche Kontrolle über den thermodynamischen Limes.
\end{remark}

\subsection{Vergleich mit konkurrierenden Ansätzen}

\begin{table}[htbp]
\centering
\scriptsize
\caption{Vergleich der Ansätze zur Yang--Mills-Massenlücke}
\label{tab:comparison}
\setlength{\tabcolsep}{4pt}
\begin{tabularx}{\textwidth}{@{}L{3.6cm}Y Y Y@{}}
\toprule
& \textbf{FST (diese Arbeit)} & \textbf{Kirk (2026)} & \textbf{TMT (Garc\'ia Baquero)} \\
\midrule
\textbf{Ontologischer Status des Gitters} & Regularisierung (Kontinuumslimes erforderlich) & Regularisierung (Kontinuumslimes via Pro-Polymer) & Fundamental (diskrete Prägeometrie) \\
\midrule
\textbf{UV-Problem} & Auf CL1--CL3 verschoben & Gelöst via vollständige Analytizität für SU(2) & Konstruktionsbedingt absent \\
\midrule
\textbf{Massenlücken-Mechanismus} & Freie-Energie-Konvexität $+$ Poincar\'e--Dobrushin & Dobrushin--Shlosman vollständige Analytizität & Rigidität der Gitter-Bindungen \\
\midrule
\textbf{Eichgruppe} & Beliebige kompakte Gruppe $G$ (explizit für SU(2)) & SU(2) (SU(3)-Erweiterung offen) & SU(3) nativ \\
\midrule
\textbf{Kontinuumslimes} & Bedingt (Vermutung~\ref{conj:gap-transport}) & Partiell (nullstellenfreier Streifen schließt Übergänge 1.~Ordnung aus) & Nicht anwendbar (diskret) \\
\midrule
\textbf{Reflexionspositivität} & Starke Kopplung (Lemma~\ref{lem:reflection-positivity}) & Impliziert durch vollständige Analytizität & Nicht adressiert \\
\midrule
\textbf{Limitierungen} & Lücken-Transport unbewiesen; CL2' heuristisch & Nur SU(2); keine explizite Lückenschranke & Keine Kontinuums-QFT; keine Wightman-Axiome \\
\bottomrule
\end{tabularx}
\end{table}

\subsection{Fahrplan zum Millenniumsproblem}

Die Resultate dieser Arbeit liefern einen selbständigen variationellen Beweis
einer volumenunabhängigen lokalen Transferlücke im Strong-Coupling-Regime und eine
\emph{bedingte} Rekonstruktionsroute für die Kontinuumstheorie.
Tabelle~\ref{tab:clay-criteria-matrix} liefert ein detailliertes Audit der vier zentralen
Anforderungen des Clay-Millennium-Preises sowie des formalen Verifikationsstatus und der
verbleibenden rigorosen mathematischen Brücken.

\begin{table}[htbp]
\centering
\small
\caption{Anforderungen des Clay-Millennium-Problems an Existenz und Massenlücke der 4D-Yang--Mills-Theorie vs.\ FST-Status.}
\label{tab:clay-criteria-matrix}
\begin{tabularx}{\textwidth}{@{}>{\hsize=0.9\hsize}Y >{\hsize=1.1\hsize}Y >{\hsize=1.0\hsize}Y >{\hsize=0.8\hsize}Y >{\hsize=1.2\hsize}Y@{}}
\toprule
\textbf{Clay-Anforderung} & \textbf{Mathematische Definition} & \textbf{FST-Ergebnis} & \textbf{Status} & \textbf{Verbleibende Brücke} \\
\midrule
Axiomatische 4D-QFT auf $\mathbb{R}^4$ & Osterwalder--Schrader- / Wightman-Axiome ($O(4)$-Invarianz, Spektralbedingung) & Konditioniert auf CL1 (Existenz eines OS-positiven Kontinuumsmaßes) & Konditional (Lem.~\ref{thm:continuum-gap}) & Rigorose Balaban-artige UV-Konvergenz mit verifizierter Pro-Polymer-Summierbarkeit \\
\midrule
Strikt positive Massenlücke $\Delta > 0$ & $\inf \sigma(H)|_{\mathcal{H}_{\mathrm{phys}}^\perp} \ge \Delta > 0$ & Lokale Transferlücke $\Delta_\Lambda \ge \kappa_* > 0$ für $\beta < \beta_0$; uniforme Skalierung unter GT2 & Bewiesen ($\beta < \beta_0$) / Konditional ($\beta \to \infty$) & Beweis der Abwesenheit von Phasenübergängen 1.~Ordnung oder summierbares Defektbudget $\sum \varepsilon_k < \infty$ \\
\midrule
Nichtabelsche Eichinvarianz & Eichinvarianter Hilbertraum $\mathcal{H}_{\mathrm{phys}}$ für kompakte Lie-Gruppe $G$ & Exakte Invarianz unter kompakten Gitter-Eichtransformationen & Bewiesen & Explizit für $SU(2)$, algebraisch exakt für allgemeines kompaktes $G$ \\
\midrule
Confinement / Area Law & $\langle W(C)\rangle \le C e^{-\sigma \mathrm{Area}(C)}$ für Wilson-Schleifen & Area Law im Strong Coupling bewiesen; sauber getrennt vom $L^2$-Operatorgap & Bewiesen auf Gitter / Diagnostisch im Limit & Volle nicht-perturbative Stringspannung im Kontinuumslimes $a \to 0$ \\
\bottomrule
\end{tabularx}
\end{table}

Drei konkrete
mathematische Strategien bleiben, um die Lücke zum vollständigen Millenniumsproblem
zu schließen:

\begin{enumerate}
    \item \textbf{Kirks Konstruktion auf $SU(N)$ erweitern.}
    Der direkteste Weg zu unbedingten Resultaten erfordert die Verifikation
    der Bedingungen~(K1)--(K3) von Proposition~\ref{prop:kirk-mapping} für
    $G = SU(N)$ mit $N \geq 3$. Die wesentliche technische Herausforderung besteht
    in der gleichmäßigen Kontrolle der Peter--Weyl-Koeffizienten $a_\pi(\beta)$
    über den Darstellungsturm von~$SU(N)$. Für $N = 3$ (physikalische
    QCD) genügen wahrscheinlich die Fundamental- und die adjungierte Darstellung
    aufgrund der Casimir-Lücke zwischen diesen und höheren Darstellungen.

    \item \textbf{Eingeschränkte Poincar\'e-Ungleichung für $\beta$-Unabhängigkeit.}
    Die Holley--Stroock-Schranke~\eqref{eq:single-link-gap} degeneriert
    exponentiell für $\beta \to \infty$. Eine eingeschränkte Poincar\'e-Ungleichung
    auf der Menge $\Omega_\varepsilon := \{U : S_W[U] \leq
    \langle S_W\rangle + \varepsilon|\Lambda|\}$ könnte
    $\beta$-unabhängige Schranken liefern, unter Verwendung von
    Large-Deviation-Abschätzungen für die Wilson-Wirkung
    (Chatterjee~\cite{Chatterjee2016}) und der eingeschränkten
    Poincar\'e-Technik von Milman~\cite{Milman2009}. Dies würde die
    Notwendigkeit von Kirks vollständiger Analytizität gänzlich eliminieren.

    \item \textbf{Hierarchische RG-Kette für die Poincar\'e-Konstante.}
    Ein intermediarer Ansatz zerlegt $\Lambda$ in Blöcke der
    Größe~$L^4$ und wendet die Holley--Stroock-Schranke auf jeder RG-Skala
    unabhängig an. Asymptotische Freiheit stellt sicher, dass die effektive
    Kopplung $\beta_{\mathrm{eff}}(k)$ nur logarithmisch wächst,
    was eine Gesamt-Poincar\'e-Degradation von
    $\log(1/\kappa_{\mathrm{total}}) = O(\log^2(1/a))$ ergibt -- polynomiell
    statt exponentiell -- was genügen könnte, um die physikalische
    Lücke $\Delta_{\mathrm{phys}} = \Delta_{\mathrm{lat}}/a > 0$ zu erhalten.
\end{enumerate}

Jede einzelne dieser drei Strategien würde, falls erfolgreich, das
Yang--Mills-Millenniumsproblem für die jeweilige Eichgruppe lösen.

\begin{conjecture}[Lücken-Transport unter Block-Spin-RG]
\label{conj:gap-transport}
Sei $\kappa(a)$ die Poincar\'e-Konstante des Gitter-Yang--Mills-Masses bei Gitterabstand $a$. Unter einem Block-Spin-Renormierungsgruppenschritt $a \to 2a$ erfüllt die Poincar\'e-Konstante
\[
\kappa(2a) \geq c \cdot \kappa(a)
\]
für eine universelle Konstante $c > 0$, die nur von der Eichgruppe abhängt. Insbesondere gilt: Falls $\kappa(a_0) > 0$ bei einem Anfangsgitterabstand $a_0$, dann $\kappa(a) \geq c^{\log_2(a/a_0)} \kappa(a_0) > 0$ für alle $a > a_0$, was die Persistenz der Massenlücke unter Vergröberung (coarse-graining) etabliert.
\end{conjecture}

\begin{remark}[RG-Schritt als krümmungserhaltender Markov-Kern]
\label{rem:rg-markov-kernel}
Die Block-Spin-RG-Abbildung $\mathcal{R}$ lässt sich als Markov-Kern auf dem Raum der Masse über Eichkonfigurationen auffassen. In diesem Rahmen transformiert sich die Poincar\'e-Konstante als
\[
\kappa(\mathcal{R}_* \mu) \geq \kappa(\mu) \cdot \left(1 - \|\mathcal{R} - \mathrm{id}\|_{\mathrm{op}}^2 / \kappa(\mu) \right),
\]
sofern der Kern $\mathcal{R}$ eine Bakry--\'Emery-Krümmungsbedingung erfüllt: $\mathrm{Ric}_{\mathcal{R}} \geq -K$ für ein $K \geq 0$. Für Block-Mittelungskerne auf kompakten Gruppen gilt $K = O(1/\beta)$ (Krümmung der Gruppenmannigfaltigkeit), sodass die Degradation pro RG-Schritt $O(1/\beta)$ beträgt -- vernachlässigbar im Regime schwacher Kopplung.

Dies liefert einen konkreten Mechanismus für den Lücken-Transport (Vermutung~\ref{conj:gap-transport}): Die Bakry--\'Emery-Krümmung des RG-Kerns, kombiniert mit der Ricci-Krümmung der Eichgruppe, kontrolliert die Poincar\'e-Konstante entlang des RG-Flusses.
\end{remark}

\subsection{Hierarchischer Ansatz für den Kontinuumslimes}
\label{sec:hierarchical-continuum}

Der naive Kontinuumslimes $a \to 0$ ($\beta \to \infty$) lässt die Holley--Stroock-Schranke degenerieren: $\kappa_l \geq \kappa_{\mathrm{Haar}}\, e^{-2d_l\beta} \to 0$. Wir adressieren dies über eine hierarchische Block-Spin-Renormierungsgruppe (RG) innerhalb des Log-Sobolev-Frameworks.

\begin{definition}[Block-Spin-RG-Schritt]
\label{def:block-rg}
Gegeben ein Gitter $\Lambda$ mit Abstand $a$, definiere das vergröberte Gitter $\Lambda' = \Lambda / 2$ mit Abstand $2a$. Die Block-Spin-Abbildung $\mathcal{R}: \mathcal{A}_\Lambda \to \mathcal{A}_{\Lambda'}$ mittelt Linkvariablen über $2^d$ Blöcke via
\[
U'_{l'} = \mathcal{P}_G\!\left( \frac{1}{|B_{l'}|} \sum_{l \in B_{l'}} U_l \right),
\]
wobei $\mathcal{P}_G$ die Projektion auf die Eichgruppe $G$ bezeichnet und $B_{l'}$ die Menge der feinen Links ist, die den vergröberten Link $l'$ bilden.
\end{definition}

\begin{proposition}[Skalenaufgelöste Freie-Energie-Zerlegung]
\label{prop:scale-resolved}
Die Gitter-Freie-Energie lässt eine teleskopierende Zerlegung zu:
\[
F_\Lambda[\mu] = \sum_{k=0}^{K} \delta F_k[\mu_k \,|\, \mu_{k+1}] + F_{\Lambda_K}[\mu_K],
\]
wobei $\mu_k$ das Maß auf Skala $k$ (Gitterabstand $2^k a$) ist, $\delta F_k = D_{\mathrm{KL}}(\mu_k \| \mu_k^{\mathrm{eq},k+1})$ die relative Entropie der Skala $k$ bedingt auf Skala $k+1$, und $K = \log_2(L/a)$ die Anzahl der RG-Schritte.
\end{proposition}

\begin{proof}[Beweisskizze]
Jede bedingte freie Energie $\delta F_k$ zerlegt sich durch die Kettenregel für die KL-Divergenz:
\[
D_{\mathrm{KL}}(\mu_k \| \nu_k) = D_{\mathrm{KL}}(\mu_{k+1} \| \nu_{k+1}) + \mathbb{E}_{\mu_{k+1}}\!\bigl[D_{\mathrm{KL}}(\mu_k^{(\cdot)} \| \nu_k^{(\cdot)})\bigr],
\]
wobei $\mu_k^{(\cdot)}$ das bedingte Maß auf Skala $k$ gegeben die vergröberte Konfiguration auf Skala $k+1$ bezeichnet. Teleskopisches Summieren über $k = 0, \ldots, K$ liefert die Zerlegung. Jeder Term $\delta F_k$ wird durch die lokale Poincar\'e-Konstante auf Skala $k$ kontrolliert: Die Varianz jeder Observablen auf Skala $k$, bedingt auf Skalen $> k$, ist durch $\kappa_k^{-1}$-mal die bedingte Dirichlet-Form beschränkt. Die Kontinuums-Massenlücke folgt, \emph{sofern} Vermutung~\ref{conj:gap-transport} (Lücken-Transport) gilt, die sicherstellt, dass die Konstanten $\kappa_k$ nicht degenerieren.
\end{proof}

\begin{remark}[Verbindung zu CL1--CL3]
\label{rem:hierarchical-cl}
Die skalenaufgelöste Zerlegung liefert eine konkrete Realisierung der abstrakten Kontinuumslimes-Bedingungen CL1--CL3. Im Einzelnen: CL1 (thermodynamischer Limes) folgt aus der teleskopierenden Schranke; CL2' (gleichmäßige Poincar\'e in physikalischen Einheiten, siehe Bemerkung~\ref{rem:cl2-relaxation}) ist äquivalent zum Lücken-Transport (Vermutung~\ref{conj:gap-transport}); CL3 (Osterwalder--Schrader) erfordert Reflexionspositivität auf jeder RG-Skala. Lemma~\ref{lem:reflection-positivity} liefert den finiten Wilson-/Heat-Kernel-Eingang, aber die Erhaltung durch die gewählte Block-Spin-Abbildung bleibt eine zusätzliche Kompatibilitätsbedingung und folgt nicht aus dem Lemma allein.
\end{remark}

\begin{proposition}[Bedingte polynomielle LSI-Skalierung via projektive Kontraktion]
\label{prop:polynomial-lsi}
Sei $\kappa_k$ die Poincar\'e-Konstante auf RG-Skala $k$. Angenommen sei
zunächst, dass eine nicht-zirkuläre Block-Spin-Konstruktion einen
OS-/RP-kompatiblen Kern, einen Vergleichssatz von Hilbert-projektiver
Kontraktion zu Poincar\'e-Konstanten-Degradation und eine physikalische
Normierung der Skalenkonstanten liefert. Falls zusätzlich der Block-Spin-Kern
$\mathcal{R}_k$ die Birkhoff-Kontraktionsschranke erfüllt
\[
\tau_B(\mathcal{R}_k) \leq 1 - \delta \quad \text{für ein } \delta > 0 \text{ unabhängig von } k,
\]
wobei $\tau_B$ den Birkhoff-Kontraktionskoeffizienten bezeichnet, dann gilt
\[
\kappa_k \geq \kappa_0 \cdot (1 - C/k^2) \quad \text{für alle } k \leq K,
\]
d.\,h., die Poincar\'e-Konstante würde höchstens polynomiell (nicht
exponentiell) unter Vergröberung degradieren. Unter denselben Zusatzannahmen
würde die Kontinuums-Massenlücke
$m_{\mathrm{phys}} \geq m_0 \cdot (1 - O((\log L)^{-2}))$ erfüllen.
\end{proposition}

\begin{proof}[Beweisskizze]
Der Birkhoff-Kontraktionskoeffizient $\tau_B(\mathcal{R}_k)$ kontrolliert die Kontraktion der Hilbert-Projektiv-Metrik zwischen aufeinanderfolgenden Massen: $d_H(\mu_k, \mu_{k+1}) \leq \tau_B \cdot d_H(\mu_{k-1}, \mu_k)$. Da die KL-Divergenz durch das Quadrat der Hilbert-Metrik beschränkt ist (bis auf Konstanten, die vom Zustandsraum abhängen), erhalten wir
\[
D_{\mathrm{KL}}(\mu_k \| \mu_{k+1}) \leq \tau_B^2 \cdot D_{\mathrm{KL}}(\mu_{k-1} \| \mu_k).
\]
Iteration dieser Kontraktion über $K$ Skalen und Verwendung der gleichmäßigen Schranke $\tau_B \leq 1 - \delta$ ergibt für die kumulative Degradation der Poincar\'e-Konstante
\[
\begin{aligned}
\prod_{k=1}^{K} \frac{\kappa_k}{\kappa_{k-1}}
&\geq \prod_{k=1}^{K} (1 - C/k^2)\\
&= \exp\!\left(\sum_{k=1}^{K} \log(1 - C/k^2)\right)\\
&\geq \exp(-C' \cdot \pi^2/6),
\end{aligned}
\]
wobei $\sum_{k=1}^{\infty} k^{-2} = \pi^2/6$ konvergiert. Daher gilt $\kappa_K \geq \kappa_0 \cdot e^{-C'\pi^2/6} > 0$, und die physikalische Massenlücke $m_{\mathrm{phys}} = \kappa_K / a_K$ wäre durch eine positive Konstante nach unten beschränkt, unabhängig von der Anzahl der RG-Schritte, sobald die oben genannten Normierungs- und Vergleichshypothesen verifiziert sind. Im vorliegenden Beitrag sind diese Hypothesen Beweisziele, keine etablierten Yang--Mills-Eingaben.
\end{proof}

\begin{proposition}[Skalenabhängige LSI via summierbaren Birkhoff-Defekt]\label{prop:kingman-lsi}
Seien $\{R_k\}_{k=1}^K$ die Block-Spin-RG-Transformationen aus
Proposition~\ref{prop:scale-resolved}, und sei
$\tau_B(R_k)$ der Birkhoff-Kontraktionskoeffizient auf Skala~$k$.
Die mittlere Kontraktionsbedingung
\begin{equation}\label{eq:mean-contraction-lsi}
  \lambda \;:=\; \limsup_{K \to \infty}\,\frac{1}{K}\sum_{k=1}^{K}\log\tau_B(R_k) \;<\; 0
\end{equation}
ist eine Diagnostik für mittlere Kontraktion. Eine positive limitierende
LSI-/Poincar\'e-Konstante folgt, wenn die stärkere summierbare
Defektbedingung
\[
  \sum_{k=1}^\infty \tau_B(R_k)^2 < \infty
\]
gilt (oder GT2 eine äquivalente summierbare Fehlerfolge liefert). Unter
dieser Zusatzannahme genügt die volumenunabhängige Poincar\'e-Konstante
\[
  \kappa_K \;\geq\; \kappa_0\,\exp\!\Bigl(\sum_{k=1}^{K}\log(1 - C\,\tau_B(R_k)^2)\Bigr)
  \;\xrightarrow{K \to \infty}\; \kappa_\infty \;>\; 0 .
\]
\end{proposition}

\begin{proof}
Nach dem subadditiven Ergodensatz von Kingman~\cite{Kingman1968}
erfüllt die subadditive Folge $S_n = \sum_{k=1}^n \log\tau_B(R_k)$
die Relation $S_n/n \to \lambda$ fast sicher und in $L^1$.
Die Bedingung $\lambda < 0$ stellt mittlere Kontraktion der komponierten
positiven Kerne sicher, impliziert aber nicht, dass
$\prod_k(1-C\tau_B(R_k)^2)$ einen positiven Grenzwert besitzt. Unter der
Summierbarkeitsannahme gilt dagegen für alle hinreichend späten Faktoren
$C\tau_B(R_k)^2\leq 1/2$, und
\[
  \sum_k |\log(1-C\tau_B(R_k)^2)|
  \leq 2C\sum_k \tau_B(R_k)^2 < \infty .
\]
Folglich konvergiert das Produkt gegen eine strikt positive Zahl, also
$\kappa_\infty>0$.
\end{proof}

\begin{remark}[Verbindung zu Bauerschmidt--Bodineau--Dagallier: uniforme LSI via Polchinski-RG]
\label{rem:bbd-polchinski}
Die uniforme Birkhoff-Kontraktionsbedingung in Proposition~\ref{prop:polynomial-lsi} besitzt ein präzises Analogon im Programm von Bauerschmidt, Bodineau und Dagallier~\cite{BBD-Survey2024,BBD-Ising2024,BBD-Phi4-2024}, die uniforme Log-Sobolev-Ungleichungen für $\varphi^4_2$, $\varphi^4_3$ und nahe-kritische Ising-Modelle mittels eines \emph{multiskalaren Bakry--\'Emery-Kriteriums} beweisen, das in Termen der Polchinski-Gleichung (Renormierungsgruppe) formuliert ist.

Ihre zentrale Einsicht ist, dass die statische Bakry--\'Emery-Krümmungsbedingung $\mathrm{Hess}(V) \geq \kappa\,\mathrm{Id}$ (die für nicht-log-konkave Masse scheitert) durch ein \emph{skalenabhängiges} Kriterium ersetzt werden kann: Es genügt, die Hessian des effektiven Potentials \emph{entlang des Polchinski-Flusses} auf jeder RG-Skala zu kontrollieren. Unter der optimalen Annahme beschränkter Suszeptibilität $\chi(\Lambda, \beta) \leq C$ beweisen sie, dass die LSI-Konstante $\rho$ uniform im Gittervolumen $|\Lambda|$ und in der Gitter-Regularisierung ist -- für alle Kopplungskonstanten in endlichem Volumen und uniform im Volumen in der gesamten Hochtemperaturphase~\cite{BBD-Phi4-2024}.

Die strukturelle Parallele zu unserem Framework ist:
\begin{enumerate}
    \item \textbf{Unser Ansatz:} Uniforme Birkhoff-Kontraktion $\tau_B(\mathcal{R}_k) \leq 1 - \delta$ für den Block-Spin-Kern $\mathcal{R}_k$ sichert polynomielle Degradation der Poincar\'e-Konstante unter Vergröberung (Proposition~\ref{prop:polynomial-lsi}).
    \item \textbf{BBD-Ansatz:} Beschränkte Suszeptibilität $\chi \leq C$ entlang des Polchinski-Flusses sichert uniforme LSI, wobei die Suszeptibilität die Rolle einer inversen Spektrallücke spielt.
    \item \textbf{Äquivalenz:} Beide Bedingungen verhindern die Degeneration der funktionalen Ungleichungskonstante unter dem RG-Fluss. Die Birkhoff-Kontraktion kontrolliert Konvergenz in der Hilbert-Projektiv-Metrik; das BBD-Kriterium kontrolliert Konvergenz in Wasserstein-/Transport-Metriken. Für skalare Feldtheorien ist das BBD-Kriterium verifiziert~\cite{BBD-Phi4-2024}.
\end{enumerate}

Für Gitter-Yang--Mills ist die Erweiterung nicht-trivial, da die Polchinski-Gleichung für $\mathrm{SU}(N)$-wertige Link-Variablen nicht in Standardform verfügbar ist. Bestehende rigorose Strong-Coupling-Arbeiten stützen volumenuniforme Poincar\'e- und Log-Sobolev-Ungleichungen in festen Gitter- oder unendlichen Volumenregimen, während der verschwindende Gitterabstand ein separater Schritt bleibt. Jüngere Heat-Kernel-Yang--Mills-Preprints schlagen mögliche volumen- und gitter-uniforme Routen über (A)~eine Bakry--\'Emery-Krümmungsschranke auf Uhlenbeck-Scheiben und (B)~ein Hochtemperatur-Dobrushin-Kriterium mit expliziten Konstanten vor; in der vorliegenden Arbeit werden diese Claims als externe Evidenz und künftige Zielaussagen behandelt, nicht als Eingaben. Falls solche Schranken in einer mit eichkovariantem Blocking kompatiblen Form bewiesen werden, könnte ein Polchinski-artiger Fluss auf $\mathrm{SU}(N)$-Bündeln einen rigorosen Weg zur Lösung der uniformen Kontraktionsbedingung (SP4) für den Kontinuumslimes liefern.

Im BBD-Framework ist die Schlüsselgröße, die die LSI-Konstante kontrolliert, die Suszeptibilität $\chi = \sum_x \langle \sigma_0 \sigma_x \rangle_c$, die in den bewiesenen skalaren Modellen eng mit inverser Korrelationslänge und Spektrallückenkontrolle zusammenhängt. Damit stützt das BBD-Programm die strukturelle Erwartung dieses Papers: \emph{Uniforme funktionale Ungleichungen und Massenlückenkontrolle sind eng gekoppelte Probleme}.
\end{remark}

\begin{corollary}[Mittlere Kontraktion als Diagnostik und summierbarer-Defekt-Kriterium]\label{cor:kingman-gap}
Unter den Voraussetzungen von Proposition~\ref{prop:polynomial-lsi} ist die mittlere Kontraktionsbedingung
\begin{equation}\label{eq:mean-contraction}
  \lambda := \limsup_{K \to \infty} \frac{1}{K} \sum_{k=1}^{K} \log \tau_B(R_k) < 0.
\end{equation}
eine nützliche Diagnostik für die Kontraktion der komponierten RG-Kerne. Sie ist für sich genommen keine untere Schranke für die limitierende Poincar\'e-Konstante. Eine positive Grenzlücke folgt, wenn die stärkere summierbare Defektbedingung
\begin{equation}\label{eq:summable-birkhoff-defect}
  \sum_{k=1}^{\infty} \tau_B(R_k)^2 < \infty
\end{equation}
(oder allgemeiner GT2 mit $\sum_k \varepsilon_k < \infty$) gilt. Dann erfüllen die Poincar\'e-Konstanten
\begin{equation}
  \kappa_K \geq \kappa_0 \cdot \exp\!\left(\sum_{k=1}^{K} \log(1 - C \tau_B(R_k)^2)\right) \xrightarrow{K \to \infty} \kappa_\infty > 0
\end{equation}
nach Verwerfen endlich vieler Anfangsskalen, auf denen $C\tau_B(R_k)^2$ noch nicht klein ist.
\end{corollary}

\begin{proof}[Beweisskizze]
Nach dem subadditiven Ergodensatz von Kingman existiert der Limes
\[
  \lambda = \lim_{K \to \infty} \frac{1}{K} \sum_{k=1}^{K} \log \tau_B(R_k)
\]
fast sicher. Die Subadditivität folgt aus $\log \tau_B(R_n \circ \cdots \circ R_1) \leq \sum_{k=1}^n \log \tau_B(R_k)$. Aus $\lambda < 0$ folgt mittlere exponentielle Kontraktion der Hilbert-Projektiv-Metrik. Da aber $K\lambda' \to -\infty$ für $\lambda'<0$, liefert diese Durchschnittsaussage keine positive untere Schranke für $\kappa_K$. Gilt dagegen \eqref{eq:summable-birkhoff-defect}, dann ist für alle hinreichend großen $k$ die Größe $C\tau_B(R_k)^2 \leq 1/2$ und
\[
  \sum_k |\log(1-C\tau_B(R_k)^2)|
  \leq 2C\sum_k \tau_B(R_k)^2 < \infty .
\]
Folglich konvergiert $\prod_k(1-C\tau_B(R_k)^2)$ gegen eine strikt positive Zahl, was $\kappa_\infty>0$ und damit die bedingte Lückenaussage liefert.
\end{proof}

\begin{remark}[Numerische Diagnostik]\label{rem:kingman-numerical}
Die mittlere Kontraktionsbedingung~\eqref{eq:mean-contraction} wird numerisch auf dem getesteten $\beta$-Gitter beobachtet (Bemerkung~\ref{rem:birkhoff-numerics}): $\lambda \approx -10$ bei $\beta=1$ (starke Kopplung), $\lambda \approx -0{,}18$ bei $\beta=8$ und $\lambda \approx -4 \times 10^{-5}$ bei $\beta=20$. Während die uniforme Bedingung $\tau_B < 1$ für $\beta \geq 8$ scheitert (wo $\tau_B(R_0) \approx 1$ auf der gröbsten Skala), motiviert der beobachtete negative Exponent das summierbare-Defekt-Programm. Er ersetzt die uniforme Birkhoff-Annahme in Proposition~\ref{prop:polynomial-lsi} jedoch nur, wenn \eqref{eq:summable-birkhoff-defect} oder GT2 separat verifiziert wird.
\end{remark}

\subsubsection*{Schwellenwert-Vermutung und diagnostisches Framework}

Das Gap-Transport-Problem lässt sich präzise als \emph{Schwellenwert-Vermutung} formulieren -- die minimale strukturelle Bedingung, die sicherstellt, dass die spektrale Lücke den Kontinuumslimes $\beta \to \infty$ überlebt.

\begin{conjecture}[Gap Transport -- GT-YM]\label{ax:gt-ym}
Sei $\mathcal{R}_k$ die Block-Spin-RG-Transformation (Definition~\ref{def:block-rg}) auf Skala $k$, und sei $\kappa_k$ die Poincar\'e-Konstante der effektiven Theorie $S_k$ auf dem Gitter $\Lambda_k$ mit Gitterabstand $2^k a$. Die \textbf{Gap-Transport-Vermutung} fordert:
\begin{enumerate}
\item[(GT1)] \textbf{Skalenkoerzivität:} $\kappa_k \geq 0$ für alle $k$, wobei $\kappa_k = 0$ nur für Eich-Nullmoden gilt.
\item[(GT2)] \textbf{Kontrollierter Transport:} Es existiert eine Defektfolge
$0 \leq \varepsilon_k \leq 1/2$, sodass
\begin{equation}\label{eq:gap-transport}
  \kappa_{k+1} \geq \kappa_k(1-\varepsilon_k),
  \qquad \text{mit } \sum_{k=0}^{\infty} \varepsilon_k < \infty.
\end{equation}
\item[(GT3)] \textbf{Kontinuumsstabilität:} Das summierbare Defektbudget
$\sum \varepsilon_k$ bleibt für $\beta \to \infty$ uniform.
\end{enumerate}
\end{conjecture}

\begin{remark}[Kerninhalt]\label{rem:gt-core}
GT-YM isoliert die zentrale Schwierigkeit: Die Massenlücke muss über Skalen transportiert werden, nicht nur auf einer initialen Skala gezeigt werden. Die Gap-Kontrollgröße $G(S_k) := \kappa_k$ ist die Poincar\'e-Konstante des bedingten Einzellink-Masses, die die Birkhoff-Kontraktion $\tau_B(R_k) = 1 - \kappa_k + O(\kappa_k^2)$ kontrolliert. Korollar~\ref{cor:kingman-gap} zeigt, dass ein negativer Kingman--Ljapunow-Exponent eine nützliche Diagnostik ist; GT2 oder ein summierbarer Birkhoff-Defekt ist die eigentliche hinreichende Bedingung für die Erhaltung einer positiven unteren Schranke.
\end{remark}

\begin{proposition}[Gap-Transport-Implikationen und Diagnostiken]\label{prop:gt-equivalences}
Die folgenden Aussagen kodieren verwandte Gap-Transport-Mechanismen, sobald
die jeweils genannten Uniformitäts-, Lokalitäts-, OS- und
Summierbarkeitsannahmen vorliegen. Sie werden nicht als unbedingte
Äquivalenzen für das Yang--Mills-Kontinuumsproblem behauptet:
\begin{enumerate}
\item[(i)] \textbf{Skalenstabiles Poincar\'e:} GT-YM gilt mit summierbarem Defektbudget.
\item[(ii)] \textbf{Exponentieller Korrelationszerfall:} Für alle eichinvarianten Observablen $f, g$ mit $\mathrm{dist}(\mathrm{supp}(f), \mathrm{supp}(g)) = r$ gilt
  $|\langle f\, g \rangle - \langle f \rangle \langle g \rangle| \leq C\, e^{-m r}$
  mit $m > 0$ uniform im Gittervolumen und $\beta$.
\item[(iii)] \textbf{Transferlücke:} Die Transfermatrix $T_k$ der Block-Spin-RG auf Skala $k$ hat eine lokale Zylindersektor-Transferlücke $\Delta_k \geq \Delta_0 > 0$ uniform in $k$.
\item[(iv)] \textbf{Summierbarer Birkhoff-Defekt:} $\sum_k \tau_B(R_k)^2 < \infty$, oder äquivalent eine verifizierte GT2-Defektfolge mit $\sum_k\varepsilon_k<\infty$.
\end{enumerate}
Die Implikationen zwischen (i), (ii) und (iii) benötigen die jeweiligen
Dobrushin--Shlosman-Vergleichssätze, Transfermatrix-Abschlüsse und uniformen
Konstanten. Item~(iv) ist ein hinreichendes Gap-Transport-Kriterium; die
schwächere numerische Bedingung $\lambda<0$ allein wird nicht als äquivalent
zur Massenlücke behauptet.
\end{proposition}

\begin{remark}[Der Konstanten-Drift-Scanner]\label{rem:drift-scanner}
Um zu lokalisieren, wo der Gap-Transport bricht, scanne man drei Typen von Budget-Erosion:
\begin{enumerate}
\item \textbf{Positivitätsbudget:} Welche Positivität (Reflexionspositivität, Konvexität der Wirkung) schwächt sich unter Blockierung ab? Im Regime starker Kopplung ($\beta < \beta_0$) gilt die Dobrushin-Bedingung $\eta < 1$ trivialerweise. Mit wachsendem $\beta$ divergiert $\eta \to \infty$ (Bemerkung~\ref{rem:dobrushin-numerics}), und die Positivität muss über die Typical-Set-Verbesserung zurückgewonnen werden.
\item \textbf{Entropiebudget:} Wo wächst die Zahl effektiver Konfigurationen schneller als die Energiebarriere? Der Instanton-Gas-Beitrag $\sim e^{-8\pi^2/g^2}$ ist exponentiell klein, akkumuliert sich aber über RG-Schritte.
\item \textbf{Ward-Budget:} Welche Eichinvarianz-Identität verliert quantitative Stärke? Die Peter--Weyl-Entwicklung kontrolliert dies: Wenn der Koeffizient $a_\pi(\beta)$ der führenden irreduziblen Darstellung driftet, degradiert die Lücke.
\end{enumerate}
\end{remark}

\begin{remark}[Defektiver Transport als Zwischenziel]\label{rem:defective-transport}
Falls der strikte Transport GT2 scheitert, genügt die defektive Version~\eqref{eq:gap-transport} mit summierbarem Fehler $\sum \varepsilon_k < \infty$ für eine Massenlücke: Der akkumulierte Defekt $\prod_k (1-\varepsilon_k)$ konvergiert gegen einen strikt positiven Grenzwert, sobald $\varepsilon_k \leq 1/2$ und $\sum \varepsilon_k < \infty$ gilt, was
\begin{equation}
m_{\mathrm{phys}} \geq m_0 \prod_{k=0}^{\infty}(1-\varepsilon_k)
\geq m_0 \exp\!\left(-2\sum_{k=0}^{\infty}\varepsilon_k\right) > 0
\end{equation}
liefert (unter Verwendung von $\log(1-x) \geq -2x$ für $x \in [0, 1/2]$). Korollar~\ref{cor:kingman-gap} gibt die entsprechende Birkhoff-Form über $\sum_k\tau_B(R_k)^2<\infty$. Numerisch (Bemerkung~\ref{rem:birkhoff-numerics}) können einzelne $\tau_B(R_k)$ auf der gröbsten Skala gegen $1$ gehen, während $\lambda < 0$ auf dem getesteten Gitter bleibt; dies motiviert, beweist aber nicht, die benötigte Summierbarkeit.
\end{remark}

\begin{remark}[No-Go-Mechanismen]\label{rem:gt-nogo}
GT-YM scheitert genau dann, wenn einer der folgenden Obstruktionen \emph{Quasi-Nullmoden} am RG-Fixpunkt erzeugt:
\begin{enumerate}
\item \textbf{Entropie besiegt Energie:} Der RG-Schritt erzeugt so viele effektive Freiheitsgrade, dass lokale Koerzivität global ,,ausgewaschen'' wird -- Lückenkollaps durch kombinatorische Explosion. Dies ist der Übergang von starker zu schwacher Kopplung.
\item \textbf{Topologische Quasi-Nullmoden:} Instantonen, Center-Vortices oder Monopole erzeugen nahezu invariante Richtungen, die im Kontinuumslimes dichter werden -- Lückenerosion als Mechanismus, nicht als Mystik.
\item \textbf{Eichfixierungsartefakte:} Falls der Transport nur nach Eichfixierung funktioniert, könnte die transportierte Größe ein Artefakt sein, nicht die physikalische Lücke. Finite Gitter-Reflexionspositivität (Lemma~\ref{lem:reflection-positivity}) liefert einen eichinvarianten OS-Test, aber die RG-Konstruktion muss weiterhin zeigen, dass diese Positivität den gewählten Transport übersteht.
\end{enumerate}
Für SU(2) im Regime starker Kopplung ist keiner dieser Mechanismen operativ (Theorem~\ref{thm:conditional-gap}). Für $\beta \to \infty$ ist Mechanismus~(1) der Hauptverdächtige: Die Holley--Stroock-Konstante degradiert wie $e^{-O(\beta)}$, und nur die Typical-Set-Verbesserung ($e^{-O(\sqrt{\beta})}$, Bemerkung~\ref{rem:restricted-significance}) oder die hierarchische RG (Proposition~\ref{prop:polynomial-lsi}) kann kompensieren.
\end{remark}

\begin{lemma}[GT-YM impliziert Massenlücke]\label{lem:gt-massgap}
Falls Conjecture~\ref{ax:gt-ym} gilt, erfüllt die physikalische Massenlücke
\begin{equation}
  m_{\mathrm{phys}} \geq m_0 \cdot \exp\!\left(-C' \sum_{k=0}^{K} \varepsilon_k / \kappa_k\right) > 0
\end{equation}
für eine explizite Konstante $m_0 > 0$, die nur von der Poincar\'e-Konstante $\kappa_0$ bei starker Kopplung abhängt (Theorem~\ref{thm:conditional-gap}).
\end{lemma}

% ======================================================================
% Wasserstein--Entropie-RG-Kontraktion (Brücke: Strategie 2 + 3)
% ======================================================================

\begin{proposition}[Wasserstein--Entropie-RG-Kontraktion (bedingt auf quantitative Transportschranke)]\label{prop:wasserstein-rg}
Man nehme an, dass der Block-Spin-Kern $R_k$, eingeschränkt auf $\mathcal{T}_\varepsilon$, die Lipschitz-Transportschranke $\mathrm{Lip}(R_k|_{\mathcal{T}}) \leq 1 - c_1/\sqrt{\beta}$ erfüllt (Conjecture~\ref{conj:tensorisation}).

\medskip
\noindent
Sei $\mu_\Lambda$ das Gitter-Eichmass bei Kopplung $\beta$,
$\mathcal{T}_\varepsilon = \{U : |W(U) - \langle W\rangle| \leq \varepsilon\}$
die typische Menge mit $\varepsilon = c\sqrt{\beta}$
(Proposition~\ref{prop:restricted-poincare}),
und $R_k$ die $k$-te Block-Spin-RG-Transformation.
Man definiere die \emph{effektive Kontraktion} über die relative Entropie auf der typischen Menge:
\begin{equation}\label{eq:wasserstein-contraction}
  \delta_k(\beta) \;:=\; 1 - \sup_{\mu,\nu \in \mathcal{P}(\mathcal{T}_\varepsilon)}
  \frac{D_{\mathrm{KL}}(R_k\mu \,\|\, R_k\nu)}{D_{\mathrm{KL}}(\mu \,\|\, \nu)}.
\end{equation}
Dann gilt:
\begin{enumerate}[(i)]
\item \textbf{Polynomielle Kontraktionsrate:}
  \begin{equation}\label{eq:polynomial-rate}
    \delta_k(\beta) \;\geq\; \frac{c''}{\sqrt{\beta}}
  \end{equation}
  für alle $\beta > 0$ und alle RG-Schritte $k$, bedingt auf die oben genannte Lipschitz-Transportschranke.
  Dies ersetzt die exponentiell degradierende Birkhoff-Schranke
  $\delta_{\mathrm{typ}} \geq c'\,e^{-c''\sqrt{\beta}}$
  aus Lemma~\ref{lem:typical-rg-old} (unten) durch eine \emph{polynomiell}
  degradierende Rate.

  \emph{Mechanismus:} Der Block-Spin-Kern $R_k$ wirkt durch Mittelung über
  einen Block von $L^d$ Link-Variablen. Auf der typischen Menge $\mathcal{T}_\varepsilon$
  ist die Oszillation des bedingten Gewichts beschränkt durch $2\varepsilon = 2c\sqrt{\beta}$,
  nicht durch $2 d_\ell \beta$ (volle Oszillation).
  Die benötigte Lipschitz-Transportkonstante
  $\mathrm{Lip}(R_k|_{\mathcal{T}}) \leq 1 - c_1/\sqrt{\beta}$
  wird hier nicht hergeleitet; sie ist der quantitative Input von
  Conjecture~\ref{conj:tensorisation}. Bestehende Lipschitz-Transportresultate
  wie Shenfelds Satz~\cite{Shenfeld2024} motivieren dieses Ziel, liefern aber
  nicht die benötigte SU(N)-Gittereichfeld-Konstante.
  Die Kontraktion der relativen Entropie erbt diese Rate:
  $D_{\mathrm{KL}}(R_k\mu \| R_k\nu) \leq (1 - c_1/\sqrt{\beta})^2 \,D_{\mathrm{KL}}(\mu\|\nu)$
  durch die verstärkte Datenverarbeitungsungleichung mittels Lipschitz-Transport.

\item \textbf{Leakage-Kontrolle:}
  Der Beitrag der schlechten Menge ist durch eine $\beta$-unabhängige Konstante beschränkt:
  \[
    \mu(\mathcal{T}_\varepsilon^c) \;\leq\; \delta_0 := 2e^{-c^2/C} \;<\; \tfrac{1}{4},
  \]
  und die Worst-Case-Kontraktion auf $\mathcal{T}_\varepsilon^c$ ist trivial ($\delta = 0$).

\item \textbf{Bedingter Kingman--Ljapunow-Exponent:}
  Der effektive Exponent erfüllt
  \begin{equation}\label{eq:lambda-polynomial}
  \begin{aligned}
    \lambda_{\mathrm{eff}}(\beta)
    &\leq -(1 - \delta_0)\,\frac{c''}{\sqrt{\beta}} + O(1/\beta)\\
    &< 0 \quad\text{für alle } \beta > 0,
  \end{aligned}
  \end{equation}
  da $1 - \delta_0 > 3/4 > 0$.
\end{enumerate}
\end{proposition}

\begin{proof}[Beweis der bedingten Implikation]
\textbf{Schritt 1 (Eingeschränkte Oszillation).}
Auf $\mathcal{T}_\varepsilon$ mit $\varepsilon = c\sqrt{\beta}$
erfüllt das bedingte Gewicht $W(U_\ell \,|\, U_{\partial\ell})$
$\mathrm{osc}(W|_{\mathcal{T}}) \leq 2\varepsilon = 2c\sqrt{\beta}$.
Dies ist die zentrale Verbesserung: Die volle Oszillation beträgt $2 d_\ell \beta = 12\beta$
(in $d=4$), aber die Einschränkung auf die typische Menge reduziert sie auf $O(\sqrt{\beta})$.

\medskip
\textbf{Schritt 2 (Lipschitz-Transport auf der typischen Menge).}
Angenommen sei die quantitative Transportschranke aus Conjecture~\ref{conj:tensorisation}
für den Block-Spin-Kern $R_k$, eingeschränkt auf Konfigurationen in
$\mathcal{T}_\varepsilon$. Die zugehörige Lipschitz-Transportkonstante ist
\[
  \mathrm{Lip}(R_k|_{\mathcal{T}})
  \;\leq\; 1 - \frac{c_1}{\sqrt{\beta}}
\]
als Annahme. Der vorherige Versuch, diese Schranke direkt aus Holley--Stroock
und einem allgemeinen Shenfeld-Transporttheorem abzuleiten, liefert für
SU(N)-Gittereichfelder keine hinreichend gute Konstante. Die Proposition
beweist daher nur die Implikation von dieser quantitativen Transportschranke
zum negativen effektiven Exponenten.

\medskip
\textbf{Schritt 3 (Kontraktion der relativen Entropie).}
Die Datenverarbeitungsungleichung liefert $D_{\mathrm{KL}}(R_k\mu \| R_k\nu) \leq D_{\mathrm{KL}}(\mu\|\nu)$
für jeden Markov-Kern. Die verstärkte Version für Lipschitz-Transport
(Shenfeld~\cite{Shenfeld2024}, Korollar~1.3) ergibt
\[
  D_{\mathrm{KL}}(R_k\mu \| R_k\nu)
  \;\leq\; \bigl(\mathrm{Lip}(R_k|_{\mathcal{T}})\bigr)^2 \cdot D_{\mathrm{KL}}(\mu\|\nu)
  \;\leq\; \Bigl(1 - \frac{c_1}{\sqrt{\beta}}\Bigr)^2 \cdot D_{\mathrm{KL}}(\mu\|\nu).
\]
Daher $\delta_k = 1 - (1 - c_1/\sqrt{\beta})^2 \geq 2c_1/\sqrt{\beta} - c_1^2/\beta \geq c''/\sqrt{\beta}$.

\medskip
\textbf{Schritt 4 (Leakage).}
Wir verwenden sub-Gausssche Konzentration auf Produkten kompakter Gruppen
(Gromov--Milman~\cite{GromovMilman1983}, Ledoux~\cite{Ledoux2001}):
Die Wilson-Wirkung hat Lipschitz-Konstante $O(\sqrt{\beta})$ pro Plaquette,
und das Produktmass auf $G^{|\Lambda|}$ erfüllt eine logarithmische
Sobolev-Ungleichung mit von $|\Lambda|$ unabhängiger Konstante, was Gausssche
Konzentration mit Varianzproxy $O(\beta)$ ergibt
und somit $\mu(|W - \langle W\rangle| > t) \leq 2\exp(-t^2/(C\beta))$.
Bei $t = \varepsilon = c\sqrt{\beta}$:
\[
  \mu(\mathcal{T}_\varepsilon^c) \;\leq\; 2e^{-c^2/C},
\]
was eine $\beta$-unabhängige Konstante $< 1/4$ für geeignetes $c$ ist.
Dies genügt für Schritt~5, da eine $\beta$-unabhängige Schranke
$\mu(\mathcal{T}_\varepsilon^c) \leq \delta_0 < 1$ zusammen mit der
polynomiellen Kontraktion $c''/\sqrt{\beta}$ auf der typischen Menge
bereits $\lambda_{\mathrm{eff}} < 0$ für alle $\beta > 0$ liefert.

\begin{remark}[Stärkere Leakage-Schranke via Brascamp--Lieb]
Das Gibbs-Maß bei Kopplung $\beta$ konzentriert sich \emph{stärker}
als das Produktmass. Brascamp--Lieb-Ungleichungen für die
eichfixierte Wirkung legen $\mu(\mathcal{T}_\varepsilon^c) \leq e^{-c'\beta}$
für großes $\beta$ nahe. Dies würde die kombinierte Ljapunow-Abschätzung
verstärken, ist aber für die Schlussfolgerung $\lambda_{\mathrm{eff}} < 0$
nicht erforderlich und bleibt im vorliegenden Setting unbewiesen.
\end{remark}

\medskip
\textbf{Schritt 5 (Kombination).}
Sei $\delta_0 := 2e^{-c^2/C} < 1/4$ die Leakage-Schranke aus Schritt~4.
Man zerlege den Ljapunow-Exponenten:
\begin{align*}
  \lambda_{\mathrm{eff}}
  &= \mu(\mathcal{T}) \cdot \log(1 - \delta_k) + \mu(\mathcal{T}^c) \cdot \log 1 \\
  &\leq (1 - \delta_0) \cdot \log\!\bigl(1 - c''/\sqrt{\beta}\bigr) \\
  &\leq -(1 - \delta_0)\,\frac{c''}{\sqrt{\beta}}
    + O\!\bigl(1/\beta\bigr) \\
  &< 0 \quad\text{für alle } \beta > 0,
\end{align*}
da $1 - \delta_0 > 3/4 > 0$. \qedhere
\end{proof}

\begin{remark}[Auflösung des Typische-Mengen-Stabilitätsvorbehalts]\label{rem:wasserstein-rg-significance}
Proposition~\ref{prop:wasserstein-rg} löst den Vorbehalt des früheren
Lemmas~\ref{lem:typical-rg-old} (Anhang) auf, indem sie die
\emph{punktweise} Typische-Mengen-Stabilitätsanforderung
$R_k(\mathcal{T}_\varepsilon) \subset \mathcal{T}_{\varepsilon'}$
durch ein \emph{maßtheoretisches} Kontraktionskriterium ersetzt.
Die zentrale Neuerung ist zweifach:
\begin{enumerate}[(a)]
\item \textbf{Relative Entropie ersetzt Birkhoff-Metrik.}
  Die projektive Birkhoff-Metrik erfordert Positivität des Kerns auf dem
  \emph{gesamten} Zustandsraum; die Kontraktion der relativen Entropie erfordert
  lediglich Lipschitz-Transport auf der \emph{typischen Menge}, wobei die schlechte Menge
  nur über ihr (exponentiell kleines) Maß beiträgt.
\item \textbf{Polynomielle ersetzt exponentielle Degradation.}
  Die Birkhoff-Kontraktion $\delta_{\mathrm{typ}} \sim e^{-c\sqrt{\beta}}$
  (aus Holley--Stroock auf $\mathcal{T}_\varepsilon$) degradiert exponentiell;
  die Wasserstein-/Entropie-Kontraktion $\delta_k \sim 1/\sqrt{\beta}$
  degradiert nur polynomiell. Dies liegt daran, dass die Lipschitz-Konstante
  des eingeschränkten Kerns durch
  $\mathrm{osc}(W|_{\mathcal{T}}) = O(\sqrt{\beta})$ kontrolliert wird,
  nicht durch $\mathrm{osc}(W) = O(\beta)$.
\end{enumerate}
Unter der bedingten Transportschranke von Proposition~\ref{prop:wasserstein-rg}
wäre die resultierende Abschätzung
$\lambda_{\mathrm{eff}}(\beta) \leq -c''/\sqrt{\beta} + o(1)$
das Yang--Mills-Analogon der RH-Even-Dominance-Lücke
$|\mathrm{cert\_gap}| \sim \sqrt{\lambda}$: Beide sind \emph{polynomiell}
im Kontrollparameter. Für Yang--Mills ist diese polynomielle
Transportabschätzung ein Zieltheorem, gestützt nur durch die getesteten
Birkhoff-/Kingman-Diagnostiken (Bemerkung~\ref{rem:birkhoff-numerics}), kein
bewiesenes analytisches Resultat.

\textbf{Kritischer Vorbehalt und partielle Auflösung via Fathi--Mikulincer--Shenfeld.}

\emph{Existenz von Lipschitz-Transport auf $\mathrm{SU}(N)$} ist nun
durch Fathi, Mikulincer und Shenfeld~\cite{FathiMikulincerShenfeld2024}
etabliert, deren Theorem~3 auf gewichtete Riemannsche Mannigfaltigkeiten $(M,g,\mu)$ mit
$\mu = e^{-V}\,\mathrm{d Vol}$ anwendbar ist, sofern die Bakry--\'Emery-Bedingung
$\mathrm{CD}(\kappa, \infty)$: $\mathrm{Ric} + \mathrm{Hess}(V) \geq \kappa\,g$
mit $\kappa > 0$ erfüllt ist.

\emph{Verifikation für $\mathrm{SU}(N)$:}
Das Haar-Maß auf $\mathrm{SU}(N)$ mit der biinvarianten Killing-Metrik
erfüllt $\mathrm{Ric} = (N/4)\,g$ (mit $V = 0$),
woraus $\mathrm{CD}(N/4, \infty)$ mit $\kappa = N/4 > 0$ folgt.
Die Wilson-Wirkungs-Störung $W = -\beta\sum_p \mathrm{Re}\,\mathrm{tr}(U_p)$
ist $L$-Lipschitz mit $L = O(\beta)$ (jede Plaquette trägt
$|\nabla \mathrm{Re}\,\mathrm{tr}(U_p)| \leq C$ bei, und jeder Link
ist an $2d(d-1) = 12$ Plaquetten in $d=4$ beteiligt).
Daher existiert nach~\cite{FathiMikulincerShenfeld2024} Theorem~3
ein Lipschitz-Transport $T: (\mathrm{SU}(N)^{|\Lambda|}, \text{Haar})
\to (\mathrm{SU}(N)^{|\Lambda|}, \mu_\beta)$ mit
dimensionsfreier Lipschitz-Konstante $\mathrm{Lip}(T) \leq e^{e^{c\beta^2}}$.

\emph{Die quantitative Lücke.}
Proposition~\ref{prop:wasserstein-rg} erfordert eine Kontraktionsrate
$\delta_k \geq c''/\sqrt{\beta}$ auf der typischen Menge.  Die globale
Fathi--Mikulincer--Shenfeld-Schranke $\mathrm{Lip}(T) \leq e^{e^{c\beta^2}}$
ist doppelt-exponentiell in $\beta$ und damit bei weitem zu schwach für die Massenlücke.
Auf der \emph{typischen Menge} $\mathcal{T}_\varepsilon$ mit
$\mathrm{osc}(W|_{\mathcal{T}}) = O(\sqrt{\beta})$ statt $O(\beta)$
reduziert sich die Lipschitz-Störungskonstante auf $L_{\mathrm{typ}} = O(\sqrt{\beta})$,
was $\mathrm{Lip}(T|_{\mathcal{T}}) \leq e^{e^{c\beta}}$ ergibt -- verbessert,
aber immer noch weit von der benötigten Kontraktion $1 - O(1/\sqrt{\beta})$ entfernt.

\emph{Zusammenfassung der verbleibenden Lücke.}
Das verfügbare Resultat verschiebt die Frage von ``existiert Lipschitz-Transport auf
$\mathrm{SU}(N)$?'' (\emph{ja}, nach~\cite{FathiMikulincerShenfeld2024})
zu ``kann die Lipschitz-Konstante von
$e^{e^{c\beta^2}}$ auf $1 - O(1/\sqrt{\beta})$ auf der typischen Menge verbessert werden?''
Dies ist eine \emph{quantitative} Frage über die Schärfe von
Transportabschätzungen auf kompakten Lie-Gruppen, keine Existenzfrage.

Drei strukturelle Hindernisse verbleiben für die quantitative Schranke:
\begin{enumerate}[(i)]
\item \emph{Gribov-Kopien} ($N \geq 2$):
  verhindern \emph{globalen} Transport mit uniformen Konstanten.
  Die Einschränkung auf die typische Menge mildert dies ab
  (der Gribov-Horizont ist eine Nullmenge für
  generisches $\beta$~\cite{Zwanziger1989}).
\item \emph{Phasenübergang} ($N \geq 3$):
  uniforme Abschätzungen in $\beta$ sind über den
  Confinement-Deconfinement-Übergang hinweg unmöglich.
  Ein stückweises Argument (Starkskopplungsregime + Schwachkopplungsregime
  mit expliziter Lücke) ist erforderlich.
\item \emph{Doppelt-exponentiell $\to$ polynomiell:}
  Die Verbesserung von $e^{e^{c L^2}}$ zu $(1 - c/L)$ erfordert
  die Ausnutzung der \emph{Produktstruktur} des Gitters
  (Tensorisierung des Transports, nicht erfasst durch die Einzelplatz-Schranke
  von~\cite{FathiMikulincerShenfeld2024}).
\end{enumerate}

Für $\mathrm{U}(1)$-Gitter-Eichtheorie (keine Gribov-Kopien,
Torus-Konfigurationsraum, kein Phasenübergang für $d \leq 4$)
kann das vollständige Programm -- einschließlich der quantitativen
Verbesserung via Tensorisierung -- einen besser handhabbaren Machbarkeitsnachweis
liefern.
\end{remark}

\subsection{Tensorisierungsprogramm für die quantitative Lücke}\label{sec:tensorisation}

Die Fathi--Mikulincer--Shenfeld-Schranke (Bemerkung~\ref{rem:wasserstein-rg-significance})
etabliert die \emph{Existenz} von Lipschitz-Transport auf $\mathrm{SU}(N)^{|\Lambda|}$,
jedoch mit einer doppelt-exponentiellen Konstante $\mathrm{Lip}(T) \leq e^{e^{c\beta^2}}$.
Proposition~\ref{prop:wasserstein-rg} erfordert eine skalenlokale Kontraktion
der Ordnung $1 - O(1/\sqrt{\beta})$. Das Schließen dieser \emph{quantitativen Lücke}
ist die zentrale verbleibende Herausforderung.

\begin{definition}[Existenz-vs-Raten-Lücke]\label{def:evr-gap}
Ein Problem weist eine \emph{Existenz-vs-Raten-Lücke} vom Typ $(\alpha, \gamma)$
auf, wenn ein Existenztheorem eine Schranke der Ordnung $e^{e^{c \cdot p^\alpha}}$
liefert, während die Anwendung $1 - O(p^{-\gamma})$ erfordert,
wobei $p$ der Kontrollparameter ist.
Die Yang--Mills-Massenlücke hat Typ $(\alpha, \gamma) = (2, 1/2)$
mit $p = \beta$.
\end{definition}

Der Standardmechanismus zur Verbesserung von Einzelplatz-Schranken zu Produkt-Schranken
ist \emph{Tensorisierung}:

\begin{conjecture}[Tensorisierter Transport auf Gitter-Eichkonfigurationen]
\label{conj:tensorisation}
Sei $\mu_\beta$ das Wilson-Gitter-Eichmass auf
$\mathrm{SU}(N)^{|\Lambda|}$ bei Kopplung $\beta$.
Für $\beta < \beta_{\mathrm{dec}}$ (unterhalb des Deconfinement-Übergangs)
erfüllt die Lipschitz-Transportkonstante
\[
  \mathrm{Lip}(T_{\mathrm{tensor}}) \;\leq\;
  \prod_{\ell \in \Lambda} \mathrm{Lip}(T_\ell|_{\mathcal{T}_\varepsilon})
  \;\leq\; \bigl(1 - c/\sqrt{\beta}\bigr)^{|\Lambda|}
\]
wobei $T_\ell$ der auf die typische Menge $\mathcal{T}_\varepsilon$
eingeschränkte Einzellink-bedingte Transport ist.
\end{conjecture}

\begin{remark}[Warum Tensorisierung die Lücke schließen könnte]
\label{rem:tensorisation-mechanism}
Die FMS-Schranke $e^{e^{c\beta^2}}$ behandelt das Gitter als
\emph{einzelne Mannigfaltigkeit} $\mathrm{SU}(N)^{|\Lambda|}$.
Der Tensorisierungsansatz nutzt die \emph{Produktstruktur} aus:
\begin{enumerate}[(i)]
\item Jeder Link $\ell$ trägt einen bedingten Transport $T_\ell$
  mit Lipschitz-Konstante $\mathrm{Lip}(T_\ell) \leq 1 + C/\sqrt{\beta}$
  bei (aus der eingeschränkten Oszillation $\mathrm{osc}(W_\ell|_{\mathcal{T}}) = O(\sqrt{\beta})$).
\item Die Dobrushin-Bedingung $\eta(\beta) < 1$ stellt sicher, dass
  bedingte Transporte \emph{kontraktiv komponieren}:
  Das Produkt $\prod \mathrm{Lip}(T_\ell)$ explodiert nicht.
\item Auf der typischen Menge liefert die effektive Dobrushin-Konstante
  $\tilde{\eta}(\beta) \leq 1 - c'/\sqrt{\beta}$ (defektives Dobrushin,
  Proposition~\ref{prop:defective-dobrushin})
  die gewünschte $1 - O(1/\sqrt{\beta})$-Kontraktion pro RG-Schritt.
\end{enumerate}
Dies ist strukturell identisch zur Tensorisierung von log-Sobolev-Ungleichungen
(Gross~1975, Ledoux~\cite{Ledoux2001}), jedoch angewandt auf
\emph{Transportabbildungen} statt auf funktionale Ungleichungen.
Der zentrale offene Schritt ist die Verifikation, dass die Dobrushin-artige
Komposition die Lipschitz-Struktur unter Eichzwängen erhält.
\end{remark}

\begin{lemma}[Typische-Mengen-RG mit kontrolliertem Leakage --- Originalversion]\label{lem:typical-rg-old}
\emph{(Ersetzt durch Proposition~\ref{prop:wasserstein-rg}.
Beibehalten zum Vergleich.)}
Die Birkhoff-Metrik-Version der Typische-Mengen-RG-Brücke
liefert $\lambda_{\mathrm{eff}} \leq -c'\,e^{-c''\sqrt{\beta}} + O(e^{-c'\beta}) < 0$,
was schwächer als~\eqref{eq:lambda-polynomial} ist, aber bereits für
einen negativen Kingman-Exponenten ausreicht.
Der Vorbehalt (punktweise Typische-Mengen-Stabilität unter RG) wird durch
Proposition~\ref{prop:wasserstein-rg} aufgelöst, die punktweise Stabilität
durch maßtheoretische Entropie-Kontraktion ersetzt.
\end{lemma}

\subsection{\texorpdfstring{$\Gamma$}{Gamma}-Konvergenz des Block-Spin-RG-Flusses}\label{sec:gamma-rg}

Die uniforme Birkhoff-Kontraktion $\tau_B(R_k) \leq 1 - \delta$ für jeden RG-Schritt ist strukturell zu stark: Sie erfordert mikroskopische Kontrolle, während die Massenlücke ein makroskopisches Grenzphänomen ist. Wir formulieren die RG als Trajektorienproblem, bei dem Kontraktion im Limes entsteht.

\subsubsection*{Schritt 1: RG-Trajektorien in der Hilbert-Metrik}

Anstatt die Block-Spin-RG als Folge diskreter Operatoren $R_1, R_2, \ldots, R_k$ zu behandeln, formulieren wir sie als Trajektorie im Raum positiver Transferoperatoren, die auf projektive Klassen von Dichten $[\mu]$ mit der Hilbert-Metrik~$d_H$ wirken:
\[
  [\mu_{k+1}] = R_k[\mu_k].
\]
Die physikalisch relevante Größe ist nicht $\tau_B(R_k)$ auf einer einzelnen Skala, sondern die \emph{asymptotische Kontraktionsrate} entlang der Trajektorie:
\begin{equation}\label{eq:lyapunov-rg}
  \lambda := \limsup_{n \to \infty} \frac{1}{n} \sum_{k=1}^{n} \log \tau_B(R_k).
\end{equation}
Ein negativer Wert von $\lambda$ ist eine Diagnostik für mittlere Kontraktion der komponierten RG-Kerne. Er ist allein nicht äquivalent zu einer Massenlücke. Eine positive Grenzlücke benötigt zusätzlich eine summierbare Defektabschätzung wie Korollar~\ref{cor:kingman-gap} oder die GT2-Bedingung in Vermutung~\ref{ax:gt-ym}.

\subsubsection*{Schritt 2: Subadditive Struktur}

Die Schlüsselbeobachtung ist Subadditivität:
\begin{equation}\label{eq:subadditive}
  \log \tau_B(R_n \circ \cdots \circ R_1) \;\leq\; \sum_{k=1}^{n} \log \tau_B(R_k).
\end{equation}
Dies ist präzise die Situation von Kingmans subadditivem Ergodensatz \cite{Kingman1968}: Bilden die RG-Schritte eine stationär ergodische Folge, liefert eine \emph{negative mittlere Kontraktion} entlang des RG-Flusses ein Drift-Signal. Einzelne RG-Schritte dürfen $\tau_B(R_k)$ nahe~$1$ haben; eine positive Grenzlücke folgt daraus aber erst mit einer separaten summierbaren Defektabschätzung.

\subsubsection*{Schritt 3: $\Gamma$-Konvergenz des effektiven Funktionals}

Betrachte das effektive RG-Freie-Energie-Funktional
\begin{equation}\label{eq:rg-functional}
  G_L[\mu] = \sum_{k=1}^{\log_L(\Lambda_{\mathrm{UV}}/\Lambda_{\mathrm{IR}})} \log \tau_B(R_k[\mu]),
\end{equation}
wobei $L$ die Blockgröße ist. Für $L \to \infty$ gilt:
\[
  G_L \;\xrightarrow{\;\Gamma\;}\; G_\infty,
\]
wobei $G_\infty$ ein \emph{deterministisches} RG-Freie-Energie-Funktional wäre. Erst wenn dieser Grenzwert zusammen mit einem summierbaren Defektbudget und physikalischer Normierung kontrolliert ist, würde daraus eine Massenlücken-Skalierung $m \sim \exp(-G_\infty)$ folgen. Entscheidend ist, dass der $\Gamma$-Limes als Diagnose auch dann sinnvoll sein kann, wenn einzelne $\tau_B(R_k)$ nahe~$1$ sind.

\subsubsection*{Schritt 4: EDP-Struktur -- warum Fehler sich nicht akkumulieren}

Der RG-Fluss besitzt eine natürliche Dissipation:
\[
  D(R_k) = \text{Entropieproduktion} + \text{Eichmittelung}.
\]
Diese Dissipation liefert eine Energie-Dissipations-Ungleichung:
\begin{equation}\label{eq:edp-rg}
  G_{k+1} - G_k \;\leq\; -D(R_k),
\end{equation}
die als Zielabschätzung verhindern müsste, dass sich Kontraktionsfehler kohärent akkumulieren. Sie ersetzt die fehlende uniforme Schranke erst dann, wenn die Dissipation in ein summierbares Defektbudget übersetzt wurde.

\subsubsection*{Schritt 5: Gribov-Kopien als integrierbare Defekte}

Singularitäten (Gribov-Kopien) erscheinen als lokale Defekte im RG-Fluss. Sturm-artige Bakry--\'Emery-Resultate auf singulären Mannigfaltigkeiten \cite{Sturm2025} zeigen, dass lokale Krümmungsdefekte zulässig sein können, wenn ihr \emph{integrierter Effekt} entlang des Flusses kontrolliert ist. Im vorliegenden Rahmen weist dies auf eine Summierbarkeitsbedingung für die induzierten Krümmungs- und Transportdefekte, nicht bloß auf $\lambda < 0$.

\begin{remark}[Kein versteckter Massenlückenbeweis]\label{rem:no-hidden-proof}
Diese Reformulierung beansprucht keine uniforme Kontrolle für alle~$\beta$. Das Strong-Coupling-Regime liefert nur den \emph{Startpunkt} des RG-Flusses; der fehlende Satz ist die Positivität der Grenzkoerzivität unter einem summierbaren Defektbudget. Das dokumentierte No-Go wird daher nicht umgangen, sondern als präzises Defekt-Summierbarkeitsproblem neu formuliert.
\end{remark}

\begin{remark}[Problemübergreifende Struktur]\label{rem:cross-edp-ym}
Dieser $\Gamma$/EDP-Mechanismus ist strukturell identisch zu:
\begin{itemize}
\item BV-Selektion aus integrierter Dissipation (Navier--Stokes, \cite{FST-NS2026}),
\item Chamelaeon-Screening als $\Gamma$-konvergente Phasenseparation (Dunkle Energie, \cite{Geiger2026-DE}),
\item thermodynamische Selektion als Ersatz algebraischer Kontrolle (BSD, \cite{Geiger2026-BSD}).
\end{itemize}
Alle teilen das Prinzip: \emph{Lokale Kontrolle $+$ Dissipation $\Rightarrow$ globale Selektion im Limes}.
\end{remark}

\begin{lemma}[RG-Gronwall: summierbare Degradation erhält die Spektrallücke]\label{lem:rg-gronwall}
Angenommen, die transportierten Poincar\'e-/LSI-Konstanten erfüllen
\begin{equation}\label{eq:rg-gronwall}
  \kappa_{k+1} \;\geq\; \kappa_k(1-\varepsilon_k),
  \qquad 0 \leq \varepsilon_k \leq \frac12,
  \qquad \sum_{k=1}^{\infty}\varepsilon_k < \infty .
\end{equation}
Dann bleibt die Grenzlücke positiv:
\[
  \inf_{N}\kappa_N \;\geq\; \kappa_0
  \prod_{k=1}^{\infty}(1-\varepsilon_k)
  \;\geq\; \kappa_0\exp\!\Bigl(-2\sum_{k=1}^{\infty}\varepsilon_k\Bigr) \;>\; 0.
\]
\end{lemma}

\begin{proof}
Iteration von \eqref{eq:rg-gronwall} ergibt
$\kappa_N \geq \kappa_0\prod_{k<N}(1-\varepsilon_k)$.
Für $0 \leq \varepsilon \leq 1/2$ gilt $\log(1-\varepsilon)\geq -2\varepsilon$.
Die vorausgesetzte Summierbarkeit liefert daher ein strikt positives unendliches Produkt.
Dies ist die deterministische Form der GT2-/summierbarer-Defekt-Bedingung, nicht eine Folge von $\lambda<0$ allein.
\end{proof}

\begin{remark}[Zusammenhang mit subadditivem RG]\label{rem:gronwall-kingman}
Lemma~\ref{lem:rg-gronwall} liefert das deterministische Gegenstück zum Kingman-subadditiven Ansatz (Abschnitt~\ref{sec:gamma-rg}): Ein negativer Ljapunow-Exponent ist das diagnostische Signal, während Summierbarkeit der induzierten Defekte die hinreichende Bedingung ist. Der Kingman-Ansatz ist allgemeiner, weil er mittlere Kontraktion erkennt; die Gronwall-Form hält genau fest, welche Zusatzabschätzung dieses Signal in eine positive Grenzlücke überführt.
\end{remark}

\begin{remark}[Spektrale Robustheit an Gribov-Kopien (Sturm)]
\label{rem:sturm-gribov}
Der Eichfeld-Konfigurationsraum $\mathcal{A}/\mathcal{G}$ besitzt konische Singularitäten an Gribov-Kopien -- Punkten, an denen die Eichbahn den Gribov-Horizont schneidet. Sturm~\cite{Sturm2025} hat gezeigt, dass Bakry--\'Emery-Krümmungsbedingungen und Spektrallücken-Abschätzungen auf Mannigfaltigkeiten mit konischen Singularitäten gültig bleiben, wobei verallgemeinerte Hardy-Ungleichungen die Standard-Krümmungs-Dimensions-Bedingungen nahe der singulären Punkte ersetzen.

Dieses Resultat liefert theoretische Unterstützung für eine mögliche Stabilitätsroute in Gegenwart von Gribov-Kopien: Man darf hoffen zu zeigen, dass die Poincar\'e-Konstante $\kappa$ an konischen Orbitraum-Singularitäten nicht degeneriert, falls die Eichorbit-Geometrie die Voraussetzungen der verfügbaren Singularraumabschätzungen erfüllt. Kombiniert mit dem BBD-Programm (Bemerkung~\ref{rem:bbd-polchinski}) ist dies Evidenz für eine Robustheitsstrategie, aber kein Beweis, dass die Kontinuums-Massenlücke unter allen topologischen Defekten stabil bleibt.
\end{remark}

\begin{center}
\fbox{\parbox{0.9\textwidth}{
\textbf{Statusmatrix.}\\[4pt]
\textbf{Bewiesenes Gitterresultat:}
\begin{itemize}
\item Lokale Transferlücke $\Delta_\Lambda^{\mathrm{loc}} \geq \kappa_* > 0$ für feste eichinvariante lokale Observablen der Wilson-Gittereichtheorie, beliebige kompakte einfache Gruppe $G$, $\beta < \beta_0(d,G)$ (Theorem~\ref{thm:lattice-gap}).
\item Für $G = SU(2)$, $d = 4$: $\beta_0 = 2\,\mathrm{artanh}(1/6) \approx 0.336$, $\kappa_* = (1 - 6\,\tanh(\beta/2)) \cdot e^{-12\beta}$ (Korollar~\ref{cor:su2-explicit}).
\item Diese lokale Zylindersektor-Lücke ist unabhängig vom Gittervolumen $|\Lambda|$.
\end{itemize}
\textbf{Bedingte Rekonstruktion:}
\begin{itemize}
\item Kontinuums-Massenlücke $\Delta \geq \Delta_0 > 0$ (Theorem~\ref{thm:continuum-gap}).
\item Hinweis: CL2 ist im Wesentlichen äquivalent zur Massenlücke; der nichttriviale Gehalt liegt in der Reduktion auf unabhängig prüfbare Bedingungen.
\end{itemize}
\textbf{Konjekturale RG-/Kontinuumsroute:}
\begin{itemize}
\item Wasserstein-Transport, Tensorisierung, good-scale-density und RP-positive Dirichlet-Form-Koerzivität sind vorgeschlagene Mechanismen für Gap-Transport; sie werden hier nicht bewiesen.
\end{itemize}
\textbf{Numerische Evidenz:}
\begin{itemize}
\item Dobrushin-Defekt- und Birkhoff-/Kingman-Rechnungen stützen die RG-Route nur auf getesteten SU(2)-Trunkierungen.
\end{itemize}
}}
\end{center}

\section{Die Massenlücke als Pattern-A-Stabilitätssystem}

Das Gitter-Massenlückenresultat (Theorem~\ref{thm:lattice-gap}) passt natürlich in ein breiteres Strukturmuster: die in \cite{FST-Unified2026} beschriebene \textbf{Pattern-A-Stabilitätslogik}, die eine gemeinsame Drei-Ebenen-Hierarchie in Spektrallückenproblemen identifiziert. Die Analogie ist heuristisch und nicht formal; sie dient nur als systematische Lesart:

\begin{enumerate}[label=\alph*)]
    \item \textbf{Analytische Rangendlichkeit (Null-Tail):} Wie der arithmetische Kern in der Riemannschen Vermutung rangendlich ist \cite{GeigerRH2026c}, ist die hier kontrollierte spektrale Instabilität des Yang--Mills-Transferoperators auf die gauge-variante Komplementklasse begrenzt; der UV-Tail wird nur im durch Poincar\'e--Dobrushin abgedeckten Regime streng kontrolliert.
    \item \textbf{Eichinvarianz als Stabilitätsbedingung:} Die mit nicht eichinvarianten Konfigurationen assoziierten Nullmoden werden durch die Wahl der \emph{physikalischen Testklasse} eichinvarianter Störungen ausgeblendet.
    \item \textbf{Beschränkte Positivität (Massenlücke):} Strikte Konvexität auf der physikalischen Klasse ist ein variationelles Stabilitätssignal; die endliche Spektrallücke $\Delta > 0$ folgt erst nach den Poincar\'e--Dobrushin- und Transferoperator-Abschätzungen.
\end{enumerate}

In dieser Lesart beruht die Massenlücke auf der Trennung gauge-varianter Nullrichtungen von der physikalischen eichinvarianten Testklasse. Sie ist eine Heuristik zum bewiesenen Gitterresultat, kein eigenständiger Kontinuumsbeweis.

\begin{remark}[Breitere Anwendungen der variationellen Methode]
Die hier eingesetzte Poincar\'e--Dobrushin--Transferoperator-Kette ist nicht spezifisch für Eichtheorien. Dieselbe Architektur kann für Gittermodelle mit (i)~kompaktem lokalem Zustandsraum, (ii)~endlicher Reichweite der Wechselwirkung und (iii)~reflexionspositiver Gewichtsfunktion versucht werden, sofern die jeweiligen Dobrushin-, Transfer- und OS-Hypothesen im konkreten Modell verifiziert sind. Vergleichsfälle umfassen $O(N)$-Sigmamodelle auf dem Gitter (wobei asymptotische-Freiheits-Heuristiken~\cite{Polyakov1975} und rigorose Resultate vom genauen Modell und Regime abhängen; für $N = 2$ (XY-Modell) liefert der Berezinskii--Kosterlitz--Thouless-Übergang algebraischen Abfall ohne Massenlücke, und für $N = 1$ (Ising) verschwindet die Massenlücke nur am kritischen Punkt $T = T_c$, mit exponentiellem Clustering für alle $T \neq T_c$), Gitter-Higgs-Modelle und Gitter-QCD mit Fermionen. Das Freie-Energie-Variationsgerüst aus Abschnitt~2 ist damit ein \emph{Kandidaten-Template} für Spektrallückenargumente, kein automatisches Beweisschema.
\end{remark}

\begin{remark}[Strukturelle Parallelen im FST-Programm]\label{rem:fst-parallels}
Die Massenlücke weist eine Drei-Ebenen-Hierarchie auf, die auch in
anderen Problemen des FST-Programms~\cite{FST-Unified2026} beobachtet wird: (i)~die
führende Ordnung (freie Energie $\mathcal{F}$) ist trivial konvex, was
das Vakuum eindeutig macht, aber die Lückengröße unsichtbar lässt; (ii)~die
KL-Divergenz erster Ordnung bestätigt $\mu_\Lambda$ als Minimierer ohne
die Krümmung zu quantifizieren; (iii)~die Massenlücke $\Delta > 0$ tritt erst
in zweiter Ordnung auf, durch die Poincar\'e--Dobrushin-Schranke und ihre
Transferoperator-Übersetzung. Im Kontinuumslimes spielen näherungsweise
topologische Sektoren\footnote{Topologische Sektoren werden durch die zweite Chern-Klasse $c_2(P) \in H^4(M;\mathbb{Z}) \cong \mathbb{Z}$ (Instantonenzahl) für Hauptfaserbündel $P$ über $S^4$ klassifiziert, oder äquivalent durch 't~Hoofts elektrischen und magnetischen Fluss $(e, m) \in Z(G) \times Z(G)$ auf dem Torus.} (definiert über Gradientenfluss~\cite{Luscher2010}) eine
Rolle analog zu den geraden/ungeraden Spektralsektoren bei der
Riemannschen Vermutung~\cite{GeigerRH2026c}: Die Einzel-Link-Poincar\'e-Konstante
hängt nur vom lokalen bedingten Maß ab und ist unempfindlich gegenüber
globaler Topologie, während ein unabhängig verifiziertes zero-free-Zertifikat
Phasenübergänge erster Ordnung entlang der RG-Trajektorie ausschließen würde.
Diese Parallelen sind heuristisch und gehen nicht in die Beweise der
Theoreme~\ref{thm:lattice-gap}--\ref{thm:continuum-gap} ein.
\end{remark}

\begin{remark}[Bedingte Abbildung auf Kirks Kontinuumskonstruktion]\label{rem:kirk-mapping}
\textbf{Vorbehalt:} Kirks Konstruktion~\cite{Kirk2026} liegt im Mai~2026 als Zenodo-Preprint vor, ist aber nicht unabhängig verifiziert und explizit von Weak-Coupling-Hypothesen abhängig. Die folgende Abbildung ist daher selbst bedingt durch eine unabhängige Prüfung der Claims und Hypothesen. Falls sie verifiziert werden, ließe sich die Konstruktion auf das Gerüst von Theorem~\ref{thm:lattice-gap} und das bedingte Theorem~\ref{thm:continuum-gap} abbilden. Die Entsprechung wäre:

\begin{table}[htbp]
\centering
\small
\caption{Strukturelle Abbildung zwischen den FST-Beweiskomponenten und Kirks Pro-Polymer-Clusterentwicklungs-Framework (Zenodo:19614728).}
\label{tab:fst-kirk-mapping}
\begin{tabularx}{\textwidth}{@{}Y Y@{}}
\toprule
\textbf{FST-Gerüst} & \textbf{Kirk-Konstruktion} \\
\midrule
Holley--Stroock-Schranke $\kappa_{\mathrm{link}} \geq \kappa_{\mathrm{Haar}}\,
e^{-2 d_\ell \beta}$ (Schritt~1)
& Einzelplaquetten-Poincar\'e via beschränkter Störung
  (derselbe Mechanismus) \\[4pt]
Dobrushin--Zegarlinski: $\eta(\beta) < 1$ für $\beta < \beta_0$ (Schritt~2)
& Dobrushin--Shlosman-vollständige Analytizität:
  Erweiterung auf \emph{alle} $\beta$ via Pro-Polymer-
  Clusterentwicklung \\[4pt]
CL1 (OS-positiver Kontinuumslimes)
& Bedingter Preprint-Claim: Anisotropie $O(a^2)$, volle $O(4)$-
  Wiederherstellung im Limes $a \to 0$ \\[4pt]
CL2 (uniformes exponentielles Clustering)
& Benötigt unabhängig verifizierte vollständige Analytizität oder eine
  explizite uniforme Clustering-Abschätzung; ein nullstellenfreier Streifen
  allein genügt nicht \\[4pt]
CL3 (uniforme Normierung)
& Würde aus verifizierten Pro-Polymer-Summierbarkeitsabschätzungen folgen \\
\bottomrule
\end{tabularx}
\end{table}

\noindent
Kirks zentrale Innovation wäre die Erweiterung von Schritt~2 über die
Strong-Coupling-Schwelle $\beta_0$ hinaus: Statt global $\eta(\beta)<1$
zu verlangen, sollen die Pro-Polymer-Abschätzungen die Akkumulation von
Abhängigkeiten zwischen entfernten Links \emph{auf jeder Skala der
Clusterentwicklung separat} kontrollieren. Falls die Hypothesen und
Abschätzungen unabhängig bestätigt werden, würde dies die exponentielle
Divergenz der Holley--Stroock-Konstanten (die in unserem Ansatz für
$\beta>\beta_0$ das Hindernis ist) in eine konvergente, polynomiell
beschränkte Resummation überführen. In diesem bedingten Szenario ergäbe
sich eine volumenunabhängige Spektrallücke
$\Delta \geq \kappa_{\mathrm{phys}} > 0$, die entlang des
Renormierungsgruppenflusses $\beta(a) \to \infty$ bestehen bleibt.

Wenn sich Kirks Konstruktion von SU(2) auf allgemeine kompakte einfache
Lie-Gruppen $G$ erweitern lässt (was Kontrolle der Peter--Weyl-Koeffizienten
$a_\pi(\beta)$ aus Lemma~\ref{lem:reflection-positivity} für alle
Darstellungen von~$G$ erfordert), wären die Annahmen~(CL1)--(CL3) von
Theorem~\ref{thm:continuum-gap} substantiiert; die Kontinuums-Massenlücke
würde dann aus unserem variationellen Gerüst zusammen mit den konstruktiven
Abschätzungen folgen.

\textbf{Anmerkung:} Kirks Preprint~\cite{Kirk2026} betrifft das Gitter-SU(2)-Yang--Mills-System mit einem Pro-Polymer-Cluster-Expansions-Framework und bleibt bis zur unabhängigen Prüfung als konditional zu behandeln. Die Skalierungslimit-Konstruktion beinhaltet auxiliare analytische Strukturen, die spezifisch für die Gitterregularisierung sind. Die Erweiterung auf die reine Kontinuums-Yang--Mills-Theorie (ohne Gitterregularisierung) und auf Eichgruppen jenseits SU(2) bleibt offen.
\end{remark}

\begin{proposition}[Kirk-Abbildung auf die Kontinuums-Massenlücke]
\label{prop:kirk-mapping}
Sei $G$ eine kompakte einfache Lie-Gruppe. Angenommen, eine unabhängig
verifizierte Konstruktion (zum Beispiel eine bestätigte Version von
Kirk~\cite{Kirk2026}) lässt sich von $SU(2)$ auf~$G$ erweitern,
d.\,h.\ die folgenden drei Bedingungen gelten:
\begin{enumerate}[label=\textup{(K\arabic*)}]
    \item \textbf{Dobrushin--Shlosman vollständige Analytizität} für die
        Wilson-Gittereichtheorie mit Gruppe~$G$ bei allen Kopplungen $\beta > 0$:
        Das Gibbs-Maß erfüllt gleichmäßigen exponentiellen Zerfall der
        Korrelationen unter allen Randbedingungen, mit einer im Volumen $|\Lambda|$
        gleichmäßigen Rate. (Für die Wilson-Wirkung werden die Peter--Weyl-Koeffizienten
        $a_\pi(\beta) \geq 0$ durch den Osterwalder--Seiler-Eingang zu Lemma~\ref{lem:reflection-positivity} bereitgestellt;
        der Inhalt von~\textup{(K1)} ist die Dobrushin--Shlosman-Bedingung am
        \emph{Maß}, nicht am Plaketten-Gewicht.)
    \item \textbf{$O(a^2)$-Anisotropie-Schranke}: Die Gitter-Schwingerfunktionen
        erfüllen $|S^{(k)}_a - S^{(k)}_{\mathrm{cont}}| \leq C_k\,a^2$ für
        eine dichte Klasse eichinvarianter lokaler Observablen, sodass
        der Kontinuumslimes die volle euklidische $O(4)$-Symmetrie wiederherstellt.
    \item \textbf{Gleichmäßiger nullstellenfreier Streifen}: Es existiert $\varepsilon > 0$
        mit $Z_\Lambda(\beta + i\tau) \neq 0$ für $|\tau| < \varepsilon$,
        gleichmäßig in $|\Lambda|$ und entlang der Trajektorie asymptotischer
        Freiheit $\beta(a)$, was Phasenübergänge erster Ordnung entlang des
        Renormierungsgruppenflusses ausschließt. Diese Bedingung ist ein
        No-first-order-transition-Guardrail und für sich genommen keine
        uniforme Clustering-Abschätzung.
\end{enumerate}
Dann gilt:
\begin{itemize}
    \item[\textup{(a)}] Die Annahmen \textup{(CL1)--(CL3)} von
        Theorem~\ref{thm:continuum-gap} sind erfüllt.
    \item[\textup{(b)}] Die Quanten-Yang--Mills-Theorie mit Eichgruppe~$G$
        auf $\mathbb{R}^4$ besitzt eine Massenlücke $\Delta > 0$.
\end{itemize}
\end{proposition}

\begin{proof}[Beweisskizze]
(K1) $\Rightarrow$ gleichmäßiges exponentielles Clustering (CL2): Vollständige Analytizität
impliziert exponentiellen Abfall trunkierter Korrelationen mit einer im Volumen gleichmäßigen
Rate~\cite{DobrushinShlosman1985}. (K2) $\Rightarrow$ OS-positiver Kontinuumslimes (CL1):
Die $O(a^2)$-Anisotropie-Schranke stellt die Konvergenz der Gitter-Schwingerfunktionen
sicher, und die Grenzfunktionen erben die OS-Positivität vom Gitter
(Lemma~\ref{lem:reflection-positivity}). (K3) schließt einen
First-order-Übergangsmechanismus entlang des RG-Flusses aus, liefert aber nicht
selbständig CL2. Die Nicht-Kollaps-Aussage für die Clustering-Rate muss aus der
verifizierten Complete-Analyticity-/uniformen Clustering-Abschätzung in~(K1)
kommen, zusammen mit den Summierbarkeitschranken, die auch CL3 liefern. Teil~(a)
ist daher bedingt durch diese verifizierten Eingaben. Teil~(b) folgt dann aus
Theorem~\ref{thm:continuum-gap}.
\end{proof}

\begin{corollary}[Hindernisse für allgemeines $G$]
Damit Kirks Konstruktion von $SU(2)$ auf eine allgemeine kompakte einfache
Gruppe~$G$ erweitert werden kann, ist das Haupthindernis~\textup{(K1)}: Die
Dobrushin--Shlosman vollständige Analytizität erfordert die Kontrolle der
Peter--Weyl-Koeffizienten $a_\pi(\beta)$ für \emph{alle} irreduziblen
Darstellungen~$\pi$ von~$G$, nicht nur für die Fundamentaldarstellung. Für
$G = SU(N)$ mit $N \geq 3$ wächst die Anzahl der Darstellungen auf einem gegebenen
Niveau polynomiell in~$N$, und die Pro-Polymer-Abschätzungen müssen gleichmäßig
in diesem Wachstum gelten.
\end{corollary}

\begin{remark}[Erweiterung auf SU(3) via Charakter-Verhältnis-Schranken]
\label{rem:su3-extension}
Kirks Pro-Polymer-Cluster-Expansion~\cite{Kirk2026} beansprucht eine conditional complete-analyticity-Route für die SU(2)-Gittereichtheorie. Falls diese Route unabhängig bestätigt wird, erfordert die Erweiterung auf SU(3) mindestens zwei Zutaten:
\begin{enumerate}
\item \emph{Charakter-Verhältnis-Schranken:} Für die irreduziblen Darstellungen $\rho_\lambda$ von SU(3) mit höchstem Gewicht $\lambda$ müssen die Charakter-Verhältnisse $\chi_\lambda(U)/\dim(\rho_\lambda)$ gleichmäßige Abfallschranken für $|\lambda| \to \infty$ erfüllen. Für SU(2) folgt dies aus der expliziten Formel $\chi_j(\theta) = \sin((2j+1)\theta)/\sin(\theta)$. Für SU(3) liefert die Weyl-Charakterformel $\chi_{(p,q)}$ in Termen von Schur-Polynomen, und die erforderlichen Schranken folgen aus Wärmekern-Abschätzungen auf der Gruppenmannigfaltigkeit:
\[
\frac{|\chi_{(p,q)}(U)|}{\dim(\rho_{(p,q)})} \leq e^{-c(p^2 + pq + q^2) d(U, \mathbf{1})^2}
\]
für eine universelle Konstante $c > 0$, die von der Normierung der Killing-Form abhängt.
\item \emph{Wärmekern-Domination:} Die Wilson-Wirkung erfüllt $e^{\beta \mathrm{Re}\,\mathrm{tr}(U)} \leq K_t(U, \mathbf{1})$ für $t = 1/(2\beta)$, wobei $K_t$ der Wärmekern auf SU(3) ist. Diese Schranke kontrolliert, kombiniert mit der Peter--Weyl-Entwicklung, die Cluster-Expansions-Koeffizienten.
\end{enumerate}
Die vollständige Analytizität für SU(3) würde dann aus dem Dobrushin--Shlosman-Kriterium mit den modifizierten Einzel-Link-Schranken folgen. Dieses Programm ist technisch anspruchsvoll und setzt voraus, dass der SU(2)-Preprint-Fall korrekt und quantitativ ausreichend stark ist.
\end{remark}

% ======================================================================
% ===================================================================
% AI Disclosure
% ===================================================================
\section*{KI-Offenlegung}
Bei der Erstellung dieses Manuskripts wurden KI-gestützte Werkzeuge
(einschließlich Claude von Anthropic und OpenAI Codex/GPT-Werkzeugen) in
unterstützender Funktion eingesetzt, insbesondere für Literaturorganisation,
Konsistenzprüfungen, Übersetzungsunterstützung, Textstrukturierung und
Formulierungshilfe. Der konzeptionelle Entwurf, die wissenschaftliche
Argumentation und die Verantwortung für den gesamten Inhalt liegen
ausschließlich beim Autor.

\section*{Anhang: Konstanten und Normen}
\label{app:constants}

Zur Referenz sammeln wir die wichtigsten Konstanten und Normierungen, die im gesamten Text verwendet werden:

\begin{itemize}
\item \textbf{Eichgruppe:} $G = \mathrm{SU}(N)$, $\dim(G) = N^2 - 1$, $\mathrm{rank}(G) = N-1$.
\item \textbf{Haar-Poincar\'e-Konstanten.}
Gradientenform:
$\kappa_{\mathrm{Haar}}(\mathrm{SU}(2)) = 3$
(Spektrum von $\Delta_{\mathrm{LB}}$ auf $S^3$:
$\lambda_l = l(l+2)$, $l \geq 1$), und
$\kappa_{\mathrm{Haar}}(\mathrm{SU}(3)) = 4/3$
(erster Casimir-Eigenwert $C_2(\mathbf{3}) = 4/3$).
\item \textbf{Links pro Plakette:} $d_l = 2(d-1) = 6$ in $d = 4$ Dimensionen.
\item \textbf{Retraktion:} $\mathrm{Retr}: \mathfrak{g} \to G$ via $X \mapsto e^X$ (Exponentialabbildung) oder Cayley-Abbildung $X \mapsto (1 + X/2)(1 - X/2)^{-1}$.
\item \textbf{Metrik auf $G$:} $d(U, V) = \|U - V\|_F / \sqrt{2N}$ (normierte Frobenius-Norm) oder $d(U, V) = \|\log(U V^\dagger)\|$ (geodätisch).
\item \textbf{Dirichlet-Form:} $\mathcal{E}(f, f) = \sum_l \int |\nabla_l f|^2 \, d\mu$, wobei $\nabla_l f(U) = \frac{d}{dt}\big|_{t=0} f(e^{tX_a} U_l)$ für eine Orthonormalbasis $\{X_a\}$ von $\mathfrak{g}$.
\item \textbf{Nachbarn:} $N_{\mathrm{nbr}} = 2d(2d - 2) = 48$ Plaketten, die sich an einem Gitterpunkt in 4D treffen (jeder der $2d = 8$ Links durch einen Punkt gehört zu $2(d-1) = 6$ Plaketten, abzüglich Doppelzählungen).
\end{itemize}

\begin{thebibliography}{99}
\small
\setlength{\itemsep}{0pt}
\setlength{\parskip}{0pt}
\setlength{\parsep}{0pt}

\bibitem{JaffeWitten2000}
A.~Jaffe and E.~Witten, \emph{Quantum Yang--Mills Theory}, in: \emph{The Millennium Prize Problems}, Clay Mathematics Institute, 2000, pp.~129--152.

\bibitem{Wilson1974}
K.~G.~Wilson, \emph{Confinement of quarks}, Phys.\ Rev.\ D \textbf{10} (1974), 2445--2459.

\bibitem{OsterwalderSchrader1973}
K.~Osterwalder and R.~Schrader, \emph{Axioms for Euclidean Green's functions}, Comm.\ Math.\ Phys.\ \textbf{31} (1973), 83--112.

\bibitem{OsterwalderSchrader1975}
K.~Osterwalder and R.~Schrader, \emph{Axioms for Euclidean Green's functions II}, Comm.\ Math.\ Phys.\ \textbf{42} (1975), 281--305.

\bibitem{OsterwalderSeiler1978}
K.~Osterwalder and E.~Seiler, \emph{Gauge field theories on a lattice}, Ann.\ Phys.\ \textbf{110} (1978), 440--471.

\bibitem{GlimmJaffe1987}
J.~Glimm and A.~Jaffe, \emph{Quantum Physics: A Functional Integral Point of View}, 2nd ed., Springer-Verlag, 1987.

\bibitem{Balaban1985a}
T.~Balaban, \emph{Ultraviolet stability of three-dimensional lattice pure gauge field theories}, Comm.\ Math.\ Phys.\ \textbf{102} (1985), 255--275.

\bibitem{Balaban1985b}
T.~Balaban, \emph{Propagators and renormalization transformations for lattice gauge theories~I}, Comm.\ Math.\ Phys.\ \textbf{95} (1984), 17--40.

\bibitem{Balaban1984c}
T.~Balaban, \emph{Propagators and renormalization transformations for lattice gauge theories~II}, Comm.\ Math.\ Phys.\ \textbf{96} (1984), 223--250.

\bibitem{CreutzJacobsRebbi1983}
M.~Creutz, L.~Jacobs, and C.~Rebbi, \emph{Monte Carlo computations in lattice gauge theories}, Phys.\ Rep.\ \textbf{95} (1983), 201--282.

\bibitem{MorningstarPeardon1999}
C.~J.~Morningstar and M.~J.~Peardon, \emph{The glueball spectrum from an anisotropic lattice study}, Phys.\ Rev.\ D \textbf{60} (1999), 034509.

\bibitem{ChenXu2006}
Y.~Chen et al., \emph{Glueball spectrum and matrix elements on anisotropic lattices}, Phys.\ Rev.\ D \textbf{73} (2006), 014516.

\bibitem{CoverThomas2006}
T.~M.~Cover and J.~A.~Thomas, \emph{Elements of Information Theory}, 2nd ed., Wiley, 2006.

\bibitem{HolleyStroock1987}
R.~Holley and D.~Stroock, \emph{Logarithmic Sobolev inequalities and stochastic Ising models}, J.~Stat.\ Phys.\ \textbf{46} (1987), 1159--1194.

\bibitem{Zegarlinski1990}
B.~Zegar\-li\'nski, \emph{Log-Sobolev inequalities for infinite one-dimensional lattice systems}, Comm.\ Math.\ Phys.\ \textbf{133} (1990), 147--162.

\bibitem{Martinelli1999}
F.~Martinelli, \emph{Lectures on Glauber dynamics for discrete spin models}, in: \emph{Lectures on Probability Theory and Statistics}, Lecture Notes in Mathematics 1717, Springer, 1999, pp.~93--191.

\bibitem{GuionnetZegarlinski2003}
A.~Guionnet and B.~Zegar\-li\'nski, \emph{Lectures on logarithmic Sobolev inequalities}, S\'eminaire de Probabilit\'es XXXVI, Lecture Notes in Mathematics 1801, Springer, 2003, pp.~1--134.

\bibitem{Kirk2026}
H.~Kirk, \emph{From Lattice Mass Gap to Continuum SU(2) Yang--Mills Theory: O(4) Restoration and Osterwalder--Schrader Reconstruction}, Zenodo preprint, 2026.
\href{https://doi.org/10.5281/zenodo.18824738}{doi:10.5281/zenodo.18824738}.

\bibitem{GeigerRH2026c}
L.~Geiger, \emph{From Landscape to Proof: The Even-Dominance Route to the Riemann Hypothesis}, Zenodo preprint, 2026.
\href{https://doi.org/10.5281/zenodo.19035640}{doi:10.5281/zenodo.19035640}.

\bibitem{FST-Unified2026}
L.~Geiger, \emph{FST Spectrum Duality: The Renormalized Free Energy Principle as Physical Instantiation of Functional Stability Theory}, Zenodo preprint, 2026.
\href{https://doi.org/10.5281/zenodo.19036190}{doi:10.5281/zenodo.19036190}.

\bibitem{DobrushinShlosman1985}
R.~L.~Dobrushin and S.~B.~Shlosman, \emph{Constructive criterion for the
uniqueness of Gibbs field}, in: \emph{Statistical Physics and Dynamical
Systems}, Birkhäuser, 1985, pp.~347--370.

\bibitem{Chatterjee2016}
S.~Chatterjee, \emph{The leading term of the Yang--Mills free energy},
J.~Funct.\ Anal.\ \textbf{271} (2016), 2944--3005.

\bibitem{Chatterjee2018}
S.~Chatterjee, \emph{Yang--Mills for probabilists}, in: \emph{Probability and Analysis in Interacting Physical Systems}, Springer Proc.\ Math.\ Stat.\ \textbf{283}, Springer, 2019, pp.~1--16. \texttt{arXiv:1803.01950}.

\bibitem{Milman2009}
E.~Milman, \emph{On the role of convexity in isoperimetry, spectral gap
and concentration}, Invent.\ Math.\ \textbf{177} (2009), 1--43.

\bibitem{Polyakov1975}
A.~M.~Polyakov, \emph{Interaction of Goldstone particles in two dimensions.
Applications to ferromagnets and massive Yang--Mills fields},
Phys.\ Lett.\ B \textbf{59} (1975), 79--81.
\texttt{doi:10.1016/0370-2693(75)90161-6}.

\bibitem{GrossWilczek1973}
D.~J.~Gross and F.~Wilczek, \emph{Ultraviolet behavior of non-abelian gauge theories},
Phys.\ Rev.\ Lett.\ \textbf{30} (1973), 1343--1346.

\bibitem{Politzer1973}
H.~D.~Politzer, \emph{Reliable perturbative results for strong interactions?},
Phys.\ Rev.\ Lett.\ \textbf{30} (1973), 1346--1349.

\bibitem{Luscher2010}
M.~Lüscher, \emph{Properties and uses of the Wilson flow in lattice QCD},
J.\ High Energy Phys.\ \textbf{2010}(8), 071.

\bibitem{BBD-Survey2024}
R.~Bauerschmidt, T.~Bodineau, and B.~Dagallier, \emph{Stochastic dynamics and the Polchinski equation: An introduction}, Probab.\ Surveys \textbf{21} (2024), 200--290. \texttt{arXiv:2307.07619}.

\bibitem{BBD-Ising2024}
R.~Bauerschmidt and B.~Dagallier, \emph{Log-Sobolev inequality for near critical Ising models}, Comm.\ Pure Appl.\ Math.\ \textbf{77} (2024), 2568--2576. \texttt{arXiv:2202.02301}.

\bibitem{BBD-Phi4-2024}
R.~Bauerschmidt and B.~Dagallier, \emph{Log-Sobolev inequality for the $\varphi^4_2$ and $\varphi^4_3$ measures}, Comm.\ Pure Appl.\ Math.\ \textbf{77}(4) (2024), 2579--2612. \texttt{arXiv:2202.02295}.

\bibitem{Sturm2025}
K.-T.~Sturm, \emph{Bakry--\'Emery, Hardy, and spectral gap estimates on manifolds with conical singularities}, Calc.\ Var.\ Partial Differential Equations (2025). \texttt{arXiv:2405.10734}.

\bibitem{Kingman1968}
J.~F.~C.~Kingman, \textit{The ergodic theory of subadditive stochastic processes}, J.~Roy.\ Statist.\ Soc.\ Ser.~B, \textbf{30} (1968), 499--510.

\bibitem{FST-NS2026}
L.~Geiger, \emph{A Thermodynamic Obstruction to Finite-Time Blow-Up in the 3D Navier--Stokes Equations}, Zenodo preprint, 2026.
\href{https://doi.org/10.5281/zenodo.19087449}{doi:10.5281/zenodo.19087449}.

\bibitem{Geiger2026-DE}
L.~Geiger, \emph{Dark Energy as Residual Vacuum Free Energy: A Thermodynamic Bound on the Cosmological Constant}, Zenodo preprint, 2026.
\href{https://doi.org/10.5281/zenodo.19036235}{doi:10.5281/zenodo.19036235}.

\bibitem{Geiger2026-BSD}
L.~Geiger, \emph{The BSD Conjecture as a Positivity Normal-Form Theorem}, Zenodo preprint, 2026.
\href{https://doi.org/10.5281/zenodo.19087443}{doi:10.5281/zenodo.19087443}.

\bibitem{BougerolLacroix1985}
P.~Bougerol \& J.~Lacroix, \emph{Products of Random Matrices with Applications to Schrödinger Operators}, Birkhäuser, 1985.

\bibitem{Geiger2026-TU}
L.~Geiger, \emph{Anomalous Dissipation and the Turbulent Cascade: A Variational Free-Energy Characterisation of Kolmogorov Scaling}, Zenodo preprint, 2026.
\href{https://doi.org/10.5281/zenodo.19056813}{doi:10.5281/zenodo.19056813}.

\bibitem{Geiger2026-PvsNP}
L.~Geiger, \emph{An Entropic Perspective on P vs NP: Reformulation via Kolmogorov Complexity}, Zenodo preprint, 2026.
\href{https://doi.org/10.5281/zenodo.19056809}{doi:10.5281/zenodo.19056809}.

\bibitem{Shenfeld2024}
S.~Shenfeld, \emph{Exact renormalization group and functional inequalities via Lipschitz transport}, arXiv:2205.01642v3, 2024.

\bibitem{FathiMikulincerShenfeld2024}
M.~Fathi, D.~Mikulincer and Y.~Shenfeld, \emph{Transportation onto log-Lipschitz perturbations}, Calc.\ Var.\ Partial Differential Equations \textbf{63} (2024), Article~61. arXiv:2305.03786.

\bibitem{Zwanziger1989}
D.~Zwanziger, \emph{Local and renormalizable action from the Gribov horizon}, Nuclear Phys.\ B \textbf{323} (1989), 513--544.

\bibitem{Ledoux2001}
M.~Ledoux, \emph{The Concentration of Measure Phenomenon}, Mathematical Surveys and Monographs \textbf{89}, American Mathematical Society, 2001.

\bibitem{GromovMilman1983}
M.~Gromov and V.\,D.~Milman, \emph{A topological application of the isoperimetric inequality}, Amer.\ J.\ Math.\ \textbf{105} (1983), 843--854.

\bibitem{Rothaus1985}
O.\,S.~Rothaus, \emph{Diffusion on compact Riemannian manifolds and logarithmic Sobolev inequalities}, J.~Funct.\ Anal.\ \textbf{62} (1985), 358--367.

\end{thebibliography}

\end{document}
